\documentclass[11pt,letterpaper]{article}
\usepackage[T1]{fontenc}
\usepackage[utf8]{inputenc}
\usepackage{lmodern}
\usepackage[margin=1in,headheight=15pt]{geometry}
\usepackage{amsmath,amssymb,amsthm,mathtools,mathrsfs}
\usepackage{microtype,booktabs,longtable,array,enumitem}
\usepackage[dvipsnames]{xcolor}
\usepackage{fancyhdr,aliascnt}
\usepackage[colorlinks=true,linkcolor=MidnightBlue,citecolor=MidnightBlue,urlcolor=MidnightBlue]{hyperref}
\usepackage[nameinlink,noabbrev]{cleveref}
\hypersetup{%
  pdftitle={Jordan Separation, Cauchy Theory, and Stokes Calculus on the Surcomplex Plane},
  pdfauthor={Research synthesis},
  pdfsubject={Surcomplex topology, Hahn differential forms, semialgebraic winding, microscopic residues}}
\pagestyle{fancy}\fancyhf{}
\fancyhead[L]{\small Surcomplex contours and Stokes calculus}
\fancyhead[R]{\small Topologies, chains, and residues}
\fancyfoot[C]{\thepage}
\setlength{\parskip}{0.3em}
\setlength{\parindent}{1.3em}
\setlist{nosep,leftmargin=2em}
\allowdisplaybreaks[2]

\newtheorem{theorem}{Theorem}[section]
\newaliascnt{lemma}{theorem}\newtheorem{lemma}[lemma]{Lemma}\aliascntresetthe{lemma}
\newaliascnt{proposition}{theorem}\newtheorem{proposition}[proposition]{Proposition}\aliascntresetthe{proposition}
\newaliascnt{corollary}{theorem}\newtheorem{corollary}[corollary]{Corollary}\aliascntresetthe{corollary}
\theoremstyle{definition}
\newaliascnt{definition}{theorem}\newtheorem{definition}[definition]{Definition}\aliascntresetthe{definition}
\newaliascnt{example}{theorem}\newtheorem{example}[example]{Example}\aliascntresetthe{example}
\theoremstyle{remark}
\newaliascnt{remark}{theorem}\newtheorem{remark}[remark]{Remark}\aliascntresetthe{remark}
\crefname{theorem}{theorem}{theorems}\crefname{lemma}{lemma}{lemmas}
\crefname{proposition}{proposition}{propositions}\crefname{corollary}{corollary}{corollaries}
\crefname{definition}{definition}{definitions}\crefname{example}{example}{examples}
\crefname{remark}{remark}{remarks}

\newcommand{\No}{\mathbf{No}}\newcommand{\SC}{\mathbf{SC}}
\newcommand{\R}{\mathbb R}\newcommand{\C}{\mathbb C}\newcommand{\Q}{\mathbb Q}
\newcommand{\N}{\mathbb N}\newcommand{\Z}{\mathbb Z}\newcommand{\DD}{\mathbb D}
\newcommand{\HH}{\mathcal H}\newcommand{\A}{\mathscr A}\newcommand{\OO}{\mathcal O}
\newcommand{\m}{\mathfrak m}\newcommand{\st}{\operatorname{st}}
\newcommand{\supp}{\operatorname{supp}}\newcommand{\Res}{\operatorname{Res}}
\newcommand{\res}{\operatorname{res}}\newcommand{\wind}{\operatorname{wind}}
\newcommand{\ord}{\operatorname{ord}}\newcommand{\Tr}{\operatorname{Tr}}
\newcommand{\red}{\operatorname{red}}\newcommand{\Ind}{\operatorname{Ind}}
\newcommand{\id}{\operatorname{id}}\newcommand{\vH}{v_{\!H}}
\newcommand{\re}{\operatorname{Re}}
\newcommand{\HInt}{\mathop{\int^{\!H}}\nolimits}
\newcommand{\HCoint}{\mathop{\oint^{\!H}}\nolimits}
\newcommand{\PInt}{\mathop{\int^{\!\mathrm{pol}}}\nolimits}
\newcommand{\ACoint}{\mathop{\oint^{\!\mathrm{alg}}}\nolimits}
\newcommand{\dR}{\mathrm{dR}}
\newcommand{\sh}{^{\sharp}}\newcommand{\dd}{\,d}
\newcommand{\abs}[1]{\left|#1\right|}
\newcolumntype{L}[1]{>{\raggedright\arraybackslash}p{#1}}
\numberwithin{equation}{section}
\setcounter{tocdepth}{2}

\begin{document}
\hypersetup{pageanchor=false}
\pagenumbering{roman}
\begin{titlepage}
\centering
\vspace*{0.6cm}
{\Large\scshape A merged research synthesis and continuation\par}
\vspace{0.75cm}
{\Huge\bfseries Jordan Separation,\par Cauchy Theory,\par and Stokes Calculus\par}
\vspace{0.4cm}
{\LARGE on the Surcomplex Plane\par}
\vspace{0.6cm}
{\large Topologies, admissible chains, polynomial and deformed Stokes,\par
multiscale contours, radius-free residues, and untransformed Leray cycles\par}
\vspace{0.6cm}
{\large September 21, 2026\par}
\vspace{0.5cm}
\begin{minipage}{0.94\textwidth}
\begin{abstract}
A counterpart of a complex-analytic theorem over $\SC=\No[i]$ must specify three
things before it has a definite meaning: a topology, a regularity algebra for the
integrand, and a category of chains with an integration pairing. The full fine
topology admits no nonconstant continuous ordinary real-parameter paths; a fixed
Hahn workspace has a different, nondiscrete but totally separated intrinsic
topology. Neither supports the literal ordinary Jordan-curve picture.

We separate the question into four theories and prove each in its own category.
First, two Jordan separation theorems: a deliberately coarse one in the
standard-part topology, whose separator is the whole boundary monad rather than a
curve, and the established semialgebraic Jordan theorem over real closed fields,
which we import rather than prove. Second, three chain categories carrying an
exact generalized Stokes theorem: polynomial simplices over any
characteristic-zero field, Hahn differential forms on ordinary manifolds with the
cohomology comparison $H^q(\Omega_H^\bullet(M))\simeq H^q_\dR(M;\C)((t^\Gamma))$,
and infinitesimally deformed chains with a support-controlled Taylor pullback.
Third, Cauchy theory as a period theory: an exact endpoint-transport formula,
integer winding at displaced centers, a Cauchy formula on deformed Jordan
contours, and moving-pole residues. Fourth, a positive-monomial rescaling theorem
that converts radius-free Hahn analytic germs into polynomial-coefficient
families, and with it microscopic circle and torus realizations of the
perturbation residue, including a determinant-transformed realization for an
arbitrary isolated complete intersection. That realization left open the
question of a cycle for the \emph{untransformed} integrand, and the question is
then settled: a ramified finite covering of the coordinate torus, built from the
leading map and one positive valuation threshold alone, carries the original form
$H\dd z/(F_1\cdots F_n)$ to the perturbation residue for every perturbation above
that threshold and every radius-free numerator, with an exact equation-adapted
lift, a pole-free Hahn homotopy, a degree normalization, a finite-sheet trace,
and an extension to formal coefficient germs.

A rational contour pairing built from semialgebraic winding numbers is
characterized by a uniqueness property and compared with Hahn integration on
explicitly verified overlaps. Root-of-unity sampling recovers the integral
coefficientwise for common-support continuous data; a separate decay hypothesis
gives valuation convergence, which need not hold without it. Berkovich, tropical,
perfectoid, definable, tame, and resurgent integration theories are compared
without identifying their distinct targets.
\end{abstract}
\end{minipage}
\vfill
\begin{minipage}{0.94\textwidth}\small
\textbf{Status.} Established results, statements reported in the two supplied
manuscripts, and deductions proved here are distinguished throughout. This is not
a claim of a unique accepted foundation for surcomplex analysis, an exhaustive
priority determination, or a formally verified proof. Symbolic checks validate the
displayed finite examples and are not machine verification of the general proofs.
All sums and mathematical workspaces used in the constructions are set-sized. The
literature search was conducted on September 21, 2026, and was targeted rather
than exhaustive; no named published conjecture is claimed to be solved.
\end{minipage}
\end{titlepage}
\setcounter{page}{2}
\begingroup\small\setlength{\parskip}{0pt}\tableofcontents\endgroup
\clearpage
\pagenumbering{arabic}
\hypersetup{pageanchor=true}

\section{The question, the sources, and the principal conclusions}
\label{contours:sec:intro}

\subsection{Three theorems, three kinds of structure}
The three classical theorems in the title concern different structures. The
Jordan curve theorem describes separation by a topological embedding of a circle.
Cauchy's integral theorem describes the vanishing of periods of holomorphic
differential forms under a homological hypothesis. Generalized Stokes relates a
differential on forms to a boundary on chains. They cooperate in ordinary complex
analysis but are not logically interchangeable: even over $\C$, a Jordan curve
need not be rectifiable, a contour need not be simple for Cauchy's theorem to
apply, and being closed does not make a cycle null-homologous in a punctured
domain. One should therefore not try to extend the three theorems to
$\SC=\No[i]$ by a single replacement of $\C$ by a larger field.

Real closedness of $\No$ supplies algebraic closedness of $\SC$, but it does not
select a topology, a function algebra, or a chain category. The useful answer is
not one unrestricted transplantation but a collection of precisely compatible
counterparts. The classical starting point is the usual contour, Cauchy, and
differential-form theory \cite{Ahlfors,Spivak,Lee}.

The positive conclusions are these. A genuine Jordan analogue is available in
\emph{semialgebraic topology} over every real closed Hahn workspace. A second,
intentionally coarse Jordan statement is available after passing to standard
parts. An exact $\SC$-valued Stokes calculus is available in three separate chain
categories: polynomial simplices, Hahn-coherent chains over an ordinary base, and
infinitesimal deformations of the latter. Cauchy theory becomes a closed-form
period theory rather than a consequence of fine connectedness. Positive
monomial rescaling makes radius-free analytic germs accessible to contour methods,
and rational differentials admit a contour pairing based on algebraic winding
numbers even when an ordinary leading contour meets the reduction of a pole.
Finally, the untransformed integrand of an arbitrary isolated complete
intersection has a genuine single cycle, once the chain category is allowed the
ramified finite coverings of the coordinate torus that it always admitted; the
cycle is not a coordinate torus, and could not have been.
These answers coexist; they are not interchangeable.

\subsection{What the supplied manuscripts provide}
We refer to \emph{Surcomplex Analysis: Infinitesimal Calculus, Hahn-Coherent
Holomorphy, and Contour Theory} as the \emph{analysis manuscript} \cite{Analysis}.
It distinguishes fine-local power-series calculus from the algebra
\[
  \HH(U)=\OO(U)((t^\Gamma))
\]
of well-ordered Hahn families of ordinary holomorphic functions on one common
ordinary domain $U$. Its section on contour functionals defines integration
coefficientwise on ordinary paths and proves Cauchy, Morera, primitive, and
integral-estimate theorems; later sections prove moving-pole residue formulas,
permit coherent affine changes of scale, and construct a global fine-analytic
logarithm. These are existing results within the supplied framework.

The second source, \emph{Infinitesimal Analytic Geometry over the Surcomplex
Numbers: Radius-Free Hahn Germs, a Nullstellensatz, and Conservation of Singular
Zeros}, is the \emph{geometry manuscript} \cite{Geometry}. It replaces
common-domain coefficient functions by
\[
  \A_n=\C\{z_1,\ldots,z_n\}((t^\Gamma)),
\]
where each coefficient is a convergent complex germ but the coefficients need not
share a radius. For isolated complete intersections it defines a perturbation
residue by a strongly summable series of ordinary local residues, and it
explicitly does \emph{not} define that functional by a surcomplex contour
integral. The gap between a well-defined residue and a representing contour is one
principal topic below.

Neither manuscript establishes a Jordan theorem for arbitrary fine-topological
curves, nor generalized Stokes for arbitrary fine-smooth functions on arbitrary
surcomplex domains. We do not silently attribute either assertion to them. We
retain their distinction between ordinary coefficient regularity, strong
summability, and fine differentiation. This article is not an independent audit of
every theorem in either supplied manuscript.

\subsection{Established external counterparts}
There are relevant results beyond the supplied sources. Berarducci and Otero
develop definable intersection numbers, winding numbers, and Jordan--Brouwer
separation over o-minimal expansions of real closed fields
\cite{BO,BOhomology,BOmeasure}; Woerheide's o-minimal homology is the other route
to definable Jordan--Brouwer \cite{Woerheide}. Eisermann gives an algebraic
winding-number construction using Cauchy indices and Sturm chains over real closed
fields \cite{Eisermann}, with corrections and boundary cases treated by Perrucci
and Roy \cite{PR}. Ohmoto and Shiota give $C^1$ semialgebraic triangulations and a
Stokes formula, and explicitly allow a larger saturated target field \cite{OS};
Ma\v r\'{\i}kov\'a and Shiota measure definable sets with values in an ordered
semiring \cite{MS}; Kaiser develops tame semialgebraic integration over real closed
fields with Archimedean value group \cite{Kaiser}. Chambert-Loir and Ducros
construct a Stokes theory for real forms on Berkovich spaces \cite{CLD}, whose
underlying topology is developed by Baker--Rumely and Temkin \cite{BR,Temkin}.
Mihara constructs analytic singular homology and an integration theory over
cosimplicial perfectoid spaces \cite{Mihara}. Cherry and Diamond develop
Schnirelman-type and $p$-adic line integrals \cite{Cherry,Diamond}; a contemporary
use of such averages appears in \cite{Ettayb}. Ducros, Hrushovski, and Loeser
construct a comparison between complex and non-Archimedean integrals \cite{DHL}.
Costin and Ehrlich study surreal extension and integration of distinguished
classes of real functions \cite{CE}.

These works answer substantial parts of the question in their own categories. A
theorem about real-valued tropical forms is not already a theorem about
$K$-valued holomorphic contour integrals. A definably connected interval is not an
ordinarily connected interval. An integral valued in a saturated extension or an
ordered semiring is not automatically a canonical element of a previously fixed
Hahn field.

\subsection{A map of the answers}
\begin{center}\small
\begin{longtable}{@{}L{0.22\textwidth}L{0.38\textwidth}L{0.32\textwidth}@{}}
\toprule
Question & Valid counterpart & Essential qualification\\\midrule\endhead
Jordan separation & Definable circle in a real closed plane
  & Definably connected components, not ordinary ones\\[0.45em]
Coarse Jordan separation & Standard-part preimage of an ordinary Jordan curve
  & A thick wall in a non-Hausdorff topology\\[0.45em]
Generalized Stokes, algebraic & Polynomial simplices over any characteristic-zero
  field & A finite factorial identity, not a measure\\[0.45em]
Generalized Stokes, analytic & Hahn forms on ordinary or infinitesimally deformed
  chains & Support-controlled pullback and integration\\[0.45em]
Cauchy vanishing & Closed Hahn-holomorphic $1$-form on an admissible
  null-homologous contour & Coherent coefficient data, or a specified algebraic
  pairing\\[0.45em]
Contour deformation & Closed-form homotopy invariance; exact endpoint corrections
  & The whole homotopy must be admissible and avoid singular coefficients\\[0.45em]
Radius-free residues & Microscopic circle, or torus after positive monomial
  rescaling & A determinant transform for a general complete intersection\\[0.45em]
Untransformed residues & One Hahn cycle on a ramified finite cover of the torus,
  built from $f$ and one valuation threshold & A convergent isolated leading map
  and a positive support bound; the cycle is not a coordinate torus\\[0.45em]
Arbitrary rational loops & Residue functional weighted by algebraic winding
  & A definition with compatibility theorems, not a Riemann limit\\[0.45em]
Sampling a circle & Coefficientwise root-of-unity recovery, or a verified
  valuation-tail criterion & Coefficientwise, valuation, and fine convergence are
  not interchangeable\\
\bottomrule
\end{longtable}
\end{center}
The affirmative rows are theorems in specified categories, not approximations to a
nonexistent unrestricted fine-topological contour theory.

\subsection{Conventions, and one name for one concept}
\label{contours:subsec:conventions}
Three of the source manuscripts merged here use three different names for one and
the same topology, and several symbols are used inconsistently across them. We fix
the following conventions and use them without exception.

\begin{enumerate}
\item \textbf{The standard-part topology.} The initial topology of the
standard-part map on the finite elements --- whose open sets are exactly the
preimages of ordinary open subsets of $\C$, whose fibres are indiscrete, and whose
Hausdorff quotient is $\C$ --- is called the \emph{standard-part topology}
throughout. It appears in the literature merged here under the names
\emph{residue topology}, \emph{reduction topology}, and \emph{shadow topology};
these are synonyms for one and the same topology, and are not used again. The
descriptive name is preferred here because ``residue'' is already committed to the
residue functionals of \cref{contours:sec:radiusfree,contours:sec:multi} and
``shadow'' is used in the analysis manuscript for the ordinary shadow of a zero.
\item \textbf{One symbol per concept.} We write $\st$ for the standard-part
(reduction) map $K^\circ\to\C$ and never $\red$; $\m_K$ for the infinitesimals of
the workspace and $\m$ only for the class of all infinitesimals of $\SC$;
$U\sh=\st^{-1}(U)$ for the thickening of an ordinary set; $v$ for the valuation of
a field element and $\vH$ for the least exponent of a Hahn coefficient family;
$\wind$ for ordinary complex winding and $\wind_{\mathrm{sa}}$ for semialgebraic
(definable, algebraic) winding, never $\operatorname{Ind}$ or $w_F$; $\res$ for an ordinary local
residue and $\Res_F$ for the perturbation residue of \eqref{contours:eq:residueone}
and \eqref{contours:eq:generalres}.
\item \textbf{Scales.} Rescaling exponents are always written $\rho_j\in
\Gamma_{>0}$ and the corresponding radii $r_j=t^{\rho_j}$. The sources use
$\rho$, $\beta$, and $s$ for the exponent and $r$, $\rho$, and $r$ for the radius;
only the first convention appears below.
\item \textbf{Multiplicities.} Three senses are kept apart and never conflated:
the order $m\ge1$ of an ordinary one-variable germ at zero; the global colength
$\mu=\dim_\C\C\{z\}/(f_1,\ldots,f_n)$ of an isolated complete intersection; and the
local multiplicity $e_a=\dim_K B_a$ at an actual displaced zero, with
$\sum_ae_a=\mu$. The coordinate exponents in a coordinate-power system are written
$d_j$.
\item \textbf{Every theorem names its ring.} Four coefficient rings occur below
under nearly identical notation, and no theorem is transferred between them. They
are tabulated in \cref{contours:subsec:rings}.
\end{enumerate}

\subsection{What is merged here, and from where}
This article merges four manuscripts of September 21, 2026. Three of them share
roughly two thirds of their content: \emph{Stokes Calculus}, \emph{Stokes Integration}, and
\emph{Stokes Formulae on the Surcomplex Plane}. The shared material --- the
fine-topology no-go results, the standard-part Jordan theorem, the imported
semialgebraic Jordan theorem, the Hahn form complex and its cohomology, deformed
Stokes, the endpoint-displacement formula, the rescaling lemma, and the
one-variable microscopic circle --- is stated once.

The differences are real and sectional. The polynomial-chain theory of
\cref{contours:sec:polynomial}, the support-preserving Poincar\'e lemma
\cref{contours:prop:poincare}, the order-valued estimate
\cref{contours:prop:ml}, the sampling theory of \cref{contours:sec:sampling}, and
the scope discussion of \cref{contours:subsec:scope} come from the second
manuscript alone. The uniqueness characterization of the algebraic pairing
\cref{contours:thm:algpairing}, the two compatibility results
\cref{contours:prop:overlap,contours:prop:comparison}, the quantitative
finite-partition obstruction \cref{contours:prop:mesh}, and above all the
unconditional general contour representation
\cref{contours:lem:residuetransform,contours:thm:generalrepresentation} come from
the first. The numbered admissible-deformation definition
\cref{contours:def:deformation}, the independent separating-torus route of
\cref{contours:thm:septorus}, and the finite-parameter appendix
\cref{contours:app:transform} come from the third.  The second and the third
both number a winding-stability theorem; \cref{contours:thm:winding} is the
second's affine-scale form, which is the stronger of the two and is the one
printed.

Two of these overlap in a way that is a genuine disagreement about status rather
than about arithmetic, and one is a disagreement about emphasis. Both are resolved
in the open in \cref{contours:subsec:scope} and \cref{contours:sec:sampling}
respectively, rather than averaged away.

A fourth manuscript, \emph{Untransformed Leray Cycles for Radius-Free Hahn
Analytic Residues} \cite{Leray}, was written afterwards and against this
article, and is integrated here as
\cref{contours:sec:untransformed,contours:sec:adapted}. It
does not overlap the other three: it takes from them the perturbation residue
\eqref{contours:eq:generalres}, the positive-monomial rescaling discipline, the
coefficient rings of \cref{contours:subsec:rings}, the support lemma
\cref{contours:lem:support} with the sharpening recorded after
\cref{contours:ex:rank}, and the standing refusal to convert strong summability
into fine convergence; and it adds the construction that
\cref{contours:subsec:scope} and \cref{contours:subsec:problems} had asked for
by name. Everything it contributes is new to this article: the ramified finite
covering \cref{contours:thm:ramifiedleray} and the target-monomial direction
\cref{contours:lem:targetmonomial} that avoids the discriminant; substitution of
a positive-support Hahn map into a radius-free or formal germ,
\cref{contours:lem:radiusfreesub}, which is not
\cref{contours:thm:rescaling}; the rooted-tree implicit-function lemma
\cref{contours:lem:positiveimplicit}; the universal protecting cycle and its
period \cref{contours:thm:untransformed}; the exact adapted lift
\cref{contours:thm:adaptedlift} with the pole-free homotopy
\cref{contours:thm:untransformedhomotopy}; the direct degree normalization
\cref{contours:thm:degreenormalization}, the finite-sheet trace
\cref{contours:cor:protectedtrace} and the ideal annihilation
\cref{contours:prop:idealannihilation}; and the formal-coefficient extension
\cref{contours:thm:formalcontour}. Accordingly
\cref{contours:subsec:scope,contours:subsec:problems} are rewritten below as the
record of a question that was open and of the exact hypotheses under which it is
now closed --- not deleted, because a reader must still be able to see what was
asked and that \cref{contours:thm:generalrepresentation} answered only the
transformed form.

\section{Workspaces, supports, and the four coefficient rings}
\label{contours:sec:foundations}

\subsection{A set-sized workspace}
Write $\SC=\No[i]$ and use the surreal normal-form monomials $t^\gamma=
\omega^{-\gamma}$. Let $\Gamma\subset\No$ be a nonzero divisible set-sized ordered
additive subgroup and put
\begin{equation}\label{contours:eq:fields}
  R_\Gamma=\R((t^\Gamma)),\qquad
  K=\C((t^\Gamma))=R_\Gamma[i],\qquad
  t^\gamma=\omega^{-\gamma}.
\end{equation}
An element is a series whose exponent support is well ordered. Its valuation
$v(a)$ is its least exponent, with $v(0)=+\infty$. The Hahn-field theorem makes
$R_\Gamma$ real closed and $K$ algebraically closed; the algebraic-closure
statement is Corollary~4 of Poonen \cite{Poonen}. The normal-form interpretation
places both fields inside $\SC$; the surreal background and restricted analytic
evaluation are discussed by van den Dries and Ehrlich \cite{vdDE}. Put
\[
  K^\circ=\{a:v(a)\ge0\},\qquad
  \m_K=\{a:v(a)>0\},\qquad
  \st:K^\circ\longrightarrow\C,
\]
where $\st$ extracts the coefficient of $t^0$. Every finite element is uniquely
$a=c+\varepsilon$ with $c\in\C$ and $\varepsilon\in\m_K$. For ordinary
$U\subset\C$ write $U\sh=\st^{-1}(U)$. The modulus $|x+iy|=\sqrt{x^2+y^2}$ is
valued in $R_\Gamma$, not in $\R$; it satisfies the algebraic triangle inequality
and is not a real-valued non-Archimedean norm.

\emph{Enlargement rule.} Every finite collection of surcomplex points, and every
set-sized coefficient presentation used below, lies in a workspace of this kind,
obtained by adjoining the groups generated by the relevant normal-form supports and
taking the divisible hull. \emph{Its limit.} This does not authorize a sum over a
proper class, nor a proper-class polynomial ring, nor does it imply that every
rescaling required later already lies in a previously chosen $\Gamma$.
Differentiation in a variable $z$ holds every element of $K$ fixed; it is not a
derivation on $\No$ regarded as a transseries field. Full-class neighborhood
statements may be interpreted in a class theory such as NBG, while all rings,
supports, and parameter manifolds used in the integration constructions are
set-sized.

\subsection{Strong summability is algebra, not a hidden limit}
\begin{definition}[Strong summability]\label{contours:def:strongsum}
A family of Hahn series is \emph{strongly summable} when the union of its supports
is well ordered \emph{and} only finitely many members contribute to any one
exponent. Its sum is defined by finite addition at each exponent. For a well
ordered $E\subset\Gamma_{>0}$, let $E^*$ denote the additive monoid it generates,
including $0$.
\end{definition}

Both clauses are load-bearing, and neither implies the other.

\begin{lemma}[Support mechanism]\label{contours:lem:support}
If $S,T\subset\Gamma$ are well ordered, then $S+T$ is well ordered and each
exponent has only finitely many decompositions $s+t$. If $E\subset\Gamma_{>0}$ is
well ordered, then $E^*$ is well ordered and a fixed exponent has only finitely
many representations as a finite ordered word in $E$ --- including only finitely
many possible word lengths. The same conclusions hold for $S+E^*$.
\end{lemma}

These are the Hahn--Neumann facts used throughout \cite{Neumann,Poonen}. Three
consequences are used constantly. First, a complex-linear map $L:V\to W$ extends
coefficientwise by $L(\sum_\gamma v_\gamma t^\gamma)=\sum_\gamma L(v_\gamma)
t^\gamma$; this cannot enlarge support and requires no norm continuity. It will be
applied to ordinary integration, differentiation, residue extraction, and chosen
linear splittings of differential complexes. Second, for positive-support $A$ the
Neumann inversion $(1+A)^{-1}=\sum_{q\ge0}(-A)^q$ is strongly summable with output
support in $E^*$. Third, $S+E^*$ is the universal support certificate: every
construction below whose contributions are indexed by one source exponent in $S$
together with a finite word in $E$ is legitimate, and its output support lies in
$S+E^*$.

Well-ordering cannot be weakened to positivity: $\{1/r:r\ge1\}$ is a positive
subset of $\Q$ that is not well ordered, and $E^*$ would then fail to be well
ordered.

\begin{example}[Strong sums without valuation limits]\label{contours:ex:rank}
For $\Gamma=\Q+\Q\omega$, every integer is less than $\omega$. The strong sum
$\sum_{n\ge0}t^n=(1-t)^{-1}$ exists, but its finite partial sums do not converge to
it in the valuation topology: their nonzero tails never have valuation greater than
$\omega$. Consequently no argument of the form ``the valuation improves by a fixed
positive amount at each iteration'' justifies any of the constructions below.
\end{example}

A second negative companion is equally load-bearing, and is stated here once
because several proofs depend on it. \emph{``Only finitely many dependencies at a
given exponent'' does not mean ``only finitely many exponents below it.''} The
latter fails below $\omega$ in $\Q+\Q\omega$, where the interval of exponents below
$\omega$ contains every positive integer. It is the former statement, supplied by
\cref{contours:lem:support}, that legitimises the finite-parameter replacement
arguments in the residue transformation proofs of \cref{contours:sec:multi} and
\cref{contours:app:transform}. It is therefore unsafe to argue that ``truncating
below $\gamma$ is a finite calculation''; the correct statement is that a
\emph{fixed output exponent} has finite combinatorial dependency.

\subsection{The four coefficient rings, named once and never interchanged}
\label{contours:subsec:rings}
Four coefficient rings appear below under nearly identical notation. They carry
nearly identical theorem statements, and confusing them is the most likely source
of a false attribution in this subject. Every theorem in this article names the
ring it is stated over, and no theorem is transferred from one to another.

\begin{center}\small
\begin{longtable}{@{}L{0.20\textwidth}L{0.24\textwidth}L{0.46\textwidth}@{}}
\toprule
Ring & Notation & Defining requirement\\\midrule\endhead
Common-domain coherent & $\HH(U)=\OO(U)((t^\Gamma))$, $U\subset\C$ ordinary open
  & Every coefficient is holomorphic on \emph{one} common ordinary domain $U$, and
    all share \emph{one} well-ordered support. This is the fixed-domain (in one
    variable, fixed-disk) ring.\\[0.5em]
Polynomial-coefficient & $\C[w_1,\ldots,w_n]((t^\Gamma))$
  & Every Hahn coefficient is an ordinary polynomial; contained in, but
    strictly smaller than, the common-domain ring with $U=\C^n$, since
    $e^{w}$ is entire and is not a polynomial. This is the uniform-support target of
    rescaling, \cref{contours:thm:rescaling}.\\[0.5em]
Radius-free germs & $\A_n=\C\{z_1,\ldots,z_n\}((t^\Gamma))$
  & Every coefficient is an ordinary convergent germ at $0$, but their convergence
    radii may tend to zero, so \emph{no common polydisk is required}.\\[0.5em]
Formal & $K[[z]]$, and $\C[[z]]((t^\Gamma))$
  & No convergence of any coefficient is required. Used for injectivity, the formal
    residue class \eqref{contours:eq:formalresclass}, and the formal-coefficient
    contour extension \cref{contours:thm:formalcontour}.\\
\bottomrule
\end{longtable}
\end{center}

One variant of the first row is used later and is not a fifth ring but the same
requirement on a larger base: for a fixed ordinary complex \emph{manifold}
$\mathsf X$, possibly disconnected, $\HH(\mathsf X)=\OO(\mathsf X)((t^\Gamma))$
demands holomorphy of every coefficient on that one fixed $\mathsf X$ and one
common well-ordered support. It appears in
\cref{contours:sec:untransformed,contours:sec:adapted}, where $\mathsf X$ is a
finite covering of a polyannulus; see \eqref{contours:eq:HXring}. No theorem
stated over $\HH(U)$ for an ordinary open $U\subset\C$ is transferred to it.

The inclusions $\C[w]((t^\Gamma))\subset\HH(\C^n)$ and
$\A_n\subset\C[[z]]((t^\Gamma))$ are strict, and the containment
$\HH(U)\subset\A_n$ (for $U$ a polydisk about $0$, after taking germs) is strict as
well. The running witness for strictness of the last containment is
\begin{equation}\label{contours:eq:shrinking}
  H(z)=\sum_{j\ge1}\frac{t^j}{1-jz}\in\A_1 ,
\end{equation}
whose coefficient poles at $1/j$ accumulate at $0$, so that no fixed positive
ordinary radius about zero carries all of its coefficient germs. Integrating it on
one fixed ordinary circle as though all coefficients were holomorphic inside that
circle would be invalid. It nevertheless evaluates at every $z\in\m_K$. The remedy
is not a uniform ordinary bound but infinitesimal rescaling; see
\cref{contours:sec:radiusfree}.

A related warning is recorded here once. Changing the order of formal completion
and Hahn extension can change the domain of the resulting expression. For instance
\begin{equation}\label{contours:eq:twoexpansions}
 \frac1{z-t}=
 \begin{cases}
 \displaystyle\sum_{n\ge0}t^n z^{-n-1},& v(t)>v(z),\\[0.8em]
 \displaystyle-\sum_{n\ge0}t^{-n-1}z^n,& v(z)>v(t),
 \end{cases}
\end{equation}
and the second expression lies in $K[[z]]$ but \emph{not} in the unscaled
radius-free ring $\A_1$, since its combined coefficient support in $t$ is not well
ordered. The first is an outer expansion enclosing the displaced pole; the second
is a Taylor expansion on a smaller region not enclosing it. Their respective
contour residues about zero are $1$ and $0$, and there is no contradiction: the
admissible regions differ.

\subsection{The perturbation residue, and why it is not a contour integral}
\label{contours:subsec:residuedef}
Fix an ordinary analytic germ $f=(f_1,\ldots,f_n)$ at $0$ with isolated common
zero, and a perturbation $F=f+E$ with $E_j\in\A_n$ and $\vH(E_j)>0$ for each
nonzero $E_j$. Zero perturbations are permitted. For $H\in\A_n$ the geometry
manuscript defines
\begin{equation}\label{contours:eq:generalres}
  \Res_F(H)=\sum_{k\in\N^n}(-1)^{|k|}
  \res_{(f_1^{k_1+1},\ldots,f_n^{k_n+1})}
  \Bigl(H\prod_{j=1}^nE_j^{k_j}\Bigr),
\end{equation}
where $\res$ is the ordinary local Grothendieck residue, a complex-linear
functional on analytic germs, extended coefficientwise, and the normalization
$(2\pi i)^{-n}$ is included in $\res$. In one variable this reads
\begin{equation}\label{contours:eq:residueone}
  \Res_F(H)=\sum_{k\ge0}(-1)^k\res_{f^{k+1}}(HE^k),
  \qquad
  \res_{f^{k+1}}(q)=\Res_{z=0}\frac{q(z)\dd z}{f(z)^{k+1}} .
\end{equation}
The positive support of $E$ and \cref{contours:lem:support} prove strong
summability, with output support contained in $\supp H+(\supp E)^*$.

We state emphatically, once, what these formulas are. \emph{Equation
\eqref{contours:eq:generalres} is a series definition, not a surcomplex contour
integral.} The geometry manuscript defines it algebraically precisely because no
contour was available on $\A_n$. Supplying a contour that represents it is a
separate task, and is the subject of
\cref{contours:sec:radiusfree,contours:sec:multi,contours:sec:untransformed,contours:sec:adapted}.

\emph{One functional, defined in two companion reports by the same series.} The
report \emph{Finite Hahn Deformations: Division, Conservation of Multiplicity,
and Residue Duality} \cite{FiniteDeformations}, in this repository at
\texttt{docs/surcomplex/finite-deformations}, defines a Hahn local-cluster
residue by exactly the expansion \eqref{contours:eq:generalres} --- the same
alternating sum of ordinary local residues of $H\prod_jE_j^{k_j}$ against
$\prod_jf_j^{k_j+1}$, applied coefficientwise, with the same sign $(-1)^{|k|}$,
the same factor $(2\pi i)^{-n}$, the same ordered tuple $F_1,\ldots,F_n$ and the
same orientation convention. That report says emphatically that \emph{the series
is the definition}, that no fine-topological Riemann integral over a surcomplex
parameterized torus is intended, and that supplying a genuine representing
contour is a separate question left to this report. Its functional is stated over
the common-domain ring $\HH(D)$ for a fixed ordinary polydisk $D$, whereas
\eqref{contours:eq:generalres} is the radius-free version over $\A_n$; the two
formulas agree wherever both are defined, and neither is a residue invented to
make a contour identity tautological. \emph{This identification is the point of
the whole exercise}: every contour theorem below ---
\cref{contours:thm:microcircle,contours:thm:microtorus,contours:thm:generalrepresentation,contours:thm:septorus,contours:thm:untransformed}
--- earns its value by computing that independently defined functional and
nothing else. The companion report's residue duality, its perfect pairing, and
its conservation of multiplicity are results about the functional and are cited,
not reproved, here; the trace identity \eqref{contours:eq:traceresidue} below,
which \cref{contours:sec:multi} takes from the supplied geometry
manuscript, is among the results that report establishes as well. The sign $(-1)^{|k|}$ and the
factor $(2\pi i)^{-n}$ on the contour side are the chosen normalization; they are
fixed so that $\Res_F(J_F)=\mu$ and $\Res_F(H)=\sum_aH(a)/J_F(a)$ at simple zeros,
where $J_F=\det(\partial F_i/\partial z_j)$.

Three normalizations involving $2\pi i$ are kept deliberately separate throughout,
and are never silently identified.
\begin{enumerate}
\item \textbf{Hahn contour.} The microscopic circle functional carries the factor
$1/(2\pi i)$ and the torus functional $1/(2\pi i)^n$ outside the integral, so that
each equals the algebraically defined perturbation residue on the nose.
\item \textbf{Local one-variable.} The residue $\res_a(R\dd z)$ is the bare
Laurent or partial-fraction coefficient $c_{a,1}$, \emph{after cancellations} in
the local finite-meromorphic germ; it is not merely the residue of one ordinary
coefficient at $\st(a)$. The factor $1/(2\pi i)$ again appears outside.
\item \textbf{Formal or algebraic.} In \eqref{contours:eq:formalresclass} and in
the algebraic pairing \eqref{contours:eq:algebraicpairing} the factor $2\pi i$ is a
\emph{chosen} normalization matching ordinary complex integration, using the
distinguished copy of $\C$ inside $K$. It is not forced by ordered-field
arithmetic; over an abstract real closed field with no distinguished copy of $\R$
one would normalize formally or work with $I/(2\pi i)$.
\end{enumerate}

\section{Three topologies, and why the naive constructions fail}
\label{contours:sec:topologies}

\subsection{The full fine topology and its restriction}
The full fine neighborhoods in $\SC$ are the balls $B(a,r)=\{z:|z-a|<r\}$ for
\emph{every} positive surreal $r$. A statement about this topology is a
class-topological statement; it is not a claim that all fine-open classes form a
set. The following conclusions are proved in the analysis manuscript
\cite[\S3]{Analysis}; the short argument is repeated because its exact scope is
crucial.

\begin{proposition}[Full-class discreteness of sets]\label{contours:prop:fullfine}
Every set-sized subset of $\SC$ is closed and discrete in the full fine topology.
Every convergent net with a set-sized index set is eventually equal to its limit.
Every fine-continuous map from an ordinary connected real interval, or from an
ordinary circle, into $\SC$ is constant. In particular, the topology induced from
the full fine topology on a fixed set-sized field $K$ is discrete.
\end{proposition}
\begin{proof}
For a point $p$ and a set $A$, the nonzero numbers $|p-a|$, $a\in A$, form a
\emph{set} of positive surreals. The cut $\{0\mid |p-a| : a\in A,\ a\ne p\}$
supplies a positive surreal $\delta$ strictly smaller than all of them, so
$B(p,\delta)$ misses $A\setminus\{p\}$. This proves closedness and discreteness.
Applying the same argument to the set of nonzero distances of the values of a net
from its proposed limit, eventual membership in the resulting ball forces eventual
equality. Finally, the image of an ordinary connected parameter space is a set; a
connected subset of a discrete space has at most one point.
\end{proof}

This is \emph{not} a theorem that a Hahn field with its own valuation topology is
discrete. The radius supplied by the cut may lie outside the workspace. The usual
real-parameter unit-circle map is not fine-continuous; calling it a contour below
will therefore always mean an admissible parameterized integration object, never a
fine-continuous loop.

\subsection{The intrinsic topology of a fixed Hahn workspace}
The internal ordered-field topology on $K$ uses radii from $R_\Gamma$. It agrees
with the natural valuation topology: monomial valuation balls refine ordered balls,
and sufficiently small ordered monomial balls refine any prescribed valuation ball,
since $v(x)>v(r)$ implies $|x|<r$ for $r>0$, while $|x|<t^{\gamma+\eta}$ with
$\eta>0$ implies $v(x)>\gamma$.

\begin{proposition}[Intrinsic total separation]\label{contours:prop:intrinsic}
The intrinsic topology on $K$ is nondiscrete and totally separated. The cosets
$a+r\m_K$, $r\ne0$, are open and closed and separate any pair of distinct points.
Every connected subset is a singleton, every continuous map from an ordinary real
interval or ordinary circle to $K$ is constant, and ordinary singular homology of
$K$ in positive degrees vanishes.
\end{proposition}
\begin{proof}
Every neighborhood of zero contains a nonzero monomial of sufficiently large
exponent, so the topology is nondiscrete. The subgroup $r\m_K$ is the valuation
ball $\{x:v(x)>v(r)\}$; it is open, and its complement is the union of the other
open cosets, so it is clopen. For $a\ne b$ the coset $a+(b-a)\m_K$ contains $a$
and not $b$, so a connected subset has at most one point, and continuous images of
connected ordinary spaces are singletons. Every ordinary singular simplex is
therefore constant, and the singular chain complex is the direct sum over points of
the singular chain complex of a point, whose positive-degree homology vanishes.
\end{proof}

The same clopen separation works in the full fine topology, using $a+(b-a)\m$ with
the full class of infinitesimals. There is consequently no ordinary topological
embedding $S^1(\R)\hookrightarrow\SC$ in either topology, and deleting a putative
circle does not leave exactly two nontrivial ordinary connected components. The
nonzero periods constructed below therefore do \emph{not} come from that singular
chain complex.

\subsection{The standard-part topology}
On $K^\circ$ --- or on the affine arena $X_{a,r}=a+rK^\circ$ for $a\in K$ and
$r\in R_\Gamma$ positive, via the surjection
\begin{equation}\label{contours:eq:quotientmap}
  q_{a,r}:X_{a,r}\longrightarrow\C,\qquad
  q_{a,r}(z)=\st\bigl((z-a)/r\bigr)
\end{equation}
--- declare the open sets to be exactly the sets $q_{a,r}^{-1}(V)$ with
$V\subset\C$ ordinary open. This is the initial topology of the standard-part map,
the \emph{standard-part topology} of \cref{contours:subsec:conventions}. Its fibres
are indiscrete and it is not Hausdorff, but its quotient by indistinguishability is
exactly the ordinary complex plane. It deliberately forgets all infinitesimal
distinctions. Its role is limited but clean: it supports a coarse Jordan theorem,
not microscopic separation.

\begin{remark}[Three topologies must not be conflated]
\label{contours:rem:threetopologies}
The intrinsic topology of $K$ uses radii in $R_\Gamma$. The topology inherited by
the set $K$ from the full class $\SC$ also permits radii outside $R_\Gamma$ and is
discrete by \cref{contours:prop:fullfine}. The standard-part topology is coarser
still. A claim about convergence or connectedness in one of them is not
automatically a claim in either of the others.
\end{remark}

\subsection{Unrestricted fundamental theorems are impossible}
\begin{proposition}[No endpoint integration for all fine-local primitives]
\label{contours:prop:noftc}
There is no family of linear endpoint integral functionals $I_{a,b}$, defined for
all fine-locally analytic derivatives and arbitrary endpoints, satisfying
$I_{a,b}(H')=H(b)-H(a)$ for every fine-locally analytic $H$ on the full surcomplex
plane.
\end{proposition}
\begin{proof}
The class $\m$ of all infinitesimals and each of its additive cosets are fine-open
and fine-closed. Hence
\[
  q(z)=\begin{cases}0,&z\in\m,\\ 1,&z\notin\m\end{cases}
\]
is locally constant, so fine-locally analytic, with $q'=0$. At $a=0$, $b=1$ the
proposed identity gives $I_{0,1}(0)=q(1)-q(0)=1$, contradicting linearity.
\end{proof}

This is the obstruction proved in \cite[\S3]{Analysis}; it is \emph{not} an
obstruction to a smaller coherent calculus. A proposed Stokes theorem must restrict
either the forms, the chains, or the differential. The Hahn theory of
\cref{contours:sec:derham,contours:sec:deformed} restricts coefficient coherence;
the semialgebraic theory of \cref{contours:sec:jordan} restricts definability; the
polynomial theory of \cref{contours:sec:polynomial} restricts to finite algebraic
data.

The analysis manuscript gives a sharper diagnostic. It constructs a globally
fine-locally analytic logarithm $\mathcal L$ on $\SC^\times$ with
$\mathcal L'=1/z$, using a selected argument cut
\cite[section on global logarithms]{Analysis}. Its restriction to ordinary points
has that argument jump and is \emph{not} a coherent primitive on
$(\C^\times)\sh$. Accordingly the coherent period
$\HCoint_{|z|=1}\dd z/z=2\pi i$ of \cref{contours:sec:cauchy} does not contradict
the existence of $\mathcal L$. Applying the coherent endpoint theorem
\cref{contours:thm:endpoints} to this noncoherent global function would be an
invalid change of category; a valid fundamental theorem must restrict the primitive
and the integration category together.

\subsection{Finite meshes cannot resolve infinitesimal tolerances}
\begin{proposition}[Finite-partition obstruction]\label{contours:prop:mesh}
Let $0=x_0<x_1<\cdots<x_N=1$ be a partition in $R_\Gamma$ with an ordinary finite
number $N$ of subintervals. Its mesh is at least $1/N$. For $f(x)=x$ the right and
left endpoint sums differ by
\begin{equation}\label{contours:eq:riemannobstruction}
  \sum_{j=1}^N(x_j-x_{j-1})^2\ \ge\ \frac1N .
\end{equation}
Consequently ordinary finite tagged partitions cannot make all tag choices agree to
a fixed positive infinitesimal error on $[0,1]_{R_\Gamma}$.
\end{proposition}
\begin{proof}
The increments $\Delta_j=x_j-x_{j-1}$ are positive and sum to one, so at least one
is at least $1/N$. The difference of the two endpoint sums is $\sum_j\Delta_j^2$,
and the algebraic identity
\[
  N\sum_j\Delta_j^2-\Bigl(\sum_j\Delta_j\Bigr)^2
  =\sum_{j<k}(\Delta_j-\Delta_k)^2\ \ge\ 0
\]
proves \eqref{contours:eq:riemannobstruction}. For any positive infinitesimal
$\epsilon$ one has $1/N>2\epsilon$ for every ordinary $N$, so the two tag choices
cannot both lie within $\epsilon$ of one proposed integral.
\end{proof}

There is a related logical pitfall. A definition asserting that ``all partitions of
mesh less than $\delta$ have sums close to $I$'' becomes vacuous when $\delta$ is
infinitesimal: there are no such ordinary finite partitions. One must specify an
actual directed approximation system and prove it has enough elements. Note also
what is \emph{not} being claimed: a polynomial antiderivative still gives
$\int_0^1x\dd x=1/2$, and the obstruction is to a particular unrestricted
finite-mesh definition, not to all integration. Hyperfinite partitions in a
nonstandard universe can be useful, but they require an additional universe,
transfer structure, and a chosen relation between that universe and the surreal or
Hahn workspace. \emph{They are not supplied merely by the field operations in
$\No$}, and taking only the ordinary standard part of such an integral would also
discard the infinitesimal information sought here.

\subsection{A contour avoiding every pole need not have an admissible pullback}
\label{contours:subsec:boundarylayer}
Consider the ordinary unit-circle trace and $a=1+t$ with $t>0$ infinitesimal. The
function $1/(z-a)$ has no pole on that trace. Away from $z=1$ its Hahn expansion
begins with $1/(z-1)$, but at $z=1$ its value is $-t^{-1}$. It cannot have a single
Hahn expansion with continuous coefficient functions on the entire ordinary circle:
the negative-exponent coefficient would vanish at all other ordinary points but not
at $1$. For $a=1-t$ the same problem occurs.

Semialgebraic geometry distinguishes these two poles as outside and inside the full
unit circle respectively. A coefficientwise contour integral based on the original
macroscopic parameterization does not automatically do so. One needs a boundary-layer
change of chart, a finer scale, or the algebraic pairing developed in
\cref{contours:sec:algebraic}. Avoidance of actual poles is necessary, but uniform
admissibility of the pulled-back coefficients is an additional hypothesis. This
example recurs in \cref{contours:subsec:notconflated} as the explicit case where
the two compatibility results must not be conflated.

\section{Jordan separation: what survives, and in which sense}
\label{contours:sec:jordan}

\subsection{A thick-wall theorem in the standard-part topology}
\begin{theorem}[Jordan separation in the standard-part topology]
\label{contours:thm:stdpartjordan}
Let $J\subset\C$ be an ordinary Jordan curve with bounded interior $D$ and exterior
$E$. Fix $a\in K$ and $r\in R_\Gamma$ with $r>0$, and use the map
\eqref{contours:eq:quotientmap}. Then
\[
  X_{a,r}\setminus q_{a,r}^{-1}(J)
  = q_{a,r}^{-1}(D)\ \sqcup\ q_{a,r}^{-1}(E)
\]
has exactly two connected components in the standard-part topology, and both have
boundary $q_{a,r}^{-1}(J)$ relative to $X_{a,r}$. In the case $a=0$, $r=1$ this
reads $K^\circ\setminus J\sh=D\sh\sqcup E\sh$ with common boundary $J\sh$.
\end{theorem}
\begin{proof}
For any subset $B\subset\C$ the subspace $q^{-1}(B)$ has precisely the saturated
open sets pulled back from $B$, since every open set of the standard-part topology
is saturated for the fibres of $q$ and every fibre is nonempty. A separation of
$q^{-1}(B)$ therefore descends to a separation of $B$, and conversely; so
$q^{-1}(B)$ is connected exactly when $B$ is. The ordinary Jordan theorem gives
$\C\setminus J=D\sqcup E$ with both sets nonempty and connected, whence $q^{-1}(D)$
and $q^{-1}(E)$ are connected. They are disjoint relatively clopen subsets covering
the complement, so they are exactly its components. Finally
\[
  \overline{q^{-1}(B)}=q^{-1}(\overline B),\qquad
  \operatorname{int}\bigl(q^{-1}(B)\bigr)=q^{-1}(\operatorname{int}B),
\]
directly from the neighborhood definition and surjectivity of $q$, and the ordinary
boundary identities give the last assertion.
\end{proof}

The removed object is the entire \emph{boundary tube}
$q^{-1}(J)=a+r(J+\m_K)$, the monad of every point of $J$, and not a
one-dimensional embedded curve or the thin trace $\{a+r\gamma(s):s\in S^1(\R)\}$ of
one parameterization. A continuous coarse path whose endpoints have standard parts
on different sides meets this thick wall. Several things are \emph{not} asserted.
Two infinitesimally separated points in the same boundary monad cannot be assigned
different sides by this construction; different points in a fibre can lie on
opposite sides of a finer-scale separator, and a finer scale is needed to
distinguish two zeros in one standard-part fibre. Neither component is being called
fine-connected. And ``exterior'' here means that the standard-part image is
ordinary unbounded: all of $K^\circ$ is bounded by any positive infinite element of
$R_\Gamma$, so we have \emph{not} obtained an unbounded exterior in the internal
ordered-field sense. This is an exact topological theorem, but it recovers only the
geometry visible at the selected scale.

\begin{example}[A boundary layer cannot be discarded silently]
\label{contours:ex:boundarylayer}
At $a=0$, $r=1$, both $1-t$ and $1+t$ reduce to the point $1$ on the unit circle,
so neither belongs to the thickened disk $\DD\sh=q^{-1}(\DD)$, although
$|1-t|<1<|1+t|$. The theorem's interior is \emph{not} the ordered-radius ball
$\{|z|<1\}$, and a Cauchy formula whose hypothesis is $z\in\DD\sh$ does not
automatically cover $z=1-t$.
\end{example}

This construction remains available at every infinitesimal or infinite affine
scale, and it explains the topology underlying the ordinary-base coherence of
\cite{Analysis} --- without asserting that the coefficient functions of a coherent
family are continuous for any of the other topologies under discussion.

\subsection{The semialgebraic replacement}
For this subsection let $R$ be a real closed field containing $\R$ and $K=R[i]$;
the case of interest is $R=R_\Gamma$. A set is \emph{semialgebraic} if it is
described by a finite Boolean combination of polynomial equations and inequalities
with parameters in $R$. A map is semialgebraic when its graph is. A set is
\emph{semialgebraically connected} (equivalently, definably connected) when it has
no separation into two nonempty relatively open semialgebraic subsets. Paths now
have parameter domain the \emph{field} interval
\[
  [0,1]_R=\{s\in R:0\le s\le1\},
\]
not just the ordinary real points in it. Definable connectedness and definable path
connectedness agree for semialgebraic sets; this is part of the standard
finite-topology theory \cite{BCR,BOhomology}.

\begin{theorem}[Semialgebraic Jordan theorem; classical input]
\label{contours:thm:sajordan}
Let $j:S^1(R)\to R^2$ be a semialgebraic homeomorphism onto its image $J$, where
$S^1(R)=\{(x,y):x^2+y^2=1\}$. Then $R^2\setminus J$ has exactly two
semialgebraically connected components. They are semialgebraically path connected,
one is bounded by an element of $R$ and the other is unbounded, and $J$ is the
common boundary of both. Every semialgebraic continuous path between the two
components meets $J$. More generally, smooth definable Jordan--Brouwer separation
holds in the o-minimal manifold framework of Berarducci--Otero.
\end{theorem}

\emph{We import this theorem, rather than claim to prove arbitrary topological
Jordan separation by elementary first-order transfer.} One route uses compatible
finite triangulation and the combinatorial separation theorem for an embedded
polygonal circle on a triangulated sphere: compactify the ambient plane
semialgebraically, triangulate the pair consisting of the ambient sphere and the
curve, observe that the curve is a finite subcomplex which is a circle and that the
adjacent two-dimensional simplices split into two sides, and compute connectivity
by chains of incident simplices. The component containing the point at infinity is
the exterior; the other has bounded realization; local incidence at the curve gives
the common-boundary conclusion. A path missing the curve cannot move from one
relatively open definable side to the other, because its parameter interval is
definably connected. The finite simplicial incidence data do not depend on which
real closed field realizes them, nor on the Archimedean property of the
coefficients. Alternatively, the definable Jordan--Brouwer theorem of
Berarducci--Otero supplies the smooth-manifold case directly, and Woerheide's
o-minimal homology the general one; the semialgebraic triangulation framework
handles the topological formulation \cite{BCR,BO,BOhomology,Woerheide,OS}. The
technical finite-topology input is being imported, not replaced by an informal
appeal to transfer of arbitrary continuous maps.

Several qualifications are essential. The theorem does \emph{not} say that those
components are connected in the unrestricted intrinsic topology: by
\cref{contours:prop:intrinsic} all their ordinary connected components are
singletons, and the nondefinable cuts that disconnect the ordered field are
excluded from the definition of semialgebraic connectedness. This is why
\cref{contours:thm:sajordan} does not contradict \cref{contours:prop:fullfine}:
$[0,1]_R$ is definably connected but need not be connected as an ordinary
topological space. ``$R$-bounded'' does not mean bounded by an ordinary real
number. The generalization in the last sentence refers to the differentiability and
definability hypotheses of the o-minimal framework; it is not a statement about
arbitrary class functions or arbitrary fine-topological subsets. And definable
results concern a finite description and its real-closed extensions; they do not
supply ordinary singular homology of a proper class. Finally, a semialgebraic
problem involving full surreal parameters may always be placed in a real closed
set-sized workspace containing those parameters, which makes this a genuine
separation theorem for surcomplex data.

\begin{example}[A microscopic circle with a genuine definable inside]
\label{contours:ex:definablecircle}
For $a\in K$ and $r\in R_\Gamma$, $r>0$, the circle $C_{a,r}=\{z\in K:|z-a|=r\}$
has semialgebraic inside $|z-a|<r$ and outside $|z-a|>r$. This remains true for
$r=t^\rho$, for an infinite $r$, and for arbitrary $a$. The two points
$a+r(1-t^\eta)$ and $a+r(1+t^\eta)$, $\eta>0$, lie on opposite definable sides
although their normalized standard parts are equal.
\end{example}

A distinction that will recur: the \emph{full semialgebraic circle} $C_{a,r}$, the
\emph{coarse thick wall} $q^{-1}(J)$ of \cref{contours:thm:stdpartjordan}, and the
\emph{ordinary-parameter contour trace} $\{a+re^{i\theta}:\theta\in[0,2\pi]_\R\}$
are three different objects, even when they are informally written with the same
equation. The last is generally only a small subset of the first, and its
parameterization is not fine-continuous.

\subsection{An integer-valued index without a contour integral}
For a semialgebraic loop $\gamma$ avoiding $b$, define its index as the definable
degree of
\[
  s\longmapsto\frac{\gamma(s)-b}{|\gamma(s)-b|}\quad\text{into }S^1(R),
\]
equivalently the corresponding definable homology invariant. With positive
orientation it is $1$ inside a semialgebraic Jordan curve and $0$ outside, it takes
values in $\Z$, and it is invariant under definable homotopies avoiding $b$. We
write it $\wind_{\mathrm{sa}}(\gamma,b)$. \emph{This is definable degree, not the
ordinary topological degree of an ordinary real-parameter map into the intrinsic
field topology.}

For piecewise polynomial or rational parameterizations there is a computationally
explicit substitute. After a generic fixed rotation of the coordinate axes, write
$\gamma-b=x+iy$ on each piece and form a Cauchy index; the rotation can be chosen
so that no piece has identically zero $y$. With the convention that a jump from
$-\infty$ to $+\infty$ contributes $+1$,
\begin{equation}\label{contours:eq:cauchyindex}
  \wind_{\mathrm{sa}}(\gamma,b)
  =\frac12\sum_{\text{pieces}}\operatorname{Ind}\Bigl(\frac{x}{y}\Bigr),
\end{equation}
with the usual half-endpoint conventions and subdivision at exceptional points; a
constant piece contributes zero and may be removed first. Sturm chains compute
these indices using field arithmetic and sign tests. Eisermann develops this over
arbitrary real closed fields \cite{Eisermann}; Perrucci and Roy supply corrected
product formulae and a systematic treatment of boundary cases \cite{PR}. Below,
every pole is disjoint from the loop, so boundary conventions are not needed, and
fractional conventions for zeros on an edge or vertex are never interpreted as
contour integrals through poles.

Three limitations are recorded. First, this degree is an ordinary integer, not a
surreal number of turns. Second, \emph{it does not by itself define the integral of
a general differential form}: the primitive of a semialgebraic function need not be
semialgebraic, and a chosen analytic expansion of a real closed field does not
automatically contain all the radius-free Hahn germs of $\A_n$. Third, the
availability of a finite algebraic prescription does not mean arbitrary surreal
inputs have computable normal forms or decidable sign tests; finite Sturm and sign
algorithms use an ordered-field arithmetic oracle, and the construction becomes an
algorithm only when the coefficient field and support arithmetic have effective
presentations.

\subsection{What Jordan separation does not supply}
Separation by itself does not define an integral. Even in a definable setting,
integrating a definable function can require new functions or a larger target
object (see \cref{contours:sec:alternatives}). Conversely, a useful period theory
can be defined on cycles with self-intersections without invoking Jordan separation
at all. Below, the standard-part theorem interprets ordinary-base domains, the
definable theorem interprets algebraic winding, and parameterized chains prove
Stokes. \emph{These are compatible uses, but not identifications of categories.}

\section{A completely algebraic Stokes theorem for polynomial chains}
\label{contours:sec:polynomial}

There is a direct counterpart of generalized Stokes that needs neither ordinary
fine paths nor infinite summation, and that lives over an arbitrary
characteristic-zero coefficient field rather than over any of the four rings of
\cref{contours:subsec:rings}. It also supplies a compatibility test for the Hahn
construction of \cref{contours:sec:derham}.

\subsection{Polynomial integration on a simplex}
For a characteristic-zero field $L$, \emph{define} a linear functional on
polynomials on the standard $n$-simplex by
\begin{equation}\label{contours:eq:simplexint}
  \PInt_{\Delta^n}u_1^{\nu_1}\cdots u_n^{\nu_n}\dd u_1\wedge\cdots\wedge \dd u_n
  =\frac{\nu_1!\cdots\nu_n!}{(|\nu|+n)!},
\end{equation}
extended $L$-linearly. The geometric notation $\Delta_R^n=\{u_i\ge0,\ \sum u_i\le1\}$
is available when $L$ contains an ordered field $R$, but formula
\eqref{contours:eq:simplexint} \emph{is} the definition, and the functional is
therefore prior to any point set. The $0$-simplex functional is evaluation. Give
the simplex the ordered vertices $(0,e_1,\ldots,e_n)$ and its boundary the
alternating face orientation.

A \emph{polynomial simplex} $\sigma:\Delta_R^n\to R^d$ is a polynomial map with
field coefficients, and for a polynomial $n$-form $\alpha$ with coefficients in $L$
we set $\PInt_\sigma\alpha=\PInt_{\Delta^n}\sigma^*\alpha$. Only finitely many
algebraic operations are involved.

\begin{theorem}[Polynomial-chain Stokes]\label{contours:thm:polstokes}
For every polynomial simplex $\sigma$ and every polynomial $(n-1)$-form $\alpha$,
\begin{equation}\label{contours:eq:polstokes}
  \PInt_\sigma d\alpha=\PInt_{\partial\sigma}\alpha .
\end{equation}
The identity extends to finite polynomial chains, and it holds over every
characteristic-zero coefficient field $L$ --- in particular over every surcomplex
workspace $K=\C((t^\Gamma))$ and over $\SC$ itself for finite data.
\end{theorem}
\begin{proof}
Pullback commutes with $d$ by the polynomial chain rule, so it suffices to prove
the identity on the standard simplex. Fix bounds on the degrees of the coefficients
of $\alpha$. Using \eqref{contours:eq:simplexint} and the affine face maps, both
sides are polynomials with rational coefficients in the finitely many coefficients
of $\alpha$. Over $\R$, formula \eqref{contours:eq:simplexint} is the ordinary
monomial simplex integral, so ordinary Stokes proves their equality for every
assignment of real coefficients. A polynomial over $\Q$ vanishing at every real
tuple is zero, by induction on the number of variables. Hence the identity is a
rational polynomial identity and holds over any characteristic-zero field. All
boundary signs and face substitutions are finite, so linearity proves the chain
assertion.
\end{proof}

\emph{This proof uses classical Stokes to certify a finite algebraic identity; it
does not transfer a statement quantifying over arbitrary integrable functions.}
Equivalent direct proofs expand the monomial factorials and cancel adjacent faces.

\begin{example}[Infinite width, infinitesimal height]\label{contours:ex:triangle}
Let $L,H>0$ be arbitrary elements of $R_\Gamma$ and orient the triangle $T$ with
vertices $0,L,iH$ counterclockwise. Since
$d(\bar z\dd z)=d\bar z\wedge \dd z=2i\dd x\wedge \dd y$, polynomial Stokes gives
\begin{equation}\label{contours:eq:triangle}
  \PInt_{\partial T}\bar z\dd z=iLH .
\end{equation}
With $L=\omega$ and $H=\omega^{-1}$ the answer is the ordinary finite $i$, despite
the extreme aspect ratio; the signed area functional is $LH/2$. No ordinary
compactness of the entire field-valued triangle is asserted.
\end{example}

The construction also handles polynomial forms on triangulated polytopes whenever
oriented internal faces cancel. It does \emph{not} define an integral for every
semialgebraic function, nor for a rational form whose denominator vanishes
somewhere on the chain; rational contour periods will require residues or
additional analytic data, which is exactly what \cref{contours:sec:algebraic}
supplies.

\section{A Hahn de Rham complex and coefficientwise Stokes}
\label{contours:sec:derham}

\subsection{The differential graded algebra of globally supported Hahn forms}
Let $M$ be an ordinary smooth manifold, with corners allowed when integration is
used, and let $\Omega^q(M;\C)$ denote its ordinary smooth complex-valued
differential forms. Define
\begin{equation}\label{contours:eq:hahnforms}
  \Omega_H^q(M)=\Omega^q(M;\C)((t^\Gamma)).
\end{equation}
An element is a family $\alpha=\sum_{\lambda\in S}\alpha_\lambda t^\lambda$ with
\emph{one global well-ordered support} $S$ and ordinary smooth coefficient forms;
smoothness up to a boundary means smooth extension in ordinary local corner charts.
The word ``globally supported'' is load-bearing: a single well-ordered support
serves the whole manifold. \emph{This is a space of coefficient data over an
ordinary manifold; it is not defined as the collection of all fine-smooth forms on
a surcomplex point set. A derivative acts on $\alpha_\lambda$, not on $t^\lambda$.}
Put
\[
  d\alpha=\sum_\lambda(d\alpha_\lambda)t^\lambda,
  \qquad
  \alpha\wedge\beta=\sum_\nu\Bigl(\sum_{\lambda+\mu=\nu}
    \alpha_\lambda\wedge\beta_\mu\Bigr)t^\nu .
\]
The inner sum is finite by \cref{contours:lem:support}, so
$\Omega_H^\bullet(M)$ is a differential graded algebra over $K$ with $d^2=0$ and
the graded Leibniz rule
$d(\alpha\wedge\beta)=d\alpha\wedge\beta+(-1)^q\alpha\wedge d\beta$ for $\alpha$
of degree $q$. Ordinary smooth pullback is coefficientwise, commutes with $d$, and
does not enlarge support. This $d$ differentiates ordinary variables while holding
every element of $K$ fixed.

For a compact oriented ordinary $q$-manifold $C$ with corners and a smooth map
$\phi:C\to M$, define
\begin{equation}\label{contours:eq:chainintegral}
  \HInt_{(C,\phi)}\alpha
  =\sum_{\lambda\in S}t^\lambda\int_C\phi^*\alpha_\lambda ,
\end{equation}
whose support is contained in $S$. Extend by finite linearity to finite piecewise
smooth chains. Finite subdivision, orientation reversal, ordinary
reparameterization, and cancellation of oppositely oriented faces have their usual
meanings. \emph{There is no infinite interchange justified by an unstated analytic
bound}: the operation is the coefficientwise extension of an ordinary linear
functional, and the output lies in $K$ even when no real-valued bound on the entire
coefficient family is available. Noncompact chains could be admitted with suitable
support conditions, but compact chains suffice for everything proved here.

\begin{theorem}[Hahn-valued generalized Stokes]\label{contours:thm:hahnstokes}
Let $C$ be a compact oriented ordinary $q$-manifold with corners, with its boundary
orientation given by the outward-normal-first convention, let $\phi:C\to M$ be
smooth, and let $\alpha\in\Omega_H^{q-1}(M)$. Then
\begin{equation}\label{contours:eq:hahnstokes}
 \boxed{\qquad
 \HInt_{(\partial C,\ \phi|_{\partial C})}\alpha
 =\HInt_{(C,\phi)}d\alpha .
 \qquad}
\end{equation}
The pairing is $K$-linear in the form and additive in the chain, the same identity
holds for finite piecewise smooth chains with their chain boundary, and the
integral is invariant under ordinary orientation-preserving smooth
reparameterization.
\end{theorem}
\begin{proof}
For each exponent $\lambda$, ordinary Stokes \cite[Thm.~4.6.1]{GH}, \cite{Lee}
gives
\[
  \int_{\partial C}(\phi|_{\partial C})^*\alpha_\lambda
  =\int_Cd(\phi^*\alpha_\lambda)
  =\int_C\phi^*(d\alpha_\lambda).
\]
Both Hahn sums have support in one well-ordered containing set, so equality of
every coefficient proves \eqref{contours:eq:hahnstokes}. For $A=\sum_\mu a_\mu
t^\mu\in K$, the coefficient of $A\alpha$ at a fixed target exponent is a finite
convolution, so ordinary complex linearity of integration applies to each finite
sum, proving $K$-linearity. For a finite chain, internal faces have equal
pulled-back forms and opposite signs and cancel coefficientwise. Ordinary change of
variables at each coefficient proves reparameterization invariance.
\end{proof}

The ordinary Stokes theorem, with its chain and orientation conventions, is the
classical input \cite{Lee,GH}. The contribution of
\cref{contours:thm:hahnstokes} is the precise Hahn coefficient setting and its
compatibility with strong summability. Its content is \emph{not} an assertion that
the full surcomplex plane is an ordinary manifold: the chain is ordinary, while its
coefficients and its integral may be genuinely nonstandard.

\subsection{A support-preserving Poincar\'e lemma}
On an ordinary star-shaped open $U\subset\R^d$ containing $0$, define the radial
homotopy operator, for $q\ge1$, by
\begin{equation}\label{contours:eq:radial}
  (h\alpha)_x(v_1,\ldots,v_{q-1})
  =\int_0^1 s^{q-1}\alpha_{sx}(x,v_1,\ldots,v_{q-1})\dd s ,
\end{equation}
and apply it to each Hahn coefficient. Ordinary smoothness at $s=0$ makes the
result an ordinary smooth form, and the Hahn support does not increase.

\begin{proposition}[Support-preserving Poincar\'e lemma]
\label{contours:prop:poincare}
The coefficientwise extension of \eqref{contours:eq:radial} to
$\Omega_H^\bullet(U)$ satisfies $dh+hd=\id-c_0^*$, where $c_0$ is the constant map
to $0$. Every closed Hahn form of positive degree on $U$ is exact, with a primitive
supported in the support of the original form. A Hahn function with zero
differential on a connected ordinary manifold is a single constant in $K$.
\end{proposition}
\begin{proof}
Differentiate the pullback under $x\mapsto sx$ with respect to $s$ and use the
ordinary chain rule; integrating from $0$ to $1$ gives the displayed homotopy
identity for each coefficient, equivalently the ordinary fundamental theorem applied
to the radial homotopy. Every operation is complex-linear, so its coefficientwise
extension preserves support. In positive degree $c_0^*=0$. In degree zero, a zero
differential makes each smooth coefficient locally constant, hence constant on a
connected ordinary base.
\end{proof}

The explicitness of the operator matters: it produces a primitive with no support
enlargement whatever, which is what later support certificates require.

\subsection{Cohomology and periods}
\begin{theorem}[Cohomology of the globally supported Hahn complex]
\label{contours:thm:cohomology}
For an ordinary smooth manifold $M$ there is a canonical $K$-linear isomorphism
\begin{equation}\label{contours:eq:cohomology}
  H^q\bigl(\Omega_H^\bullet(M),d\bigr)
  \ \simeq\ H^q_{\dR}(M;\C)((t^\Gamma)).
\end{equation}
In particular a closed Hahn form is exact if and only if each of its ordinary
coefficients is exact. If the ordinary Betti number $b_q$ is finite, the right side
is isomorphic to $K^{b_q}$ and naturally to $K\otimes_\C H^q_\dR(M;\C)$.
Integration on ordinary cycles extends the usual de Rham period pairing, and when
ordinary homology in degree $q$ is finite dimensional that pairing is nondegenerate
over $K$.
\end{theorem}
\begin{proof}
Write $Z^q=\ker d$ and $B^q=\operatorname{im}d$ in the ordinary complex. A Hahn
form is closed exactly when every coefficient lies in $Z^q$; send its class to the
Hahn series of ordinary cohomology classes, discarding zero coefficients. This has
support contained in the original support, and exact Hahn forms map to zero.
Choose a complex-linear section of $Z^q\to Z^q/B^q$; its coefficientwise extension
lifts every Hahn series of classes without enlarging support, proving surjectivity.
Choose also a complex-linear right inverse of $d$ on $B^q$; if all coefficient
classes vanish, applying it coefficientwise produces a Hahn primitive with support
in the same containing set, so the kernel is exactly the exact Hahn forms. The
auxiliary choices establish existence; the induced isomorphism is the canonical
coefficient-class map. Finite convolution proves $K$-linearity. For finite $b_q$,
expansion in a fixed ordinary basis identifies Hahn families with scalar extension
by $K$, a finite union of supports being well ordered; the ordinary de Rham theorem
supplies a nonsingular finite period matrix against a basis of ordinary homology,
and the same matrix is nonsingular over $K$.
\end{proof}

The finite-dimensional qualification matters and is not a technicality. For an
infinite-dimensional $V$, the space $V((t^\Gamma))$ generally contains families
whose coefficients span infinitely many independent directions, whereas
$K\otimes_\C V$ consists of finite sums of pure tensors. \emph{One must not
silently replace the first object by the second.} Nor have we claimed that the
global uniform-support assignment \eqref{contours:eq:hahnforms} is a sheaf for
arbitrary infinite covers: unlike holomorphic coefficient functions on a connected
domain, smooth coefficients can vanish on entire open sets, and locally chosen
supports need not have well-ordered union. For example, if $\psi_j$ are smooth
bumps supported near the positive integers, then $\sum_{j\ge1}t^{-j}\psi_j(x)$ is
locally a finite Hahn family on $\R$ but has global exponent set $\{-j:j\ge1\}$,
which is not well ordered, so it does not belong to
\eqref{contours:eq:hahnforms}. All chain integrals and homotopies below have
explicitly common supports or involve finitely many charts.

On an arbitrary ordinary manifold the classical de Rham theorem translates
\eqref{contours:eq:cohomology} into a period criterion: a closed Hahn form of
positive degree has a global Hahn primitive exactly when all its periods over
ordinary smooth cycles vanish \cite{Lee}.

\begin{example}[The surviving puncture class]\label{contours:ex:puncture}
On $\C^\times$ the class of $\dd z/z$ spans $H^1_\dR(\C^\times;\C)$, normalized by
its ordinary unit-circle period $2\pi i$; hence
$H^1(\Omega_H^\bullet(\C^\times))\simeq K\,[\dd z/z]$. This is ordinary-base
cohomology with enlarged coefficients. It is \emph{not} the positive-degree
singular cohomology of the totally separated intrinsic $K^\times$, which vanishes
by \cref{contours:prop:intrinsic}: the two complexes have different bases and
different chains, and there is no contradiction.
\end{example}

\subsection{Green, divergence, and area}
For a Hahn one-form $P\dd x+Q\dd y$ on an admissible oriented planar $2$-chain
$\Sigma$, \cref{contours:thm:hahnstokes} reads
\begin{equation}\label{contours:eq:green}
  \HInt_{\partial\Sigma}(P\dd x+Q\dd y)
  =\HInt_\Sigma(Q_x-P_y)\dd x\wedge \dd y ,
\end{equation}
and in complex notation, for a $C^1$ Hahn coefficient family $F$,
\begin{equation}\label{contours:eq:cauchygreen}
  \HInt_{\partial\Sigma}F\dd z
  =2i\HInt_\Sigma\partial_{\bar z}F\dd x\wedge \dd y .
\end{equation}
These are exact identities for the specified pullbacks. They imply Cauchy's theorem
for coherent holomorphic coefficients, but they also apply to nonholomorphic
coefficients. In three dimensions the same construction applied to a Hahn family of
vector fields $X$ with $d(\iota_X\mathrm{vol})=(\operatorname{div}X)\mathrm{vol}$
gives the divergence identity; defining flux by this form avoids introducing an
unproved independent theory of surreal surface area and unit normals.

\begin{example}[An infinitesimal rectangle has nonzero area]
\label{contours:ex:rectangle}
Let $A,B\in R_\Gamma$ be positive and parameterize a rectangle with side lengths
$A,B$ by the ordinary unit square. For $\alpha=-y\dd x+x\dd y$ one has
$d\alpha=2\dd x\wedge \dd y$; the right edge contributes $AB$, the top edge
contributes $AB$, and the other edges contribute zero, so
\begin{equation}\label{contours:eq:rectangle}
  \HInt_{\partial([0,A]\times[0,B])}\alpha=2AB
  =\HInt_{[0,A]\times[0,B]}d\alpha .
\end{equation}
The chain notation denotes the indicated affine parameterization; it does not
presume a fine-continuous ordinary square. For $A=t$ and $B=t^\omega$ the value is
$2t^{1+\omega}\ne0$, whereas the standard part alone would be zero, losing the
geometric content of the exact identity.
\end{example}

\section{Infinitesimally deformed chains, homotopy, and estimates}
\label{contours:sec:deformed}

\subsection{Admissible deformations and the Taylor-jet pullback}
\begin{definition}[Admissible positive-support deformation]
\label{contours:def:deformation}
Let $M$ be an ordinary \emph{compact} smooth parameter manifold with corners and
let $U\subset\R^d$ be ordinary open. An \emph{admissible deformation} of a smooth
map $\phi_0:M\to U$ is
\begin{equation}\label{contours:eq:deformedmap}
  \Phi=\phi_0+\eta,\qquad
  \eta=\sum_{e\in E}\eta_e t^e ,
\end{equation}
where $E\subset\Gamma_{>0}$ is \emph{one} well-ordered set, so the displacement has
strictly positive support, and the $\eta_e$ are smooth real vector-valued
coefficient functions on $M$, extending to a neighborhood in each corner chart. For
complex coordinates use real and imaginary parts; piecewise smooth
parameterizations are handled by finite subdivision.
\end{definition}

At each ordinary parameter value $\eta$ is an infinitesimal vector. \emph{The
coefficient maps need not satisfy a uniform ordinary norm bound}, and \emph{no
assertion of fine continuity of the real-parameter map $\Phi$ is intended}; indeed
a nonconstant $\phi_0$ already prevents such continuity by
\cref{contours:prop:fullfine}. The notation $\Phi$ describes an admissible
parameterized chain, not a fine-continuous map.

For a smooth coefficient function $a$ on $U$, define the \emph{Taylor-jet pullback}
\begin{equation}\label{contours:eq:jetpullback}
  T_\Phi a=\sum_{\alpha\in\N^d}
  \frac{(\partial^\alpha a)\circ\phi_0}{\alpha!}\,\eta^\alpha ,
  \qquad
  d\Phi_j=d\phi_{0,j}+d\eta_j ,
\end{equation}
and for $\omega=\sum_{\lambda,I}a_{\lambda,I}t^\lambda \dd x_{i_1}\wedge\cdots\wedge
\dd x_{i_q}$ set
\begin{equation}\label{contours:eq:formpullback}
  \Phi^*\omega=\sum_{\lambda,I}t^\lambda\,T_\Phi a_{\lambda,I}\,
  d\Phi_{i_1}\wedge\cdots\wedge d\Phi_{i_q} .
\end{equation}
The result is a form on the \emph{ordinary parameter manifold} with Hahn
coefficients. For analytic $a$ this agrees with the infinitesimal evaluation used
in the supplied manuscripts. For merely $C^\infty$ functions,
\eqref{contours:eq:jetpullback} is the \emph{specified jet prolongation}: we do not
assert ordinary convergence of its Taylor series, do not identify it with every
possible fine extension of $a$, and do not claim a canonical prolongation of every
smooth function to every fine domain. Real analyticity, when assumed below, is
imposed on the \emph{target} coefficient functions; the parameter coefficients need
only be smooth.

\begin{lemma}[Pullback and differential with a support certificate]
\label{contours:lem:pullback}
For $\omega\in\Omega_H^q(U)$ with support $S$, formula
\eqref{contours:eq:formpullback} defines an element of $\Omega_H^q(M)$ with
\[
  \supp(\Phi^*\omega)\subset S+E^* ,
\]
with only finitely many contributions at each exponent. It respects products, wedge
products, restrictions to boundary faces, ordinary smooth reparameterizations of
$M$, and the differential:
\begin{equation}\label{contours:eq:dcommute}
  d(\Phi^*\omega)=\Phi^*(d\omega).
\end{equation}
It also respects composition whenever the intervening ordinary coefficient maps are
analytic and the same positive-support substitution hypotheses hold. When the
coefficients are real analytic, it agrees with their canonical infinitesimal
evaluation at every ordinary parameter point.
\end{lemma}
\begin{proof}
Every term is indexed by one source exponent in $S$ together with a finite word of
error exponents in $E$; differentiating $\eta_e$ changes its coefficient function
but not its exponent, and exterior products introduce only finitely many coordinate
indices. \Cref{contours:lem:support} therefore gives a well-ordered total support
and only finitely many contributing words at each exponent, so each output
coefficient is a finite sum of products of ordinary smooth functions, forms, and
derivatives of ordinary functions, hence smooth. This proves existence and the
support bound.

For functions, ordinary Taylor product identities give $T_\Phi(ab)=(T_\Phi a)
(T_\Phi b)$. Differentiating \eqref{contours:eq:jetpullback} produces two groups of
terms --- differentiating $(\partial^\alpha a)\circ\phi_0$ contributes $d\phi_{0,j}$
and differentiating $\eta^\alpha$ contributes $d\eta_j$ --- and reindexing yields
\[
  d(T_\Phi a)=\sum_{j=1}^dT_\Phi(\partial_ja)\,(d\phi_{0,j}+d\eta_j),
\]
the reindexing being finite at every Hahn exponent by the first paragraph. This is
the chain rule needed for \eqref{contours:eq:dcommute}, and the exterior product
rule completes it. Equivalently, one may fix a target exponent, retain the finitely
many coefficient functions and Taylor degrees that can contribute, and verify the
identity as an ordinary finite Taylor-jet identity; no discarded higher degree can
affect that exponent. Boundary restriction and ordinary reparameterization commute
with ordinary derivatives and products, hence coefficientwise. The same finite-jet
argument, using the ordinary Taylor composition identity, proves the stated
composition compatibility. For analytic coefficients the defining Taylor jets
are exactly those of canonical infinitesimal evaluation. Equality at every
exponent proves the Hahn identities.
\end{proof}

\begin{definition}[Integral on a deformed chain]
\label{contours:def:deformedintegral}
For $\Phi$ as in \eqref{contours:eq:deformedmap} and an admissible form $\omega$ of
degree $\dim M$, put
\[
  \HInt_\Phi\omega:=\HInt_M\Phi^*\omega .
\]
Boundary parameterizations are obtained by restriction to the ordinary faces of
$M$. A finite chain is a finite linear combination of such parameterizations, whose
common faces are required to agree \emph{as parameterized Hahn data} with the
appropriate orientations.
\end{definition}

This is a definition on a parameterized object, not an attempt to integrate over its
image with a fine-topological measure.

\begin{theorem}[Stokes for Hahn-deformed chains]
\label{contours:thm:deformedstokes}
Let $M$ be compact oriented of dimension $n$, let $\Phi$ be an admissible
deformation, and let $\omega\in\Omega_H^{n-1}(U)$ on the ordinary coefficient
neighborhood $U$ --- so its coefficient forms are ordinary \emph{smooth} forms,
with no analyticity assumed. Then
\begin{equation}\label{contours:eq:deformedstokes}
 \boxed{\qquad
 \HInt_{\partial\Phi}\omega=\HInt_\Phi d\omega .
 \qquad}
\end{equation}
The same identity holds for finite sums of admissible chains with matching faces,
with the support bound of \cref{contours:lem:pullback}. The integral is invariant
under ordinary orientation-preserving smooth changes of parameter and under finite
admissible subdivision.
\end{theorem}
\begin{proof}
By \cref{contours:lem:pullback} the pulled-back form is a well-defined smooth Hahn
form on $M$ and commutes with $d$. Apply \cref{contours:thm:hahnstokes} on $M$:
\[
  \HInt_{\partial M}\Phi^*\omega
  =\HInt_Md(\Phi^*\omega)
  =\HInt_M\Phi^*(d\omega).
\]
Boundary compatibility identifies the first term with the integral on
$\partial\Phi$, since restriction to a face is the same as pulling back by the
restricted parameterization. Ordinary change of variables at every coefficient
proves reparameterization invariance, and subdivision follows by finite boundary
cancellation; internal faces cancel for finite chains.
\end{proof}

\begin{remark}[Why smoothness suffices]
Two of the three merged manuscripts state this theorem for real-analytic
coefficient forms. Smoothness is enough, and the reason is the one that runs
through the whole subject: the Taylor-jet pullback
\eqref{contours:eq:formpullback} is not summed by any convergence argument.
For each output exponent, \cref{contours:lem:support} gives only finitely many
words in the positive support $E$ together with the source support $S$.
Their lengths bound the Taylor degrees $|\alpha|$ that can contribute, so only
finitely many ordinary derivatives $\partial^\alpha a$ are ever read.
A bound by $|\alpha|\min E$ alone would not prove this in a higher-rank
ordered group, where the multiples of $\min E$ need not be cofinal.
Those exist for a smooth $a$. Requiring the ordinary Taylor series to converge
would be importing exactly the kind of limit this calculus does without.
\end{remark}

This is not merely a formal rewriting of $d^2=0$: it constructs a defined pairing
with a chain complex and proves that the pairing intertwines $d$ with the oriented
boundary. Two cautions. First, one cannot replace ``ordinary changes of parameter''
by arbitrary fine-analytic reparameterizations without proving a separate
substitution theorem and checking that the new pulled-back family remains
admissible. Second, meromorphic forms are allowed only when their denominators have
nonvanishing leading coefficient on the base image, where geometric inversion
supplies holomorphic Hahn data on a common ordinary neighborhood; \emph{passing
through a leading pole is not justified by calling its displacement
infinitesimal.}

\subsection{Affine scales and polynomial compatibility}
At an arbitrary scale write $z=a+rw$ with $a\in K$ and $r\in R_\Gamma$ positive, and
suppose the pulled-back form in the $w$ chart has admissible Hahn coefficients.
Define its integral by the chart pullback; for a one-form,
\begin{equation}\label{contours:eq:affineintegral}
  \HInt_{a+r\Phi}F(z)\dd z=\HInt_\Phi rF(a+rw)\dd w .
\end{equation}
For a real area form the Jacobian is $r^2$, and for an arbitrary constant invertible
real linear map the usual determinant appears; in several real coordinates one
inserts the ordinary algebraic Jacobian and wedge factors of the affine change.
Such factors lie in $K$ and do not disturb finite convolution. \emph{The scale
factor is not optional}: a $k$-dimensional uniform real scaling contributes $r^k$
to a $k$-form, and forgetting it would incorrectly erase microscopic areas and
volumes. This is the affine contour convention of \cite{Analysis}, here extended to
deformed chains and all form degrees. Cauchy and Stokes transfer through it with
their correct differential factors; a pole at $a+r\xi$ contributes the same residue
after substitution, since the factor $r$ in $\dd z$ cancels the scaling of a simple-pole
denominator.

A polynomial map with arbitrary Hahn coefficients need not have the specific
normalization $\phi_0+\eta$ of \eqref{contours:eq:deformedmap}. It still has a
finite polynomial pullback and can be integrated coefficientwise; expanding its
finite polynomial coefficients and applying ordinary simplex integration gives
exactly \eqref{contours:eq:simplexint}. Thus \cref{contours:thm:polstokes} and
\cref{contours:thm:hahnstokes} agree on their common domain, and polynomial
parameterizations are a useful enlargement beyond a single positive-support
normalization.

\subsection{The homotopy formula}
Let $\mathcal H:M\times[0,1]_\R\to U\sh$ be admissible in the sense of
\cref{contours:def:deformation}, with one common positive support $E$ on the
product, and write $\mathcal H_u$ for its restrictions. Define the degree-minus-one
operator
\begin{equation}\label{contours:eq:homotopyoperator}
  \mathcal K\omega=\HInt_0^1\iota_{\partial/\partial u}
  (\mathcal H^*\omega)\dd u ,
\end{equation}
the integration being coefficientwise in $u$ and leaving a form in the $M$
variables.

\begin{theorem}[Hahn homotopy formula]\label{contours:thm:homotopy}
For every admissible form $\omega$ with support $S$,
\begin{equation}\label{contours:eq:homotopyformula}
  \mathcal H_1^*\omega-\mathcal H_0^*\omega
  =d\,\mathcal K\omega+\mathcal K\,d\omega ,
\end{equation}
and $\mathcal K\omega$ has support in $S+E^*$. Consequently a closed form has equal
integrals on two admissibly homotopic cycles, and integrals of closed forms on
closed oriented parameter manifolds are invariant under such homotopies.
\end{theorem}
\begin{proof}
Write a coefficient of $\mathcal H^*\omega$ as $\beta(u)+\dd u\wedge\gamma(u)$, both
forms using only the $M$ variables. Exterior differentiation on the product gives
the ordinary coefficient identity
\[
  \beta(1)-\beta(0)=d_M\int_0^1\gamma(u)\dd u
  +\int_0^1\iota_{\partial_u}d(\mathcal H^*\omega)\dd u ,
\]
and \eqref{contours:eq:dcommute} converts the last term. The support bound of
\cref{contours:lem:pullback} is uniform over the ordinary compact parameter
cylinder, all coefficient integrations preserve support, and finite contributions at
each exponent make the calculation valid in the Hahn complex. If $d\omega=0$, the
difference is exact and its integral on a closed parameter manifold vanishes by
\cref{contours:thm:hahnstokes}; equivalently, apply
\cref{contours:thm:deformedstokes} to the cylinder with the usual opposite
orientations on its two ends.
\end{proof}

Ordinary deformations of the base map are allowed if their whole image lies in the
same ordinary coefficient domain; an infinitesimal deformation with fixed base is
automatically of this form. A homotopy passing through a pole is \emph{not}
admissible for a form singular there, and its period can change. For nonclosed forms
there is a genuine nonzero swept-region contribution: a moving boundary applied to
$\bar z\dd z$ changes the integral through
$d(\bar z\dd z)=2i\dd x\wedge \dd y$. \emph{It would be incorrect to claim that all
contour integrals are infinitesimally deformation invariant.} Closedness, not mere
existence of a derivative, is the controlling condition; \cref{contours:ex:ellipse}
exhibits the effect exactly.

\subsection{An order-valued parameter estimate}
The coefficientwise integral does not discard the surreal order. If
$A(s)=\sum_\lambda a_\lambda(s)t^\lambda$ has continuous real coefficients on an
ordinary compact interval, one common support, and $A(s)\ge0$ at every ordinary $s$,
then either $A\equiv0$ or its first nonzero coefficient function is nonnegative
everywhere and positive somewhere, so its ordinary integral is positive; hence
$\HInt A(s)\dd s\ge0$.

\begin{proposition}[An order-valued parameter estimate]\label{contours:prop:ml}
Write an admissible pullback as $B(s)\dd s$ on an ordinary interval. Suppose $m(s)$
is an ordinary nonnegative continuous function and $M\in R_\Gamma$ is nonnegative
with $|B(s)|\le M\,m(s)$ for all ordinary $s$. Then
\begin{equation}\label{contours:eq:ml}
  \Bigl|\HInt B(s)\dd s\Bigr|\ \le\ M\int m(s)\dd s .
\end{equation}
For an ordinary piecewise $C^1$ contour $\gamma$ of ordinary length $L$ and a
uniformly supported continuous Hahn family $F$ along it with $|F(\gamma(s))|\le M$,
this specializes to $\bigl|\HInt_\gamma F\dd z\bigr|\le LM$.
\end{proposition}
\begin{proof}
If $I=\HInt B\dd s\ne0$, put $u=\overline I/|I|$, so $|u|=1$ and $\re(uI)=|I|$. The
real Hahn function $M\,m(s)-\re(uB(s))$ is pointwise nonnegative and uniformly
supported; positivity and $K$-linearity give \eqref{contours:eq:ml}. The case $I=0$
is immediate, and piecewise regular data are treated on finitely many intervals. For
the contour form take $m(s)=|\gamma'(s)|$, whose integral is $L$; the integral of
$M|\gamma'(s)|-\re(uF(\gamma(s))\gamma'(s))$ is then $LM-|I|$.
\end{proof}

This is the estimate proved in the analysis manuscript, here given as a numbered
order-valued statement, and it shows why a coefficientwise construction retains an
exact surreal --- rather than merely standard-part --- inequality. The formulation
deliberately avoids assuming that the speed $|\Phi'(s)|$ possesses a globally
uniform smooth Hahn expansion through every possible zero of its leading speed.
Ordered estimates do not imply convergence of set-indexed nets in the full fine
topology.

\section{Cauchy theory as a period theory}
\label{contours:sec:cauchy}

Throughout this section the coefficient ring is the \emph{common-domain coherent
ring} $\HH(U)=\OO(U)((t^\Gamma))$ of \cref{contours:subsec:rings}, for an ordinary
open $U\subset\C$. Every theorem below is stated over that ring and over no other;
the radius-free ring $\A_n$ is treated separately in
\cref{contours:sec:radiusfree,contours:sec:multi}, and none of the statements here
is transferred to it.

\subsection{Coherent holomorphy and the vanishing theorem}
An element $F=\sum_{\lambda\in S}f_\lambda t^\lambda\in\HH(U)$ has one well-ordered
support $S$ and ordinary holomorphic coefficients $f_\lambda\in\OO(U)$ on
\emph{one} common $U$. At $z=c+\varepsilon\in U\sh$, evaluation means
\begin{equation}\label{contours:eq:coherentevaluation}
  F\sh(c+\varepsilon)
  =\sum_{\lambda,\,n\ge0}t^\lambda
   \frac{f_\lambda^{(n)}(c)}{n!}\,\varepsilon^n ,
\end{equation}
whose support is contained in $S+(\supp\varepsilon)^*$, so it is strongly defined.
Differentiation acts on the $f_\lambda$. The injectivity, composition, and gluing
properties of $\sh$ are supplied in \cite[\S5]{Analysis}. For an ordinary piecewise
$C^1$ path $\gamma$ in $U$,
\begin{equation}\label{contours:eq:Hintegral}
  \HInt_\gamma F(z)\dd z=\sum_\lambda t^\lambda\int_\gamma f_\lambda(z)\dd z ,
\end{equation}
and the support cannot increase. In real coordinates
\begin{equation}\label{contours:eq:dFdz}
  d(F\dd z)=2i\,\partial_{\bar z}F\dd x\wedge \dd y=0 ,
\end{equation}
the last equality being coefficientwise ordinary holomorphy: the Cauchy--Riemann
equation holds for each coefficient, not for an arbitrary fine-local function.

\begin{theorem}[Coherent Cauchy theorem]\label{contours:thm:cauchy}
Let $F\in\HH(U)$. If $\gamma_0$ is an ordinary piecewise smooth closed cycle
null-homologous in $U$, then $\HCoint_{\gamma_0}F\dd z=0$. If an admissible closed
contour bounds a finite admissible $2$-chain on which the coefficient functions of
$F$ are holomorphic, its period is likewise zero. In particular, for a
Hahn-deformed closed contour $\gamma=\gamma_0+\eta$ in the sense of
\cref{contours:def:deformation}, the period equals the period of $\gamma_0$,
whether or not that period is zero; and every closed contour has zero period when
$U$ is simply connected.
\end{theorem}
\begin{proof}
For an ordinary contour apply ordinary Cauchy to every coefficient, or use a smooth
filling chain with \eqref{contours:eq:dFdz} in \cref{contours:thm:hahnstokes}. The
bounding-chain assertion is \cref{contours:thm:deformedstokes} applied to the
closed form $F\dd z$. For a deformed contour the homotopy $\gamma_u=\gamma_0+u\eta$,
$0\le u\le1$, is admissible with the same positive error support, so
\cref{contours:thm:homotopy} applies to the closed form $F\dd z$; alternatively
\cref{contours:thm:endpoints} identifies the two periods directly. The simply
connected case follows from ordinary complex homology or from coefficientwise
primitives.
\end{proof}

Null-homology is a genuine hypothesis. On $U=\C^\times$ the coherent form
$\dd z/z$ has unit-circle period $2\pi i$, so ``closed contour'' alone does not
imply a vanishing integral. Conversely, self-intersections do not obstruct
Cauchy's theorem when the cycle is null-homologous: neither the contour nor a
filling chain need be embedded, and the Jordan theorem is not needed for this
vanishing statement. The Jordan theorem is useful when identifying a particular
simple boundary with its interior, or when turning winding into an inside/outside
count.

\subsection{Exact transport of an open contour}
For $c\in U$ and $\varepsilon\in\m_K$ define the infinitesimal primitive increment
\begin{equation}\label{contours:eq:endpointT}
  T_F(c,\varepsilon)=\sum_{n\ge0}\frac{F^{(n)}(c)}{(n+1)!}\,\varepsilon^{n+1},
\end{equation}
which involves no choice of global primitive.

\begin{theorem}[Exact endpoint-transport formula]\label{contours:thm:endpoints}
Let $\gamma_0:[a,b]_\R\to U$ be ordinary piecewise smooth, let $F\in\HH(U)$, and let
$\Gamma=\gamma_0+\eta$ be an admissible complex deformation, piecewise on the same
finite partition. Then
\begin{equation}\label{contours:eq:endpoints}
 \boxed{\quad
 \HInt_\Gamma F\dd z-\HInt_{\gamma_0}F\dd z
 =T_F\bigl(\gamma_0(b),\eta(b)\bigr)-T_F\bigl(\gamma_0(a),\eta(a)\bigr).
 \quad}
\end{equation}
\emph{No simple-connectedness assumption on $U$ is required.} In particular a closed
deformed contour with matching endpoint displacement has exactly the period of its
ordinary base contour, not merely the same standard part.
\end{theorem}
\begin{proof}
Differentiate $T(s)=T_F(\gamma_0(s),\eta(s))$ coefficientwise. The derivative of the
$F^{(n)}$ factor gives
\[
  \gamma_0'(s)\sum_{n\ge0}\frac{F^{(n+1)}(\gamma_0(s))}{(n+1)!}\eta(s)^{n+1}
  =\gamma_0'(s)\bigl(F\sh(\Gamma(s))-F(\gamma_0(s))\bigr),
\]
and the derivative of the displacement gives $\eta'(s)F\sh(\Gamma(s))$, so
\[
  T'(s)=F\sh(\Gamma(s))\Gamma'(s)-F(\gamma_0(s))\gamma_0'(s).
\]
All series have support in $\supp F+E^*$, and \cref{contours:lem:support} justifies
every reindexing. Integrating the ordinary parameter coefficientwise and applying
the ordinary fundamental theorem to each coefficient gives
\eqref{contours:eq:endpoints}. Splitting into finitely many smooth pieces introduces
endpoint terms which cancel at the joins; for a closed contour the two remaining
terms coincide.

Equivalently, set $\gamma_u(s)=\gamma_0(s)+u\eta(s)$ and use the identity
\begin{equation}\label{contours:eq:variation}
  \partial_u\bigl(F(\gamma_u)\partial_s\gamma_u\bigr)
  =\partial_s\bigl(F(\gamma_u)\partial_u\gamma_u\bigr),
\end{equation}
all products and derivatives being admissible by \cref{contours:lem:pullback};
integrating coefficientwise over the ordinary rectangle $[0,1]\times[a,b]$ and
Taylor expanding identifies the two endpoint terms as
\eqref{contours:eq:endpointT}. These derivatives differentiate the ordinary
coefficient functions of the parameter, not a putative fine-continuous path.
\end{proof}

Consequently the analysis manuscript's endpoint integral is the value of
\emph{every} admissible infinitesimal deformation of its base contour with those
endpoints; its endpoint correction is not an arbitrary prescription. A primitive is
not required, because only the homotopy strip and not an arbitrary filling disk is
used. In contrast, choosing a different ordinary homology class can still change
the answer by a period: a coherent global primitive eliminates those periods,
merely fine-local primitives do not. The corresponding fundamental theorem for
finite nonstandard endpoints $z_j=a_j+\varepsilon_j$ reads
\begin{equation}\label{contours:eq:endpointFTC}
  -T_F(a_1,\varepsilon_1)+\HInt_\gamma F\dd z+T_F(a_2,\varepsilon_2)
  =\text{difference of a coherent primitive at the endpoints},
\end{equation}
which does not contradict \cref{contours:prop:noftc} because the class of
admissible primitives has been restricted.

\begin{example}[The logarithm obstruction has not disappeared]
\label{contours:ex:logarithm}
The global fine-locally analytic logarithm $\mathcal L$ on $\SC^\times$ with
$\mathcal L'=1/z$ uses a selected argument cut and is \emph{not} a coherent
primitive on $(\C^\times)\sh$. Accordingly $\HCoint_{|z|=1}\dd z/z=2\pi i$ does not
contradict its existence. Applying \cref{contours:thm:endpoints} to this
noncoherent global function would be an invalid change of category.
\end{example}

\subsection{Winding at a displaced center, and the Cauchy formula}
Let $a=a_0+\xi$ with $a_0\in\C$, $\xi\in\m_K$, and $a_0$ not on the ordinary
contour $\gamma_0$. On a common ordinary collar of the contour,
\[
  \frac1{z-a}=\sum_{n\ge0}\frac{\xi^n}{(z-a_0)^{n+1}}
\]
is an admissible Hahn family. Higher-pole terms are ordinary derivatives and have
zero periods.

\begin{theorem}[Winding stability at one scale]\label{contours:thm:winding}
Let $\gamma_0$ be an ordinary closed piecewise smooth contour, let $c\notin\gamma_0$,
let $\eta$ be an admissible closed positive-support deformation, and let
$\xi\in\m_K$, $a\in K$, $r\in R_\Gamma$ with $r>0$. Set
$\Gamma=a+r(\gamma_0+\eta)$ and $b=a+r(c+\xi)$. Then the integrand is coherent on an
ordinary neighborhood of the base contour and
\begin{equation}\label{contours:eq:winding}
  \frac1{2\pi i}\HCoint_\Gamma\frac{\dd z}{z-b}=\wind(\gamma_0,c)\in\Z .
\end{equation}
\end{theorem}
\begin{proof}
The affine factor $r$ cancels by \eqref{contours:eq:affineintegral}, reducing to
$a=0$, $r=1$. Two proofs are available. By \cref{contours:thm:endpoints} the closed
period is unchanged under deformation from $\gamma_0$ to $\gamma_0+\eta$, and on the
ordinary base the $n=0$ term of the expansion has integral $2\pi i\wind(\gamma_0,c)$
while each $n>0$ term is an ordinary derivative and has zero integral; strong
summability justifies coefficientwise integration. Alternatively, on the parameter
cylinder set $g_u(s)=\gamma_0(s)-c+u(\eta(s)-\xi)$. Its leading coefficient never
vanishes, compactness of the ordinary contour supplying a positive ordinary
separation from $c$, so $1/g_u$ is a uniformly supported Hahn family; direct
differentiation gives
$\partial_u(\partial_sg_u/g_u)=\partial_s(\partial_ug_u/g_u)$, whose integral along
the closed parameter is zero coefficient by coefficient, so the period at $u=1$
equals the ordinary winding integral at $u=0$.
\end{proof}

Thus an admissible contour has an exact \emph{integer} period index, not an
unspecified surreal winding number. The theorem treats only points whose standard
part avoids the base curve; a point reducing to the curve lies in its boundary tube
and requires a finer chart, an explicit detour, or the separate algebraic definition
of \cref{contours:sec:algebraic}. This is precisely the hypothesis that fails for
$a=1\pm t$ in \cref{contours:subsec:boundarylayer}.

\begin{theorem}[Cauchy formula at displaced points and on deformed contours]
\label{contours:thm:cauchyformula}
Let $D$ be a bounded ordinary Jordan domain with piecewise smooth positively
oriented boundary and $\overline D\subset U$, let $F\in\HH(U)$, and let $\Gamma$ be
either $\partial D$ or an admissible closed infinitesimal deformation of it. For
$z=c+\varepsilon$ with $c\in D$, $\varepsilon\in\m_K$, and every $n\ge0$,
\begin{equation}\label{contours:eq:cauchyformula}
  (F^{(n)})\sh(z)=\frac{n!}{2\pi i}\HCoint_\Gamma
  \frac{F(\zeta)}{(\zeta-z)^{n+1}}\dd\zeta .
\end{equation}
The corresponding statement holds after any admissible affine rescaling.
\end{theorem}
\begin{proof}
Expand
\[
  (\zeta-c-\varepsilon)^{-n-1}
  =\sum_{k\ge0}\binom{n+k}{k}\varepsilon^k(\zeta-c)^{-n-k-1}.
\]
The ordinary point $c$ has positive ordinary distance from the compact boundary, so
all coefficient functions are holomorphic on one common collar, and
\cref{contours:lem:support} permits multiplication by $F$ and coefficientwise
integration. Ordinary Cauchy gives the coefficient $f_\lambda^{(n+k)}(c)/k!$,
yielding $\sum_{\lambda,k}t^\lambda f_\lambda^{(n+k)}(c)\varepsilon^k/k!$, which is
\eqref{contours:eq:coherentevaluation} for $F^{(n)}$. The support is contained in
$\supp F+(\supp\varepsilon)^*$. The kernel form is closed on the collar, so the
deformed case follows from \cref{contours:thm:endpoints}; affine substitution gives
the rescaled statement with the ordinary derivative and differential factors.
\end{proof}

This proves a formula for genuinely displaced points and contours. It does
\emph{not} claim that the deformed trace is the fine-topological boundary of a
fine-connected interior, and by \cref{contours:ex:boundarylayer} its hypothesis
$c\in D$ is not the same as $|z|<1$ for the unit disk.

On a simply connected ordinary domain, normalized ordinary primitives of all
$f_\lambda$ assemble into a coherent primitive with support in $S$; on a general
connected domain such a primitive exists exactly when all ordinary closed-cycle
Hahn periods vanish, and all coherent primitives of one function differ by a single
constant in $K$ \cite[\S6]{Analysis}. Morera's criterion holds for common-support
continuous coefficient families, since vanishing Hahn integrals around every
ordinary triangle means vanishing of every ordinary coefficient integral.

\subsection{Moving poles and actual zero clusters}
Let $F,G\in\HH(U)$ with $G\ne0$, and suppose the ordinary leading coefficient
function $g_0$ of $G$ after monomial normalization has no zero on the ordinary
boundary of a bounded Jordan domain $D$ with $\overline D\subset U$ and finitely
many zeros inside. The preparation and moving-pole theorem of \cite{Analysis} gives
\begin{equation}\label{contours:eq:movingres}
  \frac1{2\pi i}\HCoint_{\partial D}\frac FG\dd z
  =\sum_{b\in D\sh}\Res_b\Bigl(\frac FG\dd z\Bigr),
\end{equation}
the sum being over \emph{actual surcomplex poles}, not merely their standard parts.
The sum is finite; its summands may be infinite. As fixed in
\cref{contours:subsec:residuedef}, the residue at $b$ is the coefficient of
$(z-b)^{-1}$ in the local finite-meromorphic germ \emph{after cancellations}, and is
not merely the residue of one ordinary coefficient at $\st(b)$.

The decisive local reduction is worth recording, because it is the mechanism reused
in \cref{contours:thm:localresidue}. Preparation near a leading zero writes $G$ as a
coherent unit times a monic polynomial $P$ whose lower coefficients are
infinitesimal; division writes $F/G=J+R/P$ with $J$ coherent holomorphic and
$\deg R<\deg P$; the polynomial has finitely many infinitesimal roots, and rational
partial fractions give $R/P=\sum_b\sum_{j=1}^{m_b}A_{b,j}(z-b)^{-j}$, where $m_b$ is
the order of $b$ as a root of $P$. On a small ordinary protecting circle only the
$j=1$ terms have nonzero coefficientwise period, contributing $2\pi iA_{b,1}$, and
the holomorphic part contributes zero. This argument is valid in higher rank because
the expansions are strong sums, not sequential limits. It is stronger than calling a
coefficientwise residue at the reduced center ``the residue'' and leaving its
relation to the split points unspecified.

\begin{corollary}[Residues on a deformed protecting contour]
\label{contours:cor:deformedres}
Formula \eqref{contours:eq:movingres} remains true when $\partial D$ is replaced by
an admissible closed infinitesimal deformation lying in a coefficient collar on
which $g_0$ is nonzero. The right side is still the divisor selected by the ordinary
interior $D\sh$.
\end{corollary}
\begin{proof}
The fraction is Hahn-holomorphic on that collar, so its closed period is unchanged
by \cref{contours:thm:endpoints}.
\end{proof}

This describes the selected divisor without inventing a fine-connected interior of
the deformed trace. The same theorem applied to $F'/F$ gives the argument
principle: its normalized period is the total zero multiplicity, or zeros minus
poles for a fraction, with ordinary integer multiplicities. Positive-support
perturbations preserving the leading function preserve the total multiplicity in
each ordinary monad \cite{Analysis,Geometry}; individual roots can nevertheless
split at finer scales, which is the subject of
\cref{contours:sec:scales,contours:sec:radiusfree}.

\begin{example}[A deformed ellipse and its exact area]\label{contours:ex:ellipse}
Let $\varepsilon,\delta\in\m_K\cap R_\Gamma$ and map the ordinary unit disk by
$\Phi(x,y)=(1+\varepsilon)x+i(1+\delta)y$, with boundary
$\Gamma(\theta)=(1+\varepsilon)\cos\theta+i(1+\delta)\sin\theta$. Since the
pullback of $\dd x\wedge \dd y$ is $(1+\varepsilon)(1+\delta)\dd x\wedge \dd y$,
\cref{contours:thm:deformedstokes} applied to $\bar z\dd z$ gives
\begin{equation}\label{contours:eq:ellipse}
  \HCoint_\Gamma\bar z\dd z
  =2\pi i(1+\varepsilon)(1+\delta)
  =2\pi i(1+\varepsilon+\delta+\varepsilon\delta).
\end{equation}
Thus an infinitesimal deformation changes a nonclosed form's integral by a genuine
infinitesimal, \emph{including a second-order mixed term}. Setting
$\delta=\varepsilon$ gives $2\pi i(1+\varepsilon)^2$; further affine scaling by $r$
multiplies by $r^2$; translation adds an exact one-form with zero closed period. In
contrast $z^q\dd z$ is exact for every integer $q\ge0$ and integrates to zero over
both the circle and the ellipse. Writing $R=1+t$ and $S=1+t^\omega$ gives
$2\pi i(1+t+t^\omega+t^{1+\omega})$, a value visible only if all valuation ranks are
retained. For the undeformed circle of radius $r\in R_{\Gamma,>0}$ the same identity
reads $\HCoint_{|z|=r}\bar z\dd z=2\pi ir^2=2i\HInt_{r\DD}\dd x\dd y$, giving a
nonzero infinitesimal area for $r=t^\rho$ and an infinite one for infinite $r$. No
measure on all fine-topological subsets is asserted.
\end{example}

\begin{example}[Separating residues invisible at the coarse scale]
\label{contours:ex:cluster}
Let $t=\omega^{-1}$ and $R(z)=1/(z^2-t^2)$. Its residues at $\pm t$ are
$\pm1/(2t)$, so
\[
  \frac1{2\pi i}\HCoint_{|z|=1}R(z)\dd z=0 :
\]
both poles reduce to zero and the coarse contour sees only their cancelling sum.
Choose $r=t^\rho$ with $\rho>1$ --- \emph{including a higher Archimedean class such
as $\rho=\omega$} --- and use the contour $z=t+rw$, $|w|=1$ in the ordinary
parameter. Then
\begin{equation}\label{contours:eq:clusterpullback}
  R(z)\dd z=\frac{\dd w}{w(2t+rw)}
  =\frac1{2t}\sum_{k\ge0}\Bigl(-\frac{r}{2t}\Bigr)^kw^{k-1}\dd w ,
\end{equation}
whose relative support is positive because $v(r/t)=\rho-1>0$. Its normalized period
is $1/(2t)$, so the actual integral is $\pi i/t$; the analogous circle about $-t$
gives $-\pi i/t$. The same support-controlled calculus therefore resolves the two
infinite residues individually and recovers their cancellation on the larger
contour. For $z/(z^2-t^2)$ the two residues are $1/2$ and the coarse normalized
period is $1$. No all-scale coherence hypothesis on any transcendental function was
used: a contour only needs a valid representation at its chosen scale, so the
all-scale rigidity theorems of \cite{Analysis} pose no obstruction to this local
multiscale calculus.
\end{example}

\section{Affine scales, microscopic contours, and compatible charts}
\label{contours:sec:scales}

\subsection{A precise compatibility criterion}
\begin{proposition}[Agreement on admissible overlaps]\label{contours:prop:overlap}
Suppose two chart descriptions of the same parameterized oriented chain produce
\emph{pointwise equal} pulled-back differential forms at every ordinary parameter
point, and both pulled-back forms are common-support Hahn families with smooth
coefficients. Then their integrals agree. The same holds for a finite collection of
chart pieces with matching oriented faces.
\end{proposition}
\begin{proof}
Enlarge the workspace to contain both supports. Uniqueness of Hahn normal forms
makes each coefficient of the difference vanish at every parameter point, so the
coefficient forms are equal; their ordinary integrals agree, and so do their Hahn
sums. For finitely many pieces use additivity and cancellation of matching faces.
\end{proof}

This criterion is intentionally stronger than equality of the set-theoretic images
of two curves: a chain includes multiplicity, orientation, and the pulled-back
differential. It is also weaker than requiring one global chart for every scale, as
finitely many verified overlap identities suffice. What is \emph{not} justified is a
patchwork of arbitrary local fine germs for which common-support pullbacks and
their identities have not been established; equality of bare fine germs in
unrelated, unverified summation descriptions is not a substitute.

\subsection{Macroscopic and microscopic winding resolve different information}
A macroscopic contour whose reduction avoids a cluster assigns the same winding to
every point of that cluster. A contour centered on one actual pole and rescaled
below the finite separation distances can resolve an individual contribution. This
requires no surreal-valued winding number: ordinary integer winding numbers in
different admissible charts suffice.

\begin{example}[One coarse contour, two microscopic contours]
\label{contours:ex:expcluster}
The poles of $e^z/(z^2-t)$ are $\pm t^{1/2}$, with residues
$e^{t^{1/2}}/(2t^{1/2})$ and $-e^{-t^{1/2}}/(2t^{1/2})$. These are individually
infinite, but their sum is
\begin{equation}\label{contours:eq:expresidue}
  \frac{\sinh(t^{1/2})}{t^{1/2}}=\sum_{n\ge0}\frac{t^n}{(2n+1)!} .
\end{equation}
The ordinary unit circle and the protecting microscopic circle $|z|=t^{1/4}$ both
have integral $2\pi i$ times \eqref{contours:eq:expresidue}. A circle centered at
$t^{1/2}$ of radius $t$ isolates the positive pole: in the coordinate
$z=t^{1/2}+tw$ the second pole sits at $w=-2t^{-1/2}$, an infinite normalized
distance. Each integral is computed in its own admissible chart, and the outer
integral is the sum of the two inner residues even though the individual residues
are infinite. \emph{There is no assertion that these ordinary-parameter traces
constitute ordinary topological circles in the fine plane.}
\end{example}

\begin{example}[Two poles in different valuation ranks]
\label{contours:ex:rankcircles}
Let $\Gamma=\Q+\Q\omega$ with its inherited surreal order and put
$R(z)=1/(z-t)+1/(z-t^\omega)$, both residues being $1$. Three positively oriented
microscopic circles give
\begin{equation}\label{contours:eq:rankexample}
  \frac1{2\pi i}\HCoint_{|z|=t^\rho}R(z)\dd z=
  \begin{cases}
    2,&\rho=\tfrac12,\\
    1,&\rho=2,\\
    0,&\rho=2\omega .
  \end{cases}
\end{equation}
For the middle circle, $t^\omega$ lies inside but $t$ lies outside, because
$\omega>2>1$ reverses the magnitude order of the monomials; the last circle is
smaller than both pole scales. Each assertion follows by geometric expansion in the
appropriate ratio, or by algebraic winding and \cref{contours:thm:algpairing}.
\emph{A single real-valued valuation that discards the higher rank cannot retain all
of these distinctions.}
\end{example}

\section{Radius-free germs become coherent after microscopic rescaling}
\label{contours:sec:radiusfree}

From here to the end of \cref{contours:sec:multi} the coefficient ring is the
\emph{radius-free germ ring}
\[
  \A_n=\C\{z_1,\ldots,z_n\}((t^\Gamma))
\]
of \cref{contours:subsec:rings}, and the target of rescaling is the
polynomial-coefficient ring $\C[w_1,\ldots,w_n]((t^\Gamma))$. No theorem of
\cref{contours:sec:cauchy}, which is stated over $\HH(U)$, is applied to $\A_n$
before the rescaling theorem has moved the data into a common-domain ring.

\subsection{The missing common-radius problem}
An element $H=\sum_\lambda h_\lambda(z)t^\lambda\in\A_n$ has ordinary convergent
germs as coefficients, with no requirement that all $h_\lambda$ extend to the same
ordinary polydisk. The witness \eqref{contours:eq:shrinking} is legitimate, but its
coefficient poles at $1/j$ prevent a common ordinary disk of positive radius, so an
ordinary fixed-circle coefficient integral at zero is unavailable; applying
\cref{contours:thm:cauchyformula} to it on a fixed ordinary circle would leave a
gap. Evaluation at infinitesimal points is nevertheless strongly defined. The remedy
is rescaling.

\begin{theorem}[Positive rescaling gives polynomial coefficient families]
\label{contours:thm:rescaling}
Let $\rho_1,\ldots,\rho_n\in\Gamma_{>0}$, put $r_j=t^{\rho_j}$, and write
$rw=(r_1w_1,\ldots,r_nw_n)$. The Taylor substitution
\begin{equation}\label{contours:eq:rescaling}
  \sigma_r(H)(w):=H(rw)
  =\sum_{\lambda,\alpha}h_{\lambda,\alpha}\,
   t^{\lambda+\alpha\cdot\rho}\,w^\alpha
\end{equation}
defines an \emph{injective $K$-algebra homomorphism}
\[
  \sigma_r:\A_n\longrightarrow\C[w_1,\ldots,w_n]((t^\Gamma)) .
\]
Every fixed Hahn coefficient of \eqref{contours:eq:rescaling} is an ordinary
polynomial; the support is contained in $\supp H+\{\rho_1,\ldots,\rho_n\}^*$; the
image has coherent evaluation on the thickening of every ordinary $w$-domain; and
$\partial_{w_j}\sigma_r(H)=r_j\,\sigma_r(\partial_{z_j}H)$. More generally, for
$a\in\m_K^n$ the translated substitution
\begin{equation}\label{contours:eq:translatedrescale}
  H(a_1+r_1w_1,\ldots,a_n+r_nw_n)\in\C[w_1,\ldots,w_n]((t^\Gamma))
\end{equation}
has support contained in $\supp H+(E_a\cup\{\rho_1,\ldots,\rho_n\})^*$, where
$E_a=\bigcup_i\supp(a_i)$, and it evaluates correctly at every finite tuple
$w\in(K^\circ)^n$.
\end{theorem}
\begin{proof}
Write the Taylor germ $h_\lambda(z)=\sum_\alpha h_{\lambda,\alpha}z^\alpha$. At a
fixed Hahn exponent $\nu$, the possible contributions satisfy
$\lambda+\alpha_1\rho_1+\cdots+\alpha_n\rho_n=\nu$. \Cref{contours:lem:support}
gives finitely many source exponents $\lambda$ and finitely many finite words in the
positive alphabet $\{\rho_1,\ldots,\rho_n\}$, hence finitely many multi-indices
$\alpha$ contributing at $\nu$; the coefficient at $\nu$ is therefore an ordinary
polynomial, and the total support is well ordered. Taylor multiplication and
differentiation are valid at each exponent because the relevant regroupings are
finite, which gives the algebra and derivative statements. The Taylor embedding
$\A_n\hookrightarrow K[[z]]$ is injective, and diagonal substitution $z_j=r_jw_j$ is
injective in that formal ring because $r_j\ne0$; hence $\sigma_r$ is injective.
Polynomial coefficients are holomorphic on every ordinary domain, and evaluation at a
finite normalized point agrees with evaluation of $H$ at the infinitesimal point
$rw$ by the same finite-contribution Taylor identities. For
\eqref{contours:eq:translatedrescale}, first translate by $a$ using the strongly
summable Taylor translation of \cite[\S3]{Geometry}, whose support lies in
$\supp H+E_a^*$, then repeat the argument; evaluation at finite $w=c+\varepsilon$ is
justified by adjoining the positive supports of the infinitesimal parts.
\end{proof}

\emph{No coefficient radii enter this proof.} The argument in fact works for
arbitrary \emph{formal} complex coefficient series at infinitesimal arguments,
although the ring of interest here is the analytic-germ ring. Its analytically
relevant conclusion is not ordinary convergence of the original germs on a uniform
disk, but a \emph{new coherent representation in an infinitesimal chart}; it
neither establishes nor requires all-scale coherence of the original function. For
\eqref{contours:eq:shrinking} the substitution $z=t^{1/4}w$ gives
$H(t^{1/4}w)=\sum_{j\ge1}\sum_{q\ge0}j^qt^{j+q/4}w^q$, at a fixed exponent a finite
sum, so the singularities of the original ordinary coefficient representatives no
longer obstruct coefficientwise contour integration in the $w$-plane. The resulting
coefficients need not have any uniform degree or growth bound as the Hahn exponent
varies.

\emph{Convention.} For $r=t^\rho$ the symbol $\{|z|=r\}$ in a Hahn integral always
means the ordinary parameterized circle $z=re^{i\theta}$, $0\le\theta\le2\pi$, with
its pulled-back differential. \emph{It does not mean the full class of points of
that modulus equipped with a fine-topological integration theory}, and it is not the
full semialgebraic circle of \cref{contours:ex:definablecircle}.

\subsection{\texorpdfstring{Microscopic Cauchy formulae over $\A_n$}%
{Microscopic Cauchy formulae over the radius-free ring}}
\begin{theorem}[Cauchy formula for a radius-free germ]
\label{contours:thm:localcauchy}
Let $H\in\A_1$ and $a\in\m_K$. Choose $\rho>0$ with $\rho<v(a)$ when $a\ne0$, and
put $r=t^\rho$; if $a=0$, any $\rho\in\Gamma_{>0}$ is allowed. Then for every
$k\ge0$
\begin{equation}\label{contours:eq:localCauchy}
 \boxed{\quad
 H^{(k)}(a)=\frac{k!}{2\pi i}\HCoint_{|z|=r}
 \frac{H(z)}{(z-a)^{k+1}}\dd z .
 \quad}
\end{equation}
\end{theorem}
\begin{proof}
By \cref{contours:thm:rescaling}, $G(w)=H(rw)$ lies in $\C[w]((t^\Gamma))\subset
\HH(\C)$ with ordinary entire polynomial coefficients, and the normalized point
$a/r$ is infinitesimal. Apply \cref{contours:thm:cauchyformula} to $G$ on the
ordinary unit circle. Since $G^{(k)}(a/r)=r^kH^{(k)}(a)$ and $\dd z=r\dd w$, the
scaling factors yield \eqref{contours:eq:localCauchy}.
\end{proof}

The same proof gives a product-torus Cauchy formula for $H\in\A_n$, using one
sufficiently large infinitesimal circle around each coordinate of the evaluation
point. \emph{A product of $n$ circles is a real $n$-dimensional cycle, the
distinguished boundary of a polydisk; it is not the real $(2n-1)$-dimensional
boundary of that polydisk.} There is also a local primitive theorem: ordinary
germwise integration applied coefficientwise makes $D:\A_1\to\A_1$ surjective with
kernel $K$, and the ordinary analytic Poincar\'e homotopy applied coefficientwise
makes the analytic differential-form complex of $\A_n$ exact in positive degrees.
Thus microscopic periods of nonsingular forms vanish while poles supply residue
obstructions. \emph{These assertions use the prescribed Hahn-germ algebra, not all
fine-analytic functions.}

\subsection{Protecting radii and one-variable residues}
\begin{lemma}[A radius protecting an entire one-variable cluster]
\label{contours:lem:protect}
Let $F=f+E\in\A_1$, where $f$ has ordinary order $m\ge1$ at zero, say
$f(z)=z^mu(z)$ with $u(0)\ne0$, and $\delta=\vH(E)>0$ (with $\delta=+\infty$ if
$E=0$). Choose $\rho>0$ with $m\rho<\delta$ and put $r=t^\rho$; if $E=0$ every
positive $\rho$ works. Then $r^{-m}F(rw)$ has leading coefficient $u(0)w^m$ and no
zeros on the ordinary unit-circle parameter, and every infinitesimal zero $a$ of $F$
satisfies $v(a)>\rho$.
\end{lemma}
\begin{proof}
The higher powers in $u(rw)-u(0)$ have positive valuation, and every term of
$r^{-m}E(rw)$ has valuation at least $\delta-m\rho>0$; this proves the
leading-coefficient and boundary assertions. For a nonzero infinitesimal zero with
$v(a)\le\rho$, the term $f(a)$ has valuation $m\,v(a)\le m\rho<\delta$ while $E(a)$
has valuation at least $\delta$, so they cannot cancel. The zero $a=0$ satisfies the
conclusion by convention. Divisibility of $\Gamma$ supplies a strictly smaller
positive $\rho$ in the nonzero case, even for arbitrary rank.
\end{proof}

The interpretation is a \emph{protecting-radius condition}: the chosen circle must
be infinitesimally large relative to the perturbation scales after accounting for
the order $m$. Every zero in the originating infinitesimal cluster lies strictly
inside the protecting circle.

\begin{theorem}[Microscopic circle realization of the one-variable perturbation
residue]\label{contours:thm:microcircle}
Let $F=f+E\in\A_1$ and $H\in\A_1$ be as in \cref{contours:lem:protect}, with
$r=t^\rho$ and $m\rho<\vH(E)$. Then the microscopic circle functional is well
defined and
\begin{equation}\label{contours:eq:microcircle}
 \boxed{\qquad
 \frac1{2\pi i}\HCoint_{|z|=r}\frac{H(z)}{F(z)}\dd z=\Res_F(H),
 \qquad}
\end{equation}
with $\Res_F$ the one-variable perturbation residue
\eqref{contours:eq:residueone}. The integral means substitution $z=rw$ followed by
coefficientwise integration over the ordinary unit circle in $w$, and it is
independent of every choice of $\rho$ satisfying the displayed bound. Its support is
contained in $\supp H+(\supp E)^*$.
\end{theorem}
\begin{proof}
Write $f(z)=cz^m+O(z^{m+1})$, $c\ne0$. By \cref{contours:thm:rescaling} both
$H(rw)$ and $F(rw)$ have polynomial Hahn coefficients, and
$r^{-m}F(rw)=cw^m+P(w)$ with $\vH(P)>0$: higher terms of $f$ acquire positive
relative exponents $(j-m)\rho>0$, and the perturbation terms have relative exponent
at least $\vH(E)-m\rho>0$. On an ordinary annulus around $|w|=1$ the leading
coefficient $cw^m$ is nowhere zero, so geometric inversion gives a common-support
holomorphic Hahn family --- an element of $\HH$ of that annulus --- for
$rH(rw)/F(rw)$. Expanding instead around $f(rw)$,
\[
  \frac{rH(rw)}{F(rw)}
  =\sum_{k\ge0}(-1)^k\frac{rH(rw)E(rw)^k}{f(rw)^{k+1}},
\]
which is strongly summable on the annulus since $E(rw)/f(rw)$ has strictly positive
Hahn valuation there. Coefficientwise integration extracts the coefficient of
$w^{-1}$. For an ordinary germ $q$ the formal Laurent change of variable gives
\[
  [w^{-1}]\frac{rq(rw)}{f(rw)^{k+1}}=[z^{-1}]\frac{q(z)}{f(z)^{k+1}} ,
\]
an identity that holds for a Hahn family of germs as well, by finite contributions
after rescaling: in a term $g(z)\dd z$ of the ordinary meromorphic germ expansion,
scaling sends $z^j\dd z$ to $r^{j+1}w^j\dd w$, and the circle residue retains only
$j=-1$, with no surviving scale factor. Applying this to $q=HE^k$ and summing the
strongly summable family gives exactly \eqref{contours:eq:residueone}. Since the
right side does not involve $\rho$, radius independence follows, and the support
certificate is the one carried by \eqref{contours:eq:residueone}.
\end{proof}

Radius independence here is an \emph{exact residue identity}. It does not require a
convergent sequence of radii or a fine-continuous homotopy between scales whose
ratio might be infinite.

\begin{theorem}[Radius-free moving-pole residue theorem]
\label{contours:thm:localresidue}
Let $N,G\in\A_1$ with $G\ne0$. Normalize its leading exponent and choose a
protecting microscopic circle as in \cref{contours:lem:protect} when its leading
germ vanishes at zero; if the leading germ is a unit, choose any sufficiently small
infinitesimal circle. Then
\begin{equation}\label{contours:eq:localresidue}
  \frac1{2\pi i}\HCoint_{|z|=r}\frac{N(z)}{G(z)}\dd z
  =\sum_{a\in\m_K}\res_a\Bigl(\frac NG\dd z\Bigr),
\end{equation}
a finite sum over actual poles, all of which lie strictly inside the chosen
normalized circle. If $G=f+E$ with $f$ of order $m$, there are exactly $m$
denominator zeros counted with multiplicity, and the value agrees with
\eqref{contours:eq:residueone}.
\end{theorem}
\begin{proof}
If the leading germ is a unit then $G$ is a unit of $\A_1$ and the assertion follows
from \cref{contours:thm:localcauchy} or from coefficientwise primitives. Otherwise
Hahn-germ preparation and division \cite[\S4]{Geometry} give $N/G=J+R/P$ with
$J\in\A_1$, $P$ monic of degree $m$ with all roots infinitesimal, and $\deg R<m$,
the normalization absorbing the leading monomial and unit factors into the
numerator. These roots are the zeros of $G$ with multiplicities and satisfy the
protecting bound of \cref{contours:lem:protect}. Factor $P$ over the algebraically
closed $K$ and use finite partial fractions
$R/P=\sum_a\sum_{j=1}^{m_a}c_{a,j}(z-a)^{-j}$. The $J$ integral is zero. For each
$a$ the normalized position $a/r$ is infinitesimal, so expansion on $z=rw$,
$|w|=1$, is strongly summable, with normalized integral $1$ for $(z-a)^{-1}\dd z$
and $0$ for every $(z-a)^{-j}\dd z$ with $j>1$. Hence the integral is
$\sum_ac_{a,1}$. Preparation accounts for all infinitesimal zeros, and
\cref{contours:lem:protect} places them inside the circle; agreement with
\eqref{contours:eq:residueone} is \cref{contours:thm:microcircle}, and the same
finite set of local residues proves independence of $\rho$.
\end{proof}

This extends the common-domain moving-pole theorem \eqref{contours:eq:movingres} to
the radius-free setting without requiring a common ordinary representative for the
original infinitely many coefficient germs. It is \emph{not} a claim that the same
germ can be evaluated at arbitrary noninfinitesimal points, nor that its original
coefficient functions have acquired a common ordinary analytic radius. An arbitrary
Hahn family of meromorphic germs is a larger object than a quotient of two elements
of $\A_1$, and \emph{we do not assert that every such family admits the microscopic
contour construction.}

\begin{example}[An explicit radius-free contour value]
\label{contours:ex:radiusfreevalue}
Take $H$ from \eqref{contours:eq:shrinking} and $F(z)=z$. For any $\rho>0$,
\cref{contours:thm:microcircle} gives
\begin{equation}\label{contours:eq:radiusfreevalue}
  \frac1{2\pi i}\HCoint_{|z|=t^\rho}\frac{H(z)}{z}\dd z
  =H(0)=\sum_{j\ge1}t^j=\frac{t}{1-t} .
\end{equation}
After scaling the integrand is
$\sum_{j\ge1}\sum_{k\ge0}j^kt^{j+k\rho}w^{k-1}\dd w$; at each exponent the
coefficient is a finite Laurent polynomial and the residue selects exactly $k=0$.
This calculation is valid also when $\Gamma$ has higher rank. The output is a strong
sum, not a claim about convergence of its partial sums in any topology. With
$\Gamma=\Q$, $a=t^2$ and $r=t^{1/2}$ the same mechanism gives
\begin{equation}\label{contours:eq:radiusfreecauchy}
  \frac1{2\pi i}\HCoint_{|z|=t^{1/2}}\frac{H(z)}{z-t^2}\dd z
  =\sum_{j\ge1}\frac{t^j}{1-jt^2}
  =t+t^2+2t^3+3t^4+5t^5+9t^6+16t^7+31t^8+\cdots,
\end{equation}
the coefficient at exponent $\ell$ being the finite sum of $j^k$ over $j\ge1$,
$k\ge0$ with $j+2k=\ell$.
\end{example}

\begin{example}[A radius-free numerator over a perturbed denominator]
\label{contours:ex:radiusfreepole}
Take the numerator \eqref{contours:eq:shrinking} and $F(z)=z^2-t$, so $m=2$ and
$\vH(E)=1$. For the protecting radius $r=t^{1/4}$,
\cref{contours:thm:microcircle} gives
\begin{align}\label{contours:eq:radiusfreeexample}
  \frac1{2\pi i}\HCoint_{|z|=r}\frac{H(z)}{z^2-t}\dd z
  &=\sum_{k\ge0}t^k[z^{2k+1}]H(z)\notag\\
  &=\sum_{j\ge1}\sum_{k\ge0}j^{2k+1}t^{j+k}
   =\sum_{j\ge1}\frac{j\,t^j}{1-j^2t}\notag\\
  &=t+3t^2+12t^3+64t^4+441t^5+3855t^6+\cdots .
\end{align}
Every rearrangement is strongly valid: the coefficient at $t^\ell$ has only the
pairs $j+k=\ell$ with $j\ge1$, $k\ge0$. Equivalently the actual simple-pole formula
$\bigl(H(t^{1/2})-H(-t^{1/2})\bigr)/(2t^{1/2})$ gives $jt^j/(1-j^2t)$ for each
summand of $H$, proving agreement. \emph{No common ordinary disk for the original
numerator coefficients exists}, so the period is genuinely supplied by the
microscopic representation and not by an unavailable fixed-circle integral.
\end{example}

\section{Microscopic tori and general complete-intersection residues}
\label{contours:sec:multi}

\subsection{Coordinate-power leading systems}
Let
\begin{equation}\label{contours:eq:powerleading}
  F_j(z)=z_j^{d_j}+E_j(z),\qquad d_j\ge1,\quad E_j\in\A_n,\quad \vH(E_j)=\delta_j>0
\end{equation}
for each nonzero $E_j$ (with $\delta_j=+\infty$ if $E_j=0$). Choose positive
$\rho_j\in\Gamma$ with
\begin{equation}\label{contours:eq:protecting}
  d_j\rho_j<\delta_j\quad\text{whenever }E_j\ne0,
  \qquad r_j=t^{\rho_j} ;
\end{equation}
such choices exist by divisibility of $\Gamma$, and for a zero perturbation there is
no upper restriction on $\rho_j$. Write $T_r$ for the parameterized real $n$-torus
$z_j=r_je^{i\theta_j}$, $0\le\theta_j\le2\pi$, with ordinary parameter angles and
orientation $\dd\theta_1\wedge\cdots\wedge \dd\theta_n$; integration is defined by
pullback to the ordinary product of unit circles. \emph{It is a real
$n$-dimensional cycle in an $n$-complex-variable calculation, not the
$(2n-1)$-dimensional sphere used in some other several-variable integral
formulas.}

\begin{theorem}[Coordinate-power torus realization]\label{contours:thm:microtorus}
For $F$ as in \eqref{contours:eq:powerleading}, radii as in
\eqref{contours:eq:protecting}, and $H\in\A_n$, the integral
\begin{equation}\label{contours:eq:microtorus}
  \mathcal R_F(H):=\frac1{(2\pi i)^n}\HInt_{T_r}
  \frac{H(z)\dd z_1\wedge\cdots\wedge \dd z_n}{F_1(z)\cdots F_n(z)}
\end{equation}
is well defined and satisfies the exact coefficient formula
\begin{equation}\label{contours:eq:toruscoeff}
 \boxed{\quad
 \mathcal R_F(H)=\sum_{k\in\N^n}(-1)^{|k|}
 \bigl[z_1^{d_1(k_1+1)-1}\cdots z_n^{d_n(k_n+1)-1}\bigr]
 \Bigl(H\prod_{j=1}^nE_j^{k_j}\Bigr)
 =\Res_F(H),
 \quad}
\end{equation}
the perturbation residue \eqref{contours:eq:generalres} for $f_j=z_j^{d_j}$. Its
support is contained in $\supp H+E^*$, where $E=\bigcup_j\supp E_j$, and it is
independent of all protecting radii satisfying
\eqref{contours:eq:protecting}. Every infinitesimal common zero satisfies
$v(z_j)>\rho_j$ for every $j$.
\end{theorem}
\begin{proof}
After substituting $z_j=r_jw_j$ one has
$r_j^{-d_j}F_j(rw)=w_j^{d_j}+P_j(w)$ with $\vH(P_j)\ge\delta_j-d_j\rho_j>0$, and all
coefficient functions are polynomials by \cref{contours:thm:rescaling}. On a common
ordinary product annulus about $|w_j|=1$ all leading $w_j^{d_j}$ are nonzero, so
each reciprocal has a strong geometric expansion; \cref{contours:lem:support} proves
strong summability of the entire pulled-back form, which may therefore be integrated
coefficientwise. At multi-index $k$ the denominator contributes
$\prod_jz_j^{-d_j(k_j+1)}$ and the numerator is $H\prod_jE_j^{k_j}$; the ordinary
iterated torus integral selects the Laurent monomial $w_1^{-1}\cdots w_n^{-1}$, so a
source Taylor monomial $z^\alpha$ contributes only when
$\alpha_j=d_j(k_j+1)-1$ for every $j$, and for that monomial the scale factor in
coordinate $j$ is $r_j^{\,1+\alpha_j-d_j(k_j+1)}=1$. The surviving coefficients are
exactly \eqref{contours:eq:toruscoeff}, including the sign. Ordinary Taylor
coefficient operators do not enlarge Hahn support and the perturbations have
positive support, giving the support statement; the formula contains no $r_j$,
giving radius independence. Finally, for an infinitesimal common zero,
$d_jv(z_j)=v(E_j(z))\ge\delta_j$ whenever $z_j\ne0$, and a zero coordinate has
valuation $+\infty$.
\end{proof}

The theorem permits \emph{different valuation ranks in different coordinates}: for
example $\rho_1\in\Q_{>0}$ and $\rho_2$ proportional to $\omega$ are both
legitimate, and nothing in the proof replaces $\Gamma$ by a real-valued norm. It
does not assume that the perturbations are polynomial, nor that their coefficient
germs share an ordinary neighborhood. It supplies a contour meaning to a residue
that was deliberately defined algebraically in \cite{Geometry}.

\begin{example}[A two-variable torus and cancellation across valuation ranks]
\label{contours:ex:coupled}
Take $F_1=x^2-Ay$, $F_2=y^2-Bx$ with $A,B\in\m_K\setminus\{0\}$, and choose
$2\rho_x<v(A)$, $2\rho_y<v(B)$. Then
\begin{equation}\label{contours:eq:coupledres}
  \mathcal R_F(H)=\sum_{k,\ell\ge0}A^kB^\ell\,
  [x^{2k+1-\ell}y^{2\ell+1-k}]H ,
\end{equation}
a coefficient with a negative exponent being zero. In particular
$\mathcal R_F(1)=0$, $\mathcal R_F(xy)=1$, and $\mathcal R_F(4xy-AB)=4$. The
Jacobian is $J_F=4xy-AB$, equal to $-AB$ at the origin and to $3AB$ at the three
roots with $xy=AB$, so the four weights for numerator $1$ are $-1/(AB)$ and
$1/(3AB)$ three times, summing to zero. For $A=t$ and $B=t^\omega$ this cancellation
involves weights of valuation $-(1+\omega)$, whereas the torus result has no
negative term at all.
\end{example}

\subsection{Arbitrary isolated analytic complete intersections}
\label{contours:subsec:general}
Now let $f=(f_1,\ldots,f_n)$ be an ordinary analytic germ at zero with an isolated
common zero, i.e. $0<\mu=\dim_\C\C\{z\}/(f_1,\ldots,f_n)<\infty$, and let $F=f+E$
with all nonzero $E_j\in\A_n$ of strictly positive Hahn valuation. The perturbation
residue is \eqref{contours:eq:generalres}, where the ordinary local Grothendieck
residue is a complex-linear functional on analytic germs whose coordinate-power case
is Taylor coefficient extraction. We use the classical residue transformation law as
presented in \cite{Tsikh,CDS}.

Since $(f)$ has finite colength, choose $d_1,\ldots,d_n\ge1$ and an \emph{ordinary
analytic} matrix $A(z)$ with
\begin{equation}\label{contours:eq:transformmatrix}
  A(z)f(z)=(z_1^{d_1},\ldots,z_n^{d_n})^{\mathsf T},
\end{equation}
and set $G=AF$, so that $G_j=z_j^{d_j}+D_j$ with $D=AE$ and all nonzero $D_j$ of
positive Hahn valuation. \emph{No invertibility of $A$ is assumed.} The existence of
such $A$ and $d_j$ uses finite colength alone; it is ordinary local analytic
algebra and requires nothing from the geometry manuscript beyond the definition of
$\Res_F$.

\begin{lemma}[Transformation of the Hahn perturbation residue]
\label{contours:lem:residuetransform}
With these hypotheses, for every $H\in\A_n$,
\begin{equation}\label{contours:eq:residuetransform}
  \Res_F(H)=\Res_G(H\det A).
\end{equation}
\end{lemma}
\begin{proof}
We first specify the classical input. For a finite analytic parameter deformation
$f_u$ of an isolated complete intersection, the sum of local residues over the
nearby cluster is analytic in $u$; if $g_u=Af_u$ also has an isolated central
fibre, the classical transformation law gives
\[
  \res^{\mathrm{cluster}}_{f_u}(h)=\res^{\mathrm{cluster}}_{g_u}(h\det A).
\]
The cluster for $g_u$ can contain additional points, and \emph{the determinant is
part of the law}. Taylor expansion in $u$ produces the geometric denominator
expansions appearing in \eqref{contours:eq:generalres}. This ordinary finite-family
statement follows from the usual local residue transformation law and its parameter
version \cite{Tsikh,CDS}.

Now fix a Hahn exponent $\nu$. On either side of
\eqref{contours:eq:residuetransform}, a coefficient is built from a source
coefficient of $H$ together with a finite word of positive perturbation exponents.
\Cref{contours:lem:support} leaves only finitely many contributing words, hence
finitely many ordinary coefficient germs --- this is exactly the place where the
distinction recorded after \cref{contours:ex:rank} is load-bearing, since
infinitely many \emph{exponents} may lie below $\nu$. Replace the finitely many
retained perturbation terms by independent ordinary parameters. Together with the
finitely many entries of $A$, their germs have a common ordinary representative
neighborhood. Apply the finite-family identity there, compare the finitely many
required Taylor coefficients, and substitute the monomials $t^e$ back for the
parameters; exponent collisions cause only finite additions. Thus the two sides have
equal coefficient at $\nu$, and $\nu$ was arbitrary. A fully detailed truncation
argument for a transformation matrix that is only \emph{formal} in the parameters is
given in \cref{contours:app:transform}.
\end{proof}

\begin{theorem}[Microscopic contour representation of a general residue]
\label{contours:thm:generalrepresentation}
Choose $A$ and $d_j$ as in \eqref{contours:eq:transformmatrix} and protecting radii
for $G$ as in \eqref{contours:eq:protecting}. Then for every $H\in\A_n$
\begin{equation}\label{contours:eq:generalrepresentation}
 \boxed{\qquad
 \Res_F(H)=\frac1{(2\pi i)^n}\HInt_{T_r}
 \frac{H(z)\det A(z)\dd z_1\wedge\cdots\wedge \dd z_n}{G_1(z)\cdots G_n(z)} .
 \qquad}
\end{equation}
The value is independent of the admissible radii and of the choice of ordinary
transformation data $A,d_j$.
\end{theorem}
\begin{proof}
The integrand is admissible by \cref{contours:thm:rescaling}, and
\cref{contours:thm:microtorus} identifies its normalized torus integral with
$\Res_G(H\det A)$. Apply \cref{contours:lem:residuetransform}. Every permitted
choice gives the same independently defined left side, proving independence.
\end{proof}

\emph{Two non-claims must survive this theorem, and are stated here as part of it.}
First, \eqref{contours:eq:generalrepresentation} is a contour representation
\emph{by a transformed form}. It does \emph{not} assert that the untransformed form
$H\dd z_1\wedge\cdots\wedge \dd z_n/(F_1\cdots F_n)$ has the same integral over that
coordinate torus; in the classical theory the natural Grothendieck cycle is adapted
to the equations, typically $\{|F_j|=\epsilon_j\}$, and replacing that cycle by a
coordinate torus without the determinant transformation would in general be wrong.
Second, the theorem does \emph{not} claim that $F$ and $G$ have identical zero
schemes: multiplication by a matrix with nonunit determinant can introduce
additional zeros and multiplicities, and \emph{$\det A$ compensates for this in the
residue identity --- it is not an optional Jacobian factor.} Formula
\eqref{contours:eq:generalrepresentation} is useful precisely because it gives a
controlled coordinate-torus substitute.

The logical cost of this theorem should be recorded exactly. It uses the strong
residue expansion \eqref{contours:eq:generalres}, ordinary finite-family residue
transformation, and positive rescaling \cref{contours:thm:rescaling}. \emph{It does
not require the Nullstellensatz or faithful-flatness theorems of \cite{Geometry},
nor that manuscript's finite-algebra normal form.}

\subsection{Relation to the finite-algebra conclusions of the supplied source}
The geometry manuscript proves that the perturbed quotient
$B_F=\A_n/(F_1,\ldots,F_n)$ has $K$-dimension exactly $\mu$, and that
\begin{equation}\label{contours:eq:traceresidue}
  \Res_F(HJ_F)=\Tr(M_H)=\sum_ae_aH(a),
  \qquad J_F=\det\Bigl(\frac{\partial F_i}{\partial z_j}\Bigr),
  \qquad \sum_ae_a=\mu ,
\end{equation}
with $\Res_F$ annihilating $(F)$ and $e_a$ the local multiplicity at the
\emph{actual displaced} zero. For simple zeros this gives
$\Res_F(H)=\sum_aH(a)/J_F(a)$. These are reported companion results, not reproved
here as part of the contour construction. Combining them with
\cref{contours:thm:generalrepresentation} identifies one functional simultaneously
as a finite-algebra trace pairing and as a microscopic integral, and gives
\begin{equation}\label{contours:eq:torustrace}
  \frac1{(2\pi i)^n}\HInt_{T_r}
  \frac{H\,J_F\det A}{G_1\cdots G_n}\dd z_1\wedge\cdots\wedge \dd z_n
  =\sum_ae_aH(a),
\end{equation}
whose value at $H=1$ is the ordinary integer $\mu$. Thus this higher-dimensional
contour has a finite local-degree interpretation even when the actual points split
across several infinitesimal ranks. On a common scaled neighborhood avoiding the
divisor $G_1\cdots G_n=0$ the displayed meromorphic top form is holomorphic and
closed, so its admissible deformations obey
\cref{contours:thm:deformedstokes} and their periods agree when the intervening
chain avoids that divisor --- a genuine Stokes explanation of contour deformation in
addition to the algebraic radius-independence proof. This is \emph{not} a statement
about arbitrary fine-continuous hypersurfaces in the full class plane.

In dimension one, Stokes, the residue of a Laurent series, and the integer winding
of a circle already organize the calculation. In higher dimensions the determinant
transformation and the real $n$-torus play the corresponding roles. \emph{There is
no need to postulate a new Jordan separation theorem for a $(2n-1)$-sphere in order
to justify these residues.}

\subsection[Scope: the transformed representation, the objection, and the
resolution]{Scope: the transformed representation, the objection it drew, and the
resolution of the untransformed question}
\label{contours:subsec:scope}
This subsection records a disagreement about \emph{status} among the three
manuscripts originally merged here, and then records what became of the question
that disagreement was about. \emph{The history is kept deliberately: the
objection below was raised, it was correct about the coordinate torus, and it
identified by name a problem that a fourth manuscript later solved with a
different cycle.} Nothing in the resolution reaches back and strengthens
\cref{contours:thm:generalrepresentation}, whose own non-claim about the
untransformed form stands exactly as printed above. The disagreement about status
is the sharpest in the set, and it is not averaged.
\Cref{contours:lem:residuetransform,contours:thm:generalrepresentation} above are
proved unconditionally and are printed with their proof. One of the three sources
judged the corresponding statement \emph{unproved} and declined to assert it, and a
third reached a determinant torus formula only conditionally. Both positions are
reproduced here, because dropping the objection would overstate how much is proved
unconditionally, while dropping the theorem would understate it.

\paragraph{The objection, as raised.} For arbitrary isolated $f$, a single diagonal
rescaling need not make the leading equations coordinate powers, or even keep the
leading homogeneous system isolated. Thus \cref{contours:thm:microtorus} cannot
simply be applied after deleting its hypothesis. Moreover the usual cycle
$\{|f_i|=\epsilon_i\}$ is not in general a coordinate product torus. What is
available without further argument is an exact \emph{coefficient-dependent cycle
interpretation} of \eqref{contours:eq:generalres}: at a prescribed Hahn exponent,
\cref{contours:lem:support} retains only finitely many numerator germs, perturbation
germs, and multi-indices; those finitely many germs have a common ordinary
representative neighborhood; their ordinary residues can be computed on suitable
protecting cycles there; and the answer is independent of the representatives and
radii by ordinary local residue theory. Doing this at each exponent gives precisely
\eqref{contours:eq:generalres}. \emph{This is a compatible germwise procedure, not
one fixed ordinary cycle for the entire original infinite family.} On that basis the
source in question wrote: ``The present article proves that realization for one
variable and for the coordinate-power class above; it does not assume it for every
complete intersection,'' and listed a general realization theorem as an open
problem: given arbitrary isolated $f$ and radius-free positive-support $E$,
construct \emph{one} admissible $n$-cycle whose integral equals
\eqref{contours:eq:generalres}, possibly requiring a non-diagonal analytic change of
variables, a finite resolution-type decomposition, or chains on an enlarged analytic
space, and prove that no omitted summation or radius assumption is hidden in the
parameterization. \emph{That is the problem now solved in
\cref{contours:sec:untransformed,contours:sec:adapted}}; the last paragraph of
this subsection states under exactly which hypotheses.

\paragraph{What the objection does and does not refute.} It does not refute
\cref{contours:thm:generalrepresentation}, and the two are consistent, because they
answer different questions. The open problem asks for a cycle on which the
\emph{untransformed} integrand $H\dd z/(F_1\cdots F_n)$ is integrated;
\cref{contours:thm:generalrepresentation} supplies instead a coordinate torus for
the \emph{transformed} integrand $H\det A\dd z/(G_1\cdots G_n)$, and, as stated
immediately after that theorem, it explicitly does not claim the former. The
objection is therefore best read as a statement about which of the two
representations has been achieved. Nor does the objection's existence prove that no
single suitably nonstandard cycle is possible. Conversely,
\cref{contours:thm:generalrepresentation} does settle the representation question in
the transformed form, unconditionally, from finite colength alone. A reader should
take away both: a general isolated complete intersection \emph{does} have a
controlled microscopic contour representation, and it is \emph{not} the naive one.
The objection's positive part is also confirmed rather than overturned by what
follows: \cref{contours:ex:wrongtorus} exhibits an ordinary germ for which the
original coordinate product torus returns $0$ while the residue is $1/2$, so no
theorem about the untransformed integrand on a coordinate torus was ever
available. It is the \emph{cycle} that had to change, not
\cref{contours:thm:generalrepresentation}.

\paragraph{The conditional third route.} A second and genuinely independent route to
a determinant torus formula runs through univariate separating denominators. It is
worth recording as an alternative proof, \emph{provided its dependence is stated}:
unlike \cref{contours:thm:generalrepresentation}, it assumes the finite-algebra and
support-controlled normal-form results of the companion geometry manuscript
\cite{Geometry}, which is an unrefereed research manuscript. It is therefore
conditional on that companion rather than independent of it.

\begin{theorem}[Separating-torus realization; conditional on \cite{Geometry}]
\label{contours:thm:septorus}
\emph{Assume the finite-algebra and support-controlled normal-form results of
\cite[\S\S7--9]{Geometry}}: that $B_F=\A_n/(F_1,\ldots,F_n)$ has $K$-dimension
$\mu$, that any chosen ordinary monomial basis remains a basis, and that
multiplication matrices and normal-form operations have explicit support in $E^*$.
Let $M_i$ be the multiplication matrix of $z_i$ on $B_F$ in that basis, and put
\begin{equation}\label{contours:eq:Qi}
  Q_i(T)=\det(TI-M_i)=T^\mu+\sum_{j<\mu}q_{ij}T^j .
\end{equation}
Its exponent-zero matrix is multiplication by $z_i$ in the ordinary local algebra
$\C\{z\}/(f)$, hence nilpotent, so every nonzero $q_{ij}$ has positive valuation. By
Cayley--Hamilton $Q_i(z_i)\in(F)$; applying the supplied normal form to it, with
zero remainder, produces a matrix $\mathsf C=(C_{ij})$ with
\begin{equation}\label{contours:eq:matrixC}
  Q_i(z_i)=\sum_{j=1}^nC_{ij}(z)F_j(z),\qquad C_{ij}\in\A_n,
\end{equation}
all entries having nonnegative support contained in $E^*$. Choose
$\rho_i\in\Gamma_{>0}$ with
\begin{equation}\label{contours:eq:torusradius}
  (\mu-j)\rho_i<v(q_{ij})\ \text{ for every nonzero }q_{ij},
  \qquad r_i=t^{\rho_i} .
\end{equation}
Then for every $H\in\A_n$
\begin{equation}\label{contours:eq:septorus}
 \boxed{\quad
 \Res_F(H)=\frac1{(2\pi i)^n}\HInt_{T_r}
 \frac{H(z)\det\mathsf C(z)}{\prod_{i=1}^nQ_i(z_i)}
 \dd z_1\wedge\cdots\wedge \dd z_n ,
 \quad}
\end{equation}
independently of the allowed radii and of the chosen representing matrix
$\mathsf C$, with support contained in $\supp H+E^*$.
\end{theorem}
\begin{proof}
\emph{Existence of the integral.} By \cref{contours:thm:rescaling}, $H(rw)$ and every
entry of $\mathsf C(rw)$ have polynomial complex coefficients at each Hahn exponent.
The scaled denominator is
$r_i^{-\mu}Q_i(r_iw_i)=w_i^\mu+\sum_{j<\mu}q_{ij}r_i^{\,j-\mu}w_i^j$, whose
non-leading coefficients have strictly positive valuation by
\eqref{contours:eq:torusradius}; hence each reciprocal is a strong geometric
expansion near the ordinary product unit torus, and products with the differential
factors $r_i\dd w_i$ remain strongly supported. Moreover every root $b$ of $Q_i$
satisfies $v(b)>\rho_i$, since otherwise $b^\mu$ has strictly smaller valuation than
each $q_{ij}b^j$ and cannot cancel; so every common zero of $F$ lies strictly inside
the separating coordinate circles.

\emph{Residue transformation.} The classical rule
$\Res_F(H)=\Res_Q(H\det\mathsf C)$ for $Q=\mathsf CF$ is standard residue theory
\cite[\S0]{CDS}. Its justification here is the finite-parameter argument of
\cref{contours:app:transform}, which is needed because $\mathsf C$ arises from a
linear splitting that need not be continuous in an analytic norm, so
$\mathsf C(u)$ cannot be asserted convergent as an analytic parameter series.

\emph{The separated residue is a torus integral.} For the leading tuple
$(z_1^\mu,\ldots,z_n^\mu)$ the ordinary residues are coefficient extraction;
expanding each $Q_i^{-1}$ around $z_i^{-\mu}$ shows that $\Res_Q$ is the iterated
coefficient of $z_1^{-1}\cdots z_n^{-1}$, and after the admissible scaling,
integration of an ordinary Laurent monomial on each unit circle extracts exactly
those coefficients, the powers of $r_i$ cancelling. All rearrangements are strongly
summable by the verified support bounds. The left side is already independent of the
radii and of $\mathsf C$, so equality gives the remaining assertions; no uniqueness
of $\mathsf C$ is needed.
\end{proof}

\begin{remark}[What was changed, and what was not]
\label{contours:rem:toruslimits}
Formula \eqref{contours:eq:septorus} does \emph{not} assert that the original
quotient $H/(F_1\cdots F_n)$ is regular on the chosen product torus: a hypersurface
$F_i=0$ can meet that torus. The denominators were deliberately replaced by the
univariate $Q_i$ and $\det\mathsf C$ was inserted. \emph{Omitting the determinant is
generally wrong.} Nor has a real $n$-torus been identified with the full boundary of
a complex polydisk; the integration cycle and its real dimension are exactly those
specified. Using univariate polynomials in a zero-dimensional ideal to compute
multidimensional residues is a classical method \cite[\S0]{CDS}; the additional
conclusion here is that the method admits a legitimate support-controlled contour
interpretation for $\A_n$. No claim is made that separating denominators or residue
transformation are new ideas.
\end{remark}

\emph{Which representation is sharper, and for what.}
\Cref{contours:thm:generalrepresentation} is sharper for \emph{logical economy}: it
is unconditional, needing only finite colength, ordinary residue transformation, and
\cref{contours:thm:rescaling}, so it survives if the companion manuscript's
finite-algebra package is withdrawn. \Cref{contours:thm:septorus} is sharper for
\emph{computation and for zero-scheme bookkeeping}: its denominators are explicit
univariate characteristic polynomials whose roots and admissible radii are read off
from \eqref{contours:eq:torusradius}, and the vanishing of $\det\mathsf C$ at the
spurious points of the separated system is visible, as
\cref{contours:ex:fourpoints} shows. They are not competing proofs of one statement;
they certify different things.

\begin{example}[A four-point cluster and its separating torus]
\label{contours:ex:fourpoints}
Take $F_1=x^2-Ay$, $F_2=y^2-Bx$ with $A,B\in\m_K\setminus\{0\}$ as in
\cref{contours:ex:coupled}. The ordinary leading algebra is $\C\{x,y\}/(x^2,y^2)$
with basis $1,x,y,xy$ and $\mu=4$. The coordinate characteristic polynomials are
\[
  Q_x=x^4-A^2Bx,\qquad Q_y=y^4-AB^2y ,
\]
and an explicit transformation matrix with
$(Q_x,Q_y)^{\mathsf T}=\mathsf C\,(F_1,F_2)^{\mathsf T}$ is
\begin{equation}\label{contours:eq:exampleC}
  \mathsf C=\begin{pmatrix}x^2+Ay&A^2\\ B^2&y^2+Bx\end{pmatrix},
  \qquad
  D=\det\mathsf C=x^2y^2+Bx^3+Ay^3+ABxy-A^2B^2 .
\end{equation}
Hence
\begin{equation}\label{contours:eq:exampletorus}
  \Res_F(H)=\frac1{(2\pi i)^2}\HInt_{T_r}
  \frac{H(x,y)\,D(x,y)}{(x^4-A^2Bx)(y^4-AB^2y)}\dd x\wedge \dd y ,
\end{equation}
admissible for any positive exponents with $3\rho_x<2v(A)+v(B)$ and
$3\rho_y<v(A)+2v(B)$; with $A=t$ and $B=t^\omega$ both $\rho_x=\rho_y=1/4$ work. The
roots are $(0,0)$ and $(\varrho\zeta^j,\varrho^2\zeta^{2j}/A)$ for $j=0,1,2$ with
$\varrho^3=A^2B$ and $\zeta$ a primitive cube root of unity. The weights
$-1/(AB),\,1/(3AB),\,1/(3AB),\,1/(3AB)$ are individually infinite and sum to zero,
while $\Res_F(xy)=1$ and $\Res_F(J_F)=4$. In the quotient algebra $D=ABJ_F$, the
exact polynomial identity having a remainder in $(F_1,F_2)$.

The determinant's role is especially explicit here. \emph{The separated system
$(Q_x,Q_y)$ has sixteen simple grid points, whereas the original system has only
four.} At each of the other twelve points $\mathsf CF=0$ with $F\ne0$, so
$\det\mathsf C=0$ and the separated simple residue vanishes. Changing the
denominators without the determinant would count the wrong zero scheme. For an
independent check, expanding
$Q_x^{-1}=x^{-4}\sum_{k\ge0}(A^2B)^kx^{-3k}$ and
$Q_y^{-1}=y^{-4}\sum_{\ell\ge0}(AB^2)^\ell y^{-3\ell}$ on the chosen torus and
extracting the coefficient of $x^{-1}y^{-1}$ reproduces
\begin{equation}\label{contours:eq:monomialres}
  \Res_F(x^py^q)=
  \begin{cases}
    A^{(2p+q-3)/3}B^{(p+2q-3)/3},
    &\tfrac{2p+q-3}{3},\ \tfrac{p+2q-3}{3}\in\N,\\[0.3em]
    0,&\text{otherwise},
  \end{cases}
\end{equation}
in agreement with \eqref{contours:eq:coupledres}. The general equality is proved by
\cref{contours:thm:septorus}, not inferred from the finite tests.
\end{example}


\paragraph{The resolution, and its exact hypotheses.}
\label{contours:para:resolution}
The open problem quoted above is solved in
\cref{contours:sec:untransformed,contours:sec:adapted}, and the solution should
be read with its hypotheses attached rather than as a removal of the scope
discussion. What is now proved, by \cref{contours:thm:untransformed}, is that for
\begin{enumerate}
\item an ordinary \emph{convergent} analytic map germ $f$ with an isolated common
zero of finite colength $\mu$ --- the hypothesis of
\cref{contours:subsec:general}, unchanged;
\item one \emph{specified} positive valuation lower bound $\delta\in\Gamma_{>0}$
on the perturbation support, $\supp E\subset[\delta,\infty)$ --- a hypothesis
that \eqref{contours:eq:generalres} did not need, and the only genuinely new
one, since a positive well-ordered support always has a least element and each
individual $E$ therefore qualifies with $\delta=\min\supp E$; and
\item a radius-free numerator $H\in\A_n$, or a formal one by
\cref{contours:thm:formalcontour},
\end{enumerate}
there is one compact oriented ordinary parameter $n$-manifold, a finite covering
of the coordinate torus of total degree $\mu$, and one Hahn parameterization of
it, built from $f$ and $\delta$ alone, on which the \emph{untransformed} form
$H\dd z_1\wedge\cdots\wedge \dd z_n/(F_1\cdots F_n)$ integrates to
\eqref{contours:eq:generalres} for \emph{every} $E$ above the threshold. The
problem's final demand --- that no omitted summation or radius assumption be
hidden in the parameterization --- is met by the explicit support certificates
\eqref{contours:eq:untransformedsupport} and
\eqref{contours:eq:Upsilonsupport}, by the separation of the two finiteness
principles in \cref{contours:subsec:tworadii}, and by
\cref{contours:thm:degreenormalization,contours:thm:untransformedhomotopy},
which together fix the degree of the cycle and identify the periods of the
protecting and adapted cycles.

\emph{Four things remain exactly as they were.}
\Cref{contours:thm:generalrepresentation} still does not assert that the
untransformed form has the same \emph{coordinate-torus} integral; that assertion
is false, by \cref{contours:ex:wrongtorus}. The new cycle is a ramified finite
cover, in general disconnected, and its construction does not eliminate
$\det A$ from \cref{contours:thm:generalrepresentation} or $\det\mathsf C$ from
\cref{contours:thm:septorus}, both of which remain non-optional for their own
formulas. The objection's germwise, coefficient-dependent cycle interpretation
recorded above remains correct and remains the right description of what
\eqref{contours:eq:generalres} does by itself. And the question is closed only
under the three hypotheses just listed: an arbitrary Hahn tuple with no
convergent isolated leading map, a non-isolated leading zero, and a perturbation
support not bounded below by a positive element are all still outside the
theorem, as \cref{contours:subsec:untransformedscope} states in detail.


\section{One cycle for the untransformed integrand}
\label{contours:sec:untransformed}

\Cref{contours:subsec:scope} distinguished two questions and answered one of
them. \Cref{contours:thm:generalrepresentation} represents the perturbation
residue of an arbitrary isolated complete intersection by a coordinate torus
carrying the \emph{transformed} integrand $H\det A\dd z/(G_1\cdots G_n)$, and
explicitly declines the untransformed form. The remaining question --- one
admissible $n$-cycle on which the original form
$H\dd z_1\wedge\cdots\wedge \dd z_n/(F_1\cdots F_n)$ itself integrates to
\eqref{contours:eq:generalres} --- was recorded there and in
\cref{contours:subsec:problems} as an open problem.

This section and \cref{contours:sec:adapted} construct such a cycle. The
hypotheses are exactly those of \cref{contours:subsec:general} together with
\emph{one} positive lower bound on the perturbation support: a convergent
ordinary leading map with an isolated zero, a specified $\delta>0$ with
$\supp E\subset[\delta,\infty)$, and a radius-free numerator. No determinant
transformation is made, no equation is replaced, and no common convergence
radius is assumed for the coefficient germs of $E$ or of $H$. The cycle is
built from $f$ and $\delta$ alone and then serves \emph{every} perturbation
above that threshold and every numerator at once.

\subsection{What ``one cycle'' is allowed to be}
\label{contours:subsec:onecycle}
The answer depends on a decision about the chain category, and that decision is
made explicitly here rather than left implicit in the word ``cycle''.

\begin{definition}[Admissible Hahn cycle on an ordinary cover]
\label{contours:def:hahncycle}
An \emph{admissible Hahn $n$-cycle} consists of one compact oriented ordinary
smooth $n$-manifold $\mathsf T$ without boundary --- \emph{possibly a finite
disjoint union} --- together with one Hahn parameterization
$Z:\mathsf T\to K^n$ whose coordinates lie in
$\OO(\mathsf X)((t^\Gamma))$ for a fixed ordinary complex manifold
$\mathsf X\supset\mathsf T$ and have strictly positive support. Its integral is
the coefficientwise integral \eqref{contours:eq:chainintegral} of the pulled-back
top-degree form on $\mathsf T$, and the image is written $Z[\mathsf T]$.
\end{definition}

Three features of this definition are the ones that make the construction
possible, and each was already licensed by the chain formalism of
\cref{contours:sec:derham,contours:sec:deformed}: the parameter manifold is
ordinary and compact, so \cref{contours:thm:hahnstokes} applies to it; it need
not be connected, so a finite covering of the ordinary torus is a legitimate
single cycle; and the parameterization is a positive-support Hahn family, so it
is an admissible deformation of the constant map to the origin in the sense of
\cref{contours:def:deformation}. What is \emph{not} required is that
$Z[\mathsf T]$ be a connected coordinate product torus in the original
variables $z$; the counterexample of \cref{contours:ex:wrongtorus} shows that
requiring that much would make the problem unsolvable. A map
$Z:\mathsf T\to K^n$ that varies on a connected component cannot be continuous
for the full fine or the intrinsic workspace topology, by
\cref{contours:prop:fullfine,contours:prop:intrinsic}. In contrast, its positive
support makes $\st Z=0$, so it is automatically continuous for the standard-part
topology. That coarse continuity carries no information about the cycle's shape.

One point of care is that \cref{contours:lem:pullback} cannot supply the
pullback certificate here. That lemma pulls back a form whose coefficients are
ordinary smooth functions on \emph{one} common ordinary target neighborhood; the
integrands below have radius-free coefficient germs and no such neighborhood
exists, by \eqref{contours:eq:shrinking}. The certificate is supplied instead by
\cref{contours:lem:radiusfreesub}. Once the pullback has been formed on
$\mathsf T$, the ordinary-base Stokes theorems apply without change.

\paragraph{The data, and one threshold.}
Throughout \cref{contours:sec:untransformed,contours:sec:adapted} fix
\begin{equation}\label{contours:eq:untransformeddata}
  f\in\C\{z_1,\ldots,z_n\}^n,\quad f(0)=0,\quad
  \mu=\dim_\C\C\{z\}/(f_1,\ldots,f_n)<\infty,
  \qquad F=f+E,
\end{equation}
with $E\in\A_n^n$, and fix a \emph{specified} $\delta\in\Gamma_{>0}$ with
\begin{equation}\label{contours:eq:threshold}
  S_E:=\supp E=\bigcup_{j=1}^n\supp E_j\ \subset\ [\delta,\infty).
\end{equation}
The set $S_E$ is well ordered, and $S_E^*$ is the monoid written $E^*$ in
\cref{contours:def:strongsum} and in \cref{contours:thm:microtorus}; in the
per-coordinate notation $\delta_j=\vH(E_j)$ of \eqref{contours:eq:powerleading},
condition \eqref{contours:eq:threshold} says $\delta\le\min_j\delta_j$, but it is
stated as a single threshold on the whole support because a single threshold is
what the cycle is built from. Since a
nonempty well-ordered subset of $\Gamma_{>0}$ has a least element, every
individual positive-support perturbation satisfies
\eqref{contours:eq:threshold} with $\delta=\min S_E$; the force of stating
$\delta$ first is that one cycle then serves the whole valuation ball of
perturbations with support above $\delta$. \emph{No bound of any kind is imposed
on the complex coefficients of the germs $E_{j,\lambda}$ or $H_\lambda$, and
their convergence radii may tend to zero.} As in
\cref{contours:subsec:conventions}, $\vH$ denotes the least exponent carrying a
nonzero coefficient \emph{function}; it is not a pointwise minimum of valuations
over a parameter domain, and no such pointwise notion is used.

\paragraph{The fixed-manifold coherent ring.}
The construction needs the common-domain ring of
\cref{contours:subsec:rings} with the ordinary open $U\subset\C$ replaced by a
fixed ordinary complex manifold. For such an $\mathsf X$, possibly
disconnected, write
\begin{equation}\label{contours:eq:HXring}
  \HH(\mathsf X)=\OO(\mathsf X)((t^\Gamma)),
\end{equation}
and $\HH(\mathsf X)_+$ for the elements of strictly positive support. The ring
discipline of \cref{contours:subsec:rings} is unchanged: every coefficient is
holomorphic on \emph{one} fixed $\mathsf X$ and all share one well-ordered
support. No theorem of \cref{contours:sec:cauchy} is transferred to
$\HH(\mathsf X)$ for a higher-dimensional or disconnected $\mathsf X$, and
nothing below asserts a sheaf property for $\HH$ on arbitrary covers.

\paragraph{Notation reconciled with the source.}
The source manuscript of these two sections writes $m_j$ for the target radial
weights, $M$ for their sum, $b$ and $c$ for the two ordinary radial bounds,
$T$ and $X$ for the parameter manifold and its ambient complex manifold,
$\langle P\rangle_\N$ for the monoid generated by $P$, and $\mathcal R_F$ for
the perturbation residue. All six are renamed here to preserve the conventions
already fixed in this article. The weights are $\kappa_j$ with
$\kappa_*=\max_j\kappa_j$ and $\kappa_\Sigma=\sum_j\kappa_j$, because $m$ is
reserved for the order of a one-variable germ and $d_j$ for coordinate powers
by \cref{contours:subsec:conventions}. The radial bounds are
$b_{\mathrm{pole}}$ and $c_{\mathrm{sep}}$. The manifolds are $\mathsf T$ and
$\mathsf X$, because $T_r$ is the coordinate torus of
\cref{contours:thm:microtorus} and $T_F$ the endpoint correction
\eqref{contours:eq:endpointT}. The monoid is $P^*$. The residue is
$\Res_F$ throughout, \emph{never} $\mathcal R_F$: in this article
$\mathcal R_F$ already denotes the coordinate-torus functional
\eqref{contours:eq:microtorus}, and the entire content of the theorems below is
that a \emph{different} cycle computes $\Res_F$.

\subsection{Substitution of a positive-support Hahn map, and a nonlinear
support lemma}
\label{contours:subsec:hahnsub}
\Cref{contours:lem:support} is the only summability input, and the sharpening
recorded after \cref{contours:ex:rank} --- finitely many dependencies at an
\emph{exact} exponent, not finitely many exponents below it --- is used at every
step. The source proves \cref{contours:lem:support} by a second route, through
Higman's lemma on words over a well-quasi-ordered alphabet \cite{Higman}: a
well-order is a well-quasi-order, an embedding of words with each matched letter
no larger than its partner makes the source sum at most the target sum;
equality forces identical words because every letter is strictly positive, and both a descending
sequence of word sums and an infinite family of distinct words with one fixed sum
therefore contradict Higman's lemma. That is an independent proof of the same
classical statement \cite{Neumann,Higman}; the statement itself is not restated
here.

Two genuinely new lemmas are needed. The first replaces
\cref{contours:thm:rescaling} in a situation where the substituted map is not a
diagonal monomial rescaling but an arbitrary positive-support Hahn family of
holomorphic maps.

\begin{lemma}[Radius-free substitution into a Hahn map]
\label{contours:lem:radiusfreesub}
Let $\mathsf X$ be a fixed ordinary complex manifold, let
$W=(W_1,\ldots,W_n)\in\HH(\mathsf X)_+^n$, and let $P_W$ be the union of the
supports of its coordinates, a well-ordered subset of $\Gamma_{>0}$. For every
$H=\sum_{\lambda\in S_H}t^\lambda h_\lambda$ with formal coefficient germs
$h_\lambda\in\C[[z]]$ --- in particular for every $H\in\A_n$ --- the fully
expanded family
\begin{equation}\label{contours:eq:radiusfreesub}
  H(W)=\sum_{\lambda\in S_H}\ \sum_{\alpha\in\N^n}
  t^\lambda\,[z^\alpha]h_\lambda\;W^\alpha
\end{equation}
is strongly summable, lies in $\HH(\mathsf X)$, and has support contained in
$S_H+P_W^*$. The assignment $H\mapsto H(W)$ respects sums and products and
commutes with differentiation in the ordinary parameters of $\mathsf X$. No
common convergence radius for the $h_\lambda$ is used, and none is produced.
\end{lemma}
\begin{proof}
A fully expanded monomial $t^\lambda[z^\alpha]h_\lambda W^\alpha$ contributes at
exponents $\lambda$ plus a finite word of $|\alpha|$ letters drawn from $P_W$,
one letter for each of the $|\alpha|$ factors. Fix a target exponent $\nu$. By
\cref{contours:lem:support} only finitely many pairs consisting of an element of
$S_H$ and a finite word in $P_W$ sum to $\nu$, and only finitely many word
lengths occur; hence only finitely many $\lambda$ and finitely many multi-indices
$\alpha$ can contribute, and for each there are finitely many ways to distribute
the letters of the word among the $n$ coordinates and their repeated factors. The
coefficient of $t^\nu$ is therefore a finite sum of finite products of ordinary
holomorphic functions on $\mathsf X$, so it is holomorphic on $\mathsf X$, and
the total support is contained in the well-ordered set $S_H+P_W^*$. The same
finiteness at each exponent licenses every regrouping, which gives the ring
identities, and ordinary differentiation of a coefficient function leaves its
exponent unchanged, so termwise differentiation is a finite operation at each
exponent. \emph{Convergence of the formal series $h_\lambda$ was never used};
what is used is that each individual output coefficient is a finite polynomial
expression in finitely many Taylor coefficients of finitely many $h_\lambda$ and
finitely many coefficient functions of $W$.
\end{proof}

\begin{corollary}[Radial evaluation of an ordinary family]
\label{contours:cor:radialeval}
Let $W(q,s)=\sum_{\ell\ge1}s^\ell W_\ell(q)$ be holomorphic on
$\mathsf X\times\DD_\epsilon$ with $W(q,0)=0$. For every $\rho\in\Gamma_{>0}$
the substitution $s=t^\rho$ produces an element of $\HH(\mathsf X)_+^n$ with
support in $\N_{>0}\rho$, and \cref{contours:lem:radiusfreesub} applies to it.
Each coefficient of the result is formed from finitely many ordinary Taylor
coefficients of the $W_\ell$, even when no single ordinary $\epsilon$ serves
all the germs $h_\lambda$ simultaneously.
\end{corollary}
\begin{proof}
The exponents of the substituted family are the positive integer multiples of
$\rho$, a well-ordered positive set, and $W(q,0)=0$ removes the exponent $0$.
Apply \cref{contours:lem:radiusfreesub} with the actual coordinate support
$P_W\subset\N_{>0}\rho$, whose generated monoid is contained in $\N\rho$.
At a fixed output exponent only finitely many radial degrees $\ell$ occur,
by \cref{contours:lem:support}.
\end{proof}

\emph{This is where the direction of the argument differs from
\cref{contours:thm:rescaling}.} Rescaling moves a radius-free germ into the
polynomial-coefficient ring by substituting a monomial for each coordinate.
\Cref{contours:lem:radiusfreesub} instead substitutes an entire Hahn family of
holomorphic maps on a fixed ordinary manifold, and lands in the coherent ring of
that manifold. Its construction includes diagonal rescaling by taking
$\mathsf X=\C^n$ and $W_j=t^{\rho_j}w_j$: finite incidence at each exponent
then gives polynomial coefficients. The injectivity in
\cref{contours:thm:rescaling} uses the additional fact that this diagonal
formal substitution is invertible; a general $W$ need not give an injective
map. Neither construction produces a common ordinary radius for the
original coefficient germs.

The second lemma solves the implicit equation that will produce the
equation-adapted lift of \cref{contours:sec:adapted}. Its coefficient ring is
not required to be a field, and its proof is combinatorial rather than
iterative.

\begin{lemma}[Positive implicit equation]
\label{contours:lem:positiveimplicit}
Let $R$ be a commutative complex algebra and let
$B(V)=(B_1(V),\ldots,B_n(V))$ be a formal power series in finitely many
variables $V=(V_1,\ldots,V_n)$ whose coefficients lie in $R((t^\Gamma))$, and
suppose the union of the supports of \emph{all} coefficients of $B$ is a
well-ordered subset $W\subset\Gamma_{>0}$. Then the equation
\begin{equation}\label{contours:eq:positiveimplicit}
  V=B(V)
\end{equation}
has exactly one solution with strictly positive coordinate supports. Its support
is contained in $W^*\setminus\{0\}$, it is given by a strongly summable family of
finite rooted trees, and each of its coefficients is computed by finitely many
coefficient operations in $R$.
\end{lemma}
\begin{proof}
Expand repeated substitution into finite ordered rooted trees. A vertex selects
one coefficient of one component of $B$; the multi-index of the selected monomial
prescribes how many children it has and which component each child computes.
Nullary vertices select constant terms. Assign to every vertex the exponent of
its selected coefficient, an element of $W$, and give a tree the sum of its
vertex exponents as its weight.

Fix a weight $\nu$. By \cref{contours:lem:support} only finitely many finite
ordered words in $W$ sum to $\nu$, and only finitely many lengths occur, so the
number of vertices of a contributing tree is bounded. There are finitely many
ordered rooted trees with at most that many vertices, finitely many ways to
colour their vertices by components, and finitely many compatible multi-indices,
since a vertex cannot have more children than there are other vertices. Hence the
tree family is strongly summable, and its sum has support in
$W^*\setminus\{0\}$; grafting all trees below a common root is a permitted
regrouping and proves \eqref{contours:eq:positiveimplicit}.

For uniqueness, suppose $V\ne V'$ both solve \eqref{contours:eq:positiveimplicit}
with positive coordinate supports, and let $\nu_0$ be the least exponent at which
some coordinate of $V$ and $V'$ differ. Every term of $B(V)-B(V')$ contains such a
difference multiplied by a coefficient of $B$, whose exponent is strictly
positive, and possibly by further coordinates of positive support. No such term
can contribute at $\nu_0$, so $V$ and $V'$ agree there, a contradiction.
\end{proof}

\begin{remark}[Not a Picard limit]
\label{contours:rem:notpicard}
\Cref{contours:lem:positiveimplicit} does not identify its support-defined sum
with a limit of iterates. That distinction is the one recorded in
\cref{contours:ex:rank}: the statement remains valid when the integer multiples
of the least coefficient valuation are \emph{not} cofinal in $\Gamma_{>0}$, where
no iteration argument based on a fixed improvement of valuation is available.
The lemma also needs no norm on $R$ and no completeness assumption; it needs only
positivity and well-ordering of the coefficient supports.
\end{remark}

\subsection{A fixed ramified family of ordinary Leray cycles}
\label{contours:subsec:ramified}
\emph{Everything in this subsection is ordinary complex analysis, and its radial
variable $s$ is an ordinary complex number.} The substitution
$s=t^\rho$ is made only in \cref{contours:subsec:protecting}. This order of
operations is the whole mechanism: it prevents the interchange of ordinary
analytic convergence with Hahn summability that
\cref{contours:subsec:rings,contours:ex:rank} warn against, and it is the reason
the construction never needs a common radius.

The classical input is the finite mapping theorem in its standard form
\cite[Chapters~II--III]{GrauertRemmert}. For $f$ as in
\eqref{contours:eq:untransformeddata} one may choose bounded representatives
$U_f$ of the source and $V_f$ of the target with $f:U_f\to V_f$ proper and finite,
$f^{-1}(0)=\{0\}$, and degree $\mu$, and a nonzero convergent germ $\Delta(w)$
vanishing on the branch locus, so that $f$ is an unramified covering of degree
$\mu$ over $V_f\setminus\{\Delta=0\}$. Concretely, take a small closed source ball
whose boundary avoids $f^{-1}(0)$ and then a target polydisk whose inverse image
stays clear of that boundary; properness follows, and the fibres are finite
because a positive-dimensional compact analytic subset of a domain in $\C^n$
cannot occur. If $f$ is already unramified near $0$ take $\Delta=1$. These are
ordinary finite-map facts, and none of them is a statement about a Hahn analytic
space.

\begin{lemma}[A discriminant-avoiding monomial direction in the \emph{target}]
\label{contours:lem:targetmonomial}
For every nonzero convergent germ $\Delta(w_1,\ldots,w_n)$ there are positive
integers $q_1,\ldots,q_n$, an integer $d\ge0$, a multi-index $\alpha$, and
$c\ne0$ such that on a fixed polyannulus
$\mathsf A=\{\varrho_1<|\xi_j|<\varrho_2\}$ around the ordinary torus
\begin{equation}\label{contours:eq:targetmonomial}
  \Delta(\tau^{q_1}\xi_1,\ldots,\tau^{q_n}\xi_n)
  =\tau^d\bigl(c\,\xi^\alpha+\tau R(\xi,\tau)\bigr),
\end{equation}
and, after shrinking the $\tau$ disk and if necessary the polyannulus, the
bracket is nowhere zero on a neighborhood of $\overline{\mathsf A}\times\{0\}$. Hence the
displayed target points avoid the branch locus for every sufficiently small
$\tau\ne0$.
\end{lemma}
\begin{proof}
Let $d_0$ be the least total degree of a nonzero term of $\Delta$. Choose
positive rational weights very close to $1$ and generic on the finitely many
monomials of total degree $d_0$, so that exactly one of them has least weighted
degree; taking the weights close enough to $1$ also makes every monomial of total
degree at least $d_0+1$ have strictly larger weighted degree than all
degree-$d_0$ monomials. Clearing denominators gives positive integers $q_j$ and a
unique least-weight monomial $c\,w^\alpha$ of weighted degree $d$. Substituting
in the convergent series gives \eqref{contours:eq:targetmonomial}, every other
weighted degree exceeding $d$ by a positive integer, and the remainder $R$ is
holomorphic for bounded $\xi$ and small $\tau$. On a closed polyannulus
$c\,\xi^\alpha$ is bounded away from zero, so compactness gives uniform
nonvanishing of the bracket for small $\tau$.
\end{proof}

The monomial direction is chosen in the \emph{target}, not in the source. This is
what allows an arbitrary isolated $f$: the lemma does \emph{not} require the
lowest weighted homogeneous parts of the $f_j$ to form an isolated complete
intersection, and it is exactly the step at which the objection recorded in
\cref{contours:subsec:scope} --- that a diagonal source rescaling need not keep
the leading system isolated or of coordinate-power shape --- is bypassed rather
than contradicted.

\begin{theorem}[Ramified Leray family; ordinary analysis]
\label{contours:thm:ramifiedleray}
For $f$ as in \eqref{contours:eq:untransformeddata} there exist a polyannulus
$\mathsf A$ around the ordinary torus $(S^1)^n$, a finite unramified holomorphic
covering $p:\mathsf X\to\mathsf A$ of total degree $\mu$ --- with $\mathsf X$ possibly
disconnected --- positive integers $\kappa_1,\ldots,\kappa_n$, and a
holomorphic map
\[
  Z:\mathsf X\times\DD_\epsilon\longrightarrow\C^n
\]
with
\begin{equation}\label{contours:eq:ordinaryZ}
  Z(q,0)=0,\qquad f_j\bigl(Z(q,s)\bigr)=s^{\kappa_j}p_j(q)
  \quad(1\le j\le n),
\end{equation}
where $p=(p_1,\ldots,p_n)$. For $s\ne0$ the fibre of $p$ over $\xi$ is
identified with the set of all $\mu$ inverse images in $U_f$ of
$(s^{\kappa_j}\xi_j)_j$. The restriction
$\mathsf T=p^{-1}\bigl((S^1)^n\bigr)$ is a compact oriented real
$n$-manifold without boundary, a finite disjoint union of tori, and
$Z(\mathsf T,s)$ is the ordinary Leray cycle of $f$ for target radii
$|s|^{\kappa_j}$. One may take $\kappa_j=Nq_j$ with $q_j$ as in
\cref{contours:lem:targetmonomial} and $N$ dividing
$\operatorname{lcm}(1,\ldots,\mu)$.
\end{theorem}
\begin{proof}
Pull the finite map $f:U_f\to V_f$ back along
\[
  A\times\DD^*\longrightarrow V_f,\qquad
  (\xi,\tau)\longmapsto(\tau^{q_j}\xi_j)_j .
\]
By \cref{contours:lem:targetmonomial} the image avoids the branch locus, so the
pullback is an unramified holomorphic covering of degree $\mu$, and its monodromy
is a permutation representation of
$\pi_1(A\times\DD^*)\simeq\Z^n\times\Z$. Let $\sigma$ be the permutation induced
by the radial generator and $N$ its order, so $N$ divides
$\operatorname{lcm}(1,\ldots,\mu)$ by the elementary order bound for a
permutation of $\mu$ letters. The substitution $\tau=s^N$ replaces $\sigma$ by
$\sigma^N=\id$, so the resulting covering has trivial monodromy along the radial
factor and is therefore pulled back from a finite covering $p:\mathsf X\to\mathsf A$;
the classification of coverings gives an isomorphism of the total space with
$\mathsf X\times\DD^*$, holomorphic because locally it identifies two
holomorphic local inverses of the same base map. \emph{No global labelling of
all $\mu$ sheets over $\mathsf A$ is required, and $\mathsf X$ may be disconnected.}

Composing with the source-coordinate projection gives a holomorphic
$Z:\mathsf X\times\DD^*\to U_f$ satisfying \eqref{contours:eq:ordinaryZ} for
$s\ne0$, with $\kappa_j=Nq_j$. It is bounded, because $U_f$ was chosen bounded, so
the removable-singularity theorem with holomorphic parameters extends it across
$s=0$: in a local chart, the Cauchy integral formula in $s$ kills every negative
Laurent coefficient, the remaining coefficients are holomorphic in $q$, and
extensions agree on overlaps. Properness of $f:U_f\to V_f$ keeps the inverse
images of a sufficiently small compact target neighborhood inside a compact
subset of $U_f$, so the extended values still lie in $U_f$. Continuity gives
$f(Z(q,0))=0$, whence $Z(q,0)=0$ because $f^{-1}(0)=\{0\}$.

For fixed $s\ne0$ every inverse image of the scaled target torus occurs exactly
once in this covering, and the source-coordinate map is injective on
$\mathsf T$: equality of source points forces equality of target points and then
equality of the sheet. Orient $\mathsf T$ so that
$p^*(d\arg\xi_1\wedge\cdots\wedge d\arg\xi_n)$ is positive; multiplying the
$j$th target coordinate by $s^{\kappa_j}$ only adds a constant phase, so the
induced orientation is the Leray orientation
$d\arg f_1\wedge\cdots\wedge d\arg f_n$ fixed in
\cref{contours:subsec:residuedef}. A finite covering of a compact torus is
compact with finitely many components, each a torus.
\end{proof}

\begin{lemma}[Coefficientwise constant-period identity]
\label{contours:lem:constantperiod}
For $h\in\C\{z\}$ and $k\in\N^n$ set
\begin{equation}\label{contours:eq:ordinaryform}
  \beta_{h,k}(s)=
  \frac{h\bigl(Z(q,s)\bigr)\,
    d_qZ_1(q,s)\wedge\cdots\wedge d_qZ_n(q,s)}
  {s^{\sum_j\kappa_j(k_j+1)}\,p_1(q)^{k_1+1}\cdots p_n(q)^{k_n+1}},
\end{equation}
where $d_q$ differentiates the angular parameters only. For each fixed $h$ and
$k$ this form has a convergent Laurent expansion in $s$ with finitely many
negative powers, on one ordinary neighborhood of $\mathsf T$ and for
sufficiently small $s$, and
\begin{equation}\label{contours:eq:constantperiod}
  \frac1{(2\pi i)^n}\int_{\mathsf T}\beta_{h,k}(s)
  =\res_{(f_1^{k_1+1},\ldots,f_n^{k_n+1})}(h).
\end{equation}
Equivalently, every nonconstant radial Laurent coefficient integrates to zero
and the constant one integrates to $(2\pi i)^n$ times the ordinary residue. The
admissible $s$ disk may depend on $h$ and on $k$; \emph{no simultaneous disk for
an infinite Hahn family of germs is asserted.}
\end{lemma}
\begin{proof}
Compactness of $\mathsf T$ and $Z(q,0)=0$ give an ordinary $s$ disk on which
$h(Z(q,s))$ is holomorphic on one neighborhood of $\mathsf T$. The remaining
numerator factors are holomorphic there, and the denominator contributes only
the displayed finite negative power of $s$, since the $p_j$ are nowhere zero on
$\mathsf A$. Ordinary compact-uniform convergence justifies integrating that single
Laurent expansion term by term. For small positive real $s$ the cycle
$Z(\mathsf T,s)$ is an ordinary regular Leray cycle for $f$ with target radii
$|s|^{\kappa_j}$, so its normalized integral is the ordinary local Grothendieck
residue of $h\dd z/\prod_jf_j^{k_j+1}$ \cite{GriffithsHarris,TajimaNabeshima}.
The period has this same value for every sufficiently small positive real $s$.
The identity theorem makes it constant on the punctured $s$ disk, and uniqueness
of the Laurent expansion gives the coefficient statement.
\end{proof}

\subsection{The universal protecting cycle}
\label{contours:subsec:protecting}
Now, and only now, the radial variable becomes a Hahn monomial. Set
\begin{equation}\label{contours:eq:kappastar}
  \kappa_*=\max_j\kappa_j,\qquad \kappa_\Sigma=\sum_j\kappa_j,
\end{equation}
and choose $\rho\in\Gamma_{>0}$ with
\begin{equation}\label{contours:eq:protectingrho}
  \kappa_*\rho<\delta ;
\end{equation}
divisibility of $\Gamma$ supplies one, for example $\rho=\delta/(\kappa_*+1)$.
The direction of \eqref{contours:eq:protectingrho} is the one already recorded
in \cref{contours:lem:protect}: a \emph{smaller} positive exponent is a
\emph{larger} infinitesimal, so the target radii must be large in modulus
relative to the perturbation, with strictly smaller valuations. Expanding the ordinary family of
\cref{contours:thm:ramifiedleray} as $Z(q,s)=\sum_{\ell\ge1}s^\ell Z_\ell(q)$,
put
\begin{equation}\label{contours:eq:Zzero}
  Z_0=\sum_{\ell\ge1}t^{\ell\rho}Z_\ell
  \ \in\ \HH(\mathsf X)_+^n ,
  \qquad\text{so}\qquad
  f_j(Z_0)=t^{\kappa_j\rho}p_j .
\end{equation}
This is legitimate by \cref{contours:cor:radialeval}, all coordinate supports lie
in $\N_{>0}\rho$, and $(\mathsf T,Z_0)$ is an admissible Hahn cycle in the sense
of \cref{contours:def:hahncycle}. \emph{It depends only on $f$ and $\delta$.}

For the particular $E$ under consideration put
\begin{equation}\label{contours:eq:Pzero}
  P_0=\{\rho\}\cup\bigcup_{j=1}^n\bigl(S_E-\kappa_j\rho\bigr),
  \qquad \mathcal M_0=P_0^* .
\end{equation}
Finite unions and translates of well-ordered sets are well ordered, and
\eqref{contours:eq:protectingrho} makes every element of $P_0$ strictly
positive.

\begin{proposition}[Admissibility of the untransformed pullback]
\label{contours:prop:admissiblepullback}
For every $H\in\A_n$ the pullback by $Z_0$ of the untransformed form
$H\dd z_1\wedge\cdots\wedge \dd z_n/(F_1\cdots F_n)$ is a Hahn family of
holomorphic $n$-forms on $\mathsf X$ with support contained in
\begin{equation}\label{contours:eq:untransformedsupport}
  S_H-\kappa_\Sigma\rho+\mathcal M_0 ,
\end{equation}
each coefficient being a finite sum of ordinary holomorphic $n$-forms on
$\mathsf X$. The assertion holds for the \emph{fully expanded} family, before any
Taylor, perturbation, or radial index has been regrouped.
\end{proposition}
\begin{proof}
By \cref{contours:lem:radiusfreesub}, $H(Z_0)\in\HH(\mathsf X)$ with support in
$S_H+\N\rho$, and likewise $\supp E_j(Z_0)\subset S_E+\N\rho$. Put
\[
  Q_j=\frac{E_j(Z_0)}{t^{\kappa_j\rho}p_j},
\]
whose support lies in $S_E-\kappa_j\rho+\N\rho$, a positive subset of
$\mathcal M_0$; the functions $p_j^{-1}$ are ordinary holomorphic units on the
same $\mathsf X$. Hence
\begin{equation}\label{contours:eq:strongdenominator}
  \frac1{F_j(Z_0)}
  =t^{-\kappa_j\rho}p_j^{-1}\sum_{r\ge0}(-Q_j)^r
\end{equation}
is a strongly summable identity in $\HH(\mathsf X)$, by the Neumann inversion
recorded after \cref{contours:lem:support}. Each angular derivative of $Z_0$
still has support in $\N_{>0}\rho$, so multiplying the numerator and the
differential factors by \eqref{contours:eq:strongdenominator} gives
\eqref{contours:eq:untransformedsupport}.

For the stronger indexed assertion, expand a contribution carrying $k_j$ factors
from $E_j$. Its exponent has the form
\begin{equation}\label{contours:eq:exponentaccounting}
  \lambda-\kappa_\Sigma\rho
  +\sum_{j=1}^n\sum_{a=1}^{k_j}\bigl(e_{j,a}-\kappa_j\rho\bigr)+r\rho,
  \qquad \lambda\in S_H,\ e_{j,a}\in S_E,\ r\ge0,
\end{equation}
and every letter after the fixed shift $-\kappa_\Sigma\rho$ belongs to $P_0$. At
an exact total exponent, \cref{contours:lem:support} permits only finitely many
words, hence finitely many multi-indices $k$, finitely many perturbation
coefficient germs and finitely many radial degrees; and for a fixed radial
degree, substitution into $Z(q,s)$ reads only finitely many Taylor monomials,
because $Z(q,0)=0$. This is local finiteness before regrouping, which is what
the integration step requires.
\end{proof}

\begin{theorem}[Universal untransformed contour]
\label{contours:thm:untransformed}
Let $f$ and $\delta$ be as in
\eqref{contours:eq:untransformeddata}--\eqref{contours:eq:threshold}, and let
$\mathsf T$, $p$, $\kappa_j$ and $Z_0$ be as in
\cref{contours:thm:ramifiedleray} and \eqref{contours:eq:Zzero}. Then for
\emph{every} $E\in\A_n^n$ with $\supp E\subset[\delta,\infty)$ and \emph{every}
$H\in\A_n$, the untransformed form has a well-defined coefficientwise pullback
and
\begin{equation}\label{contours:eq:untransformed}
 \boxed{\qquad
 \Res_F(H)=\frac1{(2\pi i)^n}
 \HInt_{Z_0[\mathsf T]}
 \frac{H(z)\dd z_1\wedge\cdots\wedge \dd z_n}{F_1(z)\cdots F_n(z)} ,
 \qquad}
\end{equation}
with $\Res_F$ the perturbation residue \eqref{contours:eq:generalres}. The cycle
$(\mathsf T,Z_0)$ is constructed from $f$ and $\delta$ alone, independently of
$E$ and of $H$; no determinant transformation of $F$ is made and no equation is
replaced; and all resulting coefficient functions are holomorphic on one common
ordinary neighborhood of $\mathsf T$, with the explicit support certificate
\eqref{contours:eq:untransformedsupport}.
\end{theorem}
\begin{proof}
Expand the reciprocal denominators by
\eqref{contours:eq:strongdenominator}. The proposed normalized integral becomes
\begin{equation}\label{contours:eq:expandedintegral}
  \frac1{(2\pi i)^n}\sum_{k\in\N^n}(-1)^{|k|}
  \HInt_{\mathsf T}
  \frac{\bigl(H\,E_1^{k_1}\cdots E_n^{k_n}\bigr)(Z_0)\;
    dZ_{0,1}\wedge\cdots\wedge dZ_{0,n}}
  {\prod_{j=1}^nf_j(Z_0)^{k_j+1}} ,
\end{equation}
and the whole expanded family is strongly summable by
\cref{contours:prop:admissiblepullback}, so integration and every regrouping are
legitimate --- the coefficientwise integral \eqref{contours:eq:chainintegral}
being a complex-linear operation applied at each exponent, with no interchange of
an infinite limit.

Fix one selection of Hahn coefficients from $H$ and from the finitely many
perturbation factors occurring in one term of index $k$. Their ordinary product
is a single convergent germ $h$, and \emph{before} the substitution $s=t^\rho$
the angular pullback of the corresponding summand is exactly the ordinary form
\eqref{contours:eq:ordinaryform}. By \cref{contours:lem:constantperiod} its
integral kills every nonconstant radial Laurent coefficient and returns
$\res_{(f_1^{k_1+1},\ldots,f_n^{k_n+1})}(h)$. Multiplying by the selected Hahn
monomial and summing the strongly summable family reproduces
\eqref{contours:eq:generalres} term for term, including the sign $(-1)^{|k|}$
and the normalization $(2\pi i)^{-n}$ fixed in
\cref{contours:subsec:residuedef}.

Only $f$ and $\delta$ entered the construction of $\mathsf T$, $p$ and $Z_0$.
The support certificate \eqref{contours:eq:Pzero} depends on $S_E$, but the cycle
does not, and $H$ was arbitrary throughout.
\end{proof}

\begin{corollary}[Cancellation of the auxiliary scales]
\label{contours:cor:scalecancel}
Although the pulled-back form carries exponents in the larger set
\eqref{contours:eq:untransformedsupport}, the integral
\eqref{contours:eq:untransformed} has support contained in $\supp H+(\supp E)^*$
and is independent of the chosen discriminant $\Delta$, the target weights
$q_j$, the radial ramification order $N$, the identification of the covering,
and the protecting exponent $\rho$.
\end{corollary}
\begin{proof}
The value equals the independently defined series
\eqref{contours:eq:generalres}, whose support bound was established in
\cref{contours:subsec:residuedef} and whose definition mentions none of those
choices.
\end{proof}

\emph{This is independence of the period only.} It does not assert that two
differently scaled or differently ramified cycles are joined by a pole-free
homotopy in some larger chain category; the only homotopy proved below is the
one of \cref{contours:thm:untransformedhomotopy}, between the protecting cycle
and its own adapted lift at a fixed scale. A chain-level comparison across
weights, covers and scales remains open; see
\cref{contours:subsec:untransformedscope}.

\subsection{Why the missing common radius causes no gap}
\label{contours:subsec:tworadii}
Two different finiteness principles are used in the proof of
\cref{contours:thm:untransformed}, and the construction works precisely because
they are kept apart.

\begin{enumerate}
\item \textbf{For a fixed ordinary coefficient product}, ordinary compactness of
$\mathsf T$ supplies a small ordinary $s$ disk and proves the Laurent period
identity \eqref{contours:eq:constantperiod}. This is an ordinary analytic
statement about one germ; the disk may shrink as the germ changes.
\item \textbf{For a fixed Hahn exponent}, \cref{contours:lem:support} supplies a
finite list of contributing coefficient products, multi-indices and radial
degrees. This is an algebraic statement about supports; it says nothing about
any disk.
\end{enumerate}

No step asks for one ordinary disk that serves all coefficient germs
simultaneously, which is what \eqref{contours:eq:shrinking} shows to be
unavailable. The common parameter domain of the cycle comes from the finite map
$f$ and its covering --- objects built before any perturbation is named --- and
not from an intersection of infinitely many convergence neighborhoods. This is
also why the conclusion is stronger than a representation assembled separately
for each numerator or each perturbation: one cycle serves an entire valuation
ball of perturbations whose coefficient germs may have radii tending to zero.
It is the same distinction as the one recorded after \cref{contours:ex:rank} and
used in \cref{contours:lem:residuetransform}, applied here to a contour rather
than to a transformation law.

\section{The equation-adapted lift, its normalizations, and its limits}
\label{contours:sec:adapted}

The protecting cycle of \cref{contours:thm:untransformed} is adapted to the
\emph{leading} equations: it satisfies $f_j(Z_0)=t^{\kappa_j\rho}p_j$ exactly,
while $F_j(Z_0)$ differs from that by a positive-support error. This section
corrects it, on the same finite cover, to a cycle on which the \emph{perturbed}
equations hold exactly, connects the two by a pole-free Hahn homotopy, and
derives the normalizations that \cref{contours:subsec:problems} named as the
acceptance tests for any proposed cycle. It closes with the example, the
pathologies, the construction guide, and the exact boundaries of the result.

\subsection{Two uniform ordinary radial bounds}
\label{contours:subsec:radialbounds}
The following bounds are ordinary, depend only on the finite map and its ramified
family of \cref{contours:thm:ramifiedleray}, and are fixed before any Hahn
coefficient germ is named.

\begin{lemma}[Finite radial bounds]
\label{contours:lem:radialbounds}
After shrinking the polyannulus and the ordinary $s$ disk around the same compact
$\mathsf T$, there are integers $b_{\mathrm{pole}},c_{\mathrm{sep}}\ge0$ such
that
\begin{equation}\label{contours:eq:Jbound}
  \bigl(D_zf(Z(q,s))\bigr)^{-1}=s^{-b_{\mathrm{pole}}}\,\mathsf J(q,s)
\end{equation}
with $\mathsf J$ holomorphic on the fixed product neighborhood, and such that
for any distinct $q,q'\in\mathsf T$ with $p(q)=p(q')$ at least one coordinate of
$Z(q,s)-Z(q',s)$ has a nonzero Taylor coefficient in $s$ of degree at most
$c_{\mathrm{sep}}$. If $\mu=1$ one may take $c_{\mathrm{sep}}=0$.
\end{lemma}
\begin{proof}
The determinant $a(q,s)=\det D_zf(Z(q,s))$ is holomorphic across $s=0$ and
nowhere zero for $s\ne0$, because $Z(\mathsf T,s)$ then lies over regular target
values. Near a point $(q_0,0)$, Weierstrass preparation in $s$ writes $a$ as a
unit times a monic polynomial in $s$; every root of that polynomial must be $0$,
since no zero with $s\ne0$ occurs in the product neighborhood, so locally
$a=s^ru$ with $u$ a holomorphic unit and $r$ locally constant along
$\mathsf X\times\{0\}$. A finite cover of the compact $\mathsf T$ by such
neighborhoods, together with the adjugate formula for the inverse, gives a
common finite pole order, and shrinking to a common collar of $\mathsf T$ and a
common $s$ disk gives \eqref{contours:eq:Jbound}.

For the separation bound, consider the off-diagonal fibre product
$\mathsf T\times_{(S^1)^n}\mathsf T\setminus\operatorname{diag}(\mathsf T)$,
which is compact, being a union of components of a finite cover. At each of its
points the two holomorphic source-coordinate vectors differ for $s\ne0$, so some
Taylor coefficient of their difference is nonzero, and nonvanishing of that
coefficient persists on an ordinary neighborhood. A finite subcover gives a
uniform upper bound $c_{\mathrm{sep}}$ on a degree at which a difference is
already detected.
\end{proof}

Fix now
\begin{equation}\label{contours:eq:Lrho}
  L=\max\{b_{\mathrm{pole}},c_{\mathrm{sep}},\kappa_*\}+1,
  \qquad
  \rho\in\Gamma_{>0}\ \text{with}\ (L+b_{\mathrm{pole}})\rho<\delta ,
\end{equation}
for example $\rho=\delta/\bigl[2(L+b_{\mathrm{pole}}+1)\bigr]$, available by
divisibility. Rebuild $Z_0$ by \eqref{contours:eq:Zzero} with this possibly
smaller exponent; since $L+b_{\mathrm{pole}}\ge\kappa_*$, the protecting
inequality \eqref{contours:eq:protectingrho} still holds and
\cref{contours:thm:untransformed} is unaffected. Put
\begin{equation}\label{contours:eq:Pfull}
  P=\{\rho\}\cup\bigl(S_E-(L+b_{\mathrm{pole}})\rho\bigr),
  \qquad \mathcal M=P^* ,
\end{equation}
a well-ordered positive generating set with $P_0\subset\mathcal M$, because
$L+b_{\mathrm{pole}}\ge\kappa_j$ for every $j$.

\subsection{The scaled implicit equation and the exact lift}
\label{contours:subsec:exactlift}
Write $s=t^\rho$ from here on as an algebraic abbreviation only --- the ordinary
analysis is finished --- and set
\[
  J_0=D_zf(Z_0),\qquad \Phi=Z_0+t^{L\rho}V .
\]
The requirement $F(\Phi)=f(Z_0)$ is equivalent to the fixed-point equation
\begin{align}
  V&=B(V),\label{contours:eq:Bequation}\\
  B(V)&=-t^{-L\rho}J_0^{-1}\Bigl(
    f\bigl(Z_0+t^{L\rho}V\bigr)-f(Z_0)-t^{L\rho}J_0V
    +E\bigl(Z_0+t^{L\rho}V\bigr)\Bigr),
  \label{contours:eq:Bdefinition}
\end{align}
a formal power series in $V$ with coefficients in $\HH(\mathsf X)$.

\begin{theorem}[Equation-adapted lift on the same cover]
\label{contours:thm:adaptedlift}
Equation \eqref{contours:eq:Bequation} has exactly one solution with strictly
positive coordinate supports. The correction $\Upsilon=t^{L\rho}V$ satisfies
\begin{equation}\label{contours:eq:Upsilonsupport}
  \supp\Upsilon\subset L\rho+\bigl(\mathcal M\setminus\{0\}\bigr),
  \qquad
  \supp\Upsilon\subset S_E-b_{\mathrm{pole}}\rho+\mathcal M ,
\end{equation}
so $\vH(\Upsilon)>L\rho$ and $\vH(\Upsilon)\ge\delta-b_{\mathrm{pole}}\rho$ when
$\Upsilon\ne0$, and $\Upsilon=0$ when $E=0$. The parameterization
$\Phi=Z_0+\Upsilon$ lies in $\HH(\mathsf X)_+^n$ on the \emph{same} fixed
ordinary parameter domain, is an admissible Hahn cycle on the same finite cover
$\mathsf T$ in the sense of \cref{contours:def:hahncycle}, and satisfies the
perturbed target equations exactly:
\begin{equation}\label{contours:eq:adapted}
  F_j\bigl(\Phi(q)\bigr)=t^{\kappa_j\rho}p_j(q)
  \qquad (q\in\mathsf T,\ 1\le j\le n).
\end{equation}
It is the unique parameterization of the form $Z_0+\Upsilon$ with
$\vH(\Upsilon)>L\rho$ satisfying \eqref{contours:eq:adapted}. No shrinking of the
parameter domain depending on $E$ occurs.
\end{theorem}
\begin{proof}
Consider the coefficients of \eqref{contours:eq:Bdefinition}. A coefficient of
total $V$-degree $r\ge2$ coming from the ordinary nonlinear part has exponents in
$\bigl((r-1)L-b_{\mathrm{pole}}\bigr)\rho+\N\rho$, positive and contained in
$\mathcal M$ because $L>b_{\mathrm{pole}}$. A coefficient of total degree
$r\ge0$ coming from the perturbation has exponents in
$S_E-(L+b_{\mathrm{pole}})\rho+rL\rho+\N\rho$, again a positive subset of
$\mathcal M$. Both assertions follow from Taylor expansion together with
\eqref{contours:eq:Jbound}, and all coefficient functions are holomorphic on the
same $\mathsf X$ by \cref{contours:lem:radiusfreesub}. Hence the union of the
supports of all coefficients of $B$ is a well-ordered subset of
$\mathcal M\setminus\{0\}$, and \cref{contours:lem:positiveimplicit} applies with
$R=\OO(\mathsf X)$, giving existence, uniqueness in the positive ball, and
$\supp V\subset\mathcal M\setminus\{0\}$. The first bound in
\eqref{contours:eq:Upsilonsupport} is immediate.

For the second bound, every finite rooted tree in the expansion of
\cref{contours:lem:positiveimplicit} has at least one nullary vertex, and the
constant coefficient of $B$ is exactly $-t^{-L\rho}J_0^{-1}E(Z_0)$, with support
in $S_E-(L+b_{\mathrm{pole}})\rho+\N\rho$. Removing one such vertex exponent from
a tree's weight leaves exponents in $\mathcal M$, so
$\supp V\subset S_E-(L+b_{\mathrm{pole}})\rho+\mathcal M$; multiplying by
$t^{L\rho}$ gives the second bound. If $E=0$ that constant coefficient vanishes,
every tree has weight belonging to no admissible word, and $V=0$.

All substitutions used to pass between $F(\Phi)=f(Z_0)$ and
\eqref{contours:eq:Bequation} are covered by the same support bounds and
\cref{contours:lem:radiusfreesub}, so reversing the algebra gives
$F(\Phi)=f(Z_0)=\bigl(t^{\kappa_j\rho}p_j\bigr)_j$ \emph{exactly}, not modulo a
truncation. Uniqueness among corrections with $\vH(\Upsilon)>L\rho$ is
uniqueness in \cref{contours:lem:positiveimplicit}.
\end{proof}

\begin{corollary}[Regular protected sheets]
\label{contours:cor:protectedsheets}
The matrix $D_zF(\Phi)$ is invertible in the Hahn matrix algebra over
$\HH(\mathsf X)$. For each fixed ordinary $\xi\in(S^1)^n$ the $\mu$ values
$\Phi(q)$ with $p(q)=\xi$ are pairwise distinct. In local holomorphic
coordinates on $\mathsf X$ the angular derivative of $\Phi$ has full rank over
the Hahn coefficient ring.
\end{corollary}
\begin{proof}
Write $J_0^{-1}D_zF(\Phi)=I+J_0^{-1}\bigl(D_zF(\Phi)-J_0\bigr)$. The ordinary
part of the difference contains at least one factor $\Upsilon$, so after
multiplication by $J_0^{-1}$ its exponents exceed
$(L-b_{\mathrm{pole}})\rho>0$; the perturbation derivative contributes exponents
at least $\delta-b_{\mathrm{pole}}\rho>0$. A matrix Neumann series inverts the
right-hand side, and $J_0$ is invertible by \eqref{contours:eq:Jbound}. For two
distinct sheets, \cref{contours:lem:radialbounds} detects a nonzero coefficient
of $Z_0(q)-Z_0(q')$ at an exponent at most $c_{\mathrm{sep}}\rho$, while the
difference of the corrections has all exponents strictly above
$L\rho>c_{\mathrm{sep}}\rho$ and cannot cancel it. Finally, differentiating
\eqref{contours:eq:adapted} gives
\[
  D_zF(\Phi)\,d\Phi
  =\operatorname{diag}\bigl(t^{\kappa_1\rho},\ldots,t^{\kappa_n\rho}\bigr)\,dp,
\]
whose right-hand side has invertible determinant in local coordinates, since $p$
is an unramified covering and the monomials are units in $K$.
\end{proof}

\emph{These are statements about the protected sheets only.} They do not count
the zeros of $F$ in any larger surreal domain, they do not assert that every
zero of $F$ is simple, and they do not assume that the displaced zero scheme is
reduced; singular leading zeros are allowed throughout, and the cycle lies over
regular nonzero target values. In particular \eqref{contours:eq:adapted} is an
exact identity about \emph{one parameterized cycle}. \emph{The image
$\Phi[\mathsf T]$ is not identified with the entire class of surcomplex
solutions of $|F_j|=t^{\kappa_j\rho}$, and no such identification is needed for
any integral below.} The distinction is the same one drawn in
\cref{contours:ex:definablecircle} between a full semialgebraic circle and an
ordinary-parameter contour trace.

\subsection{A pole-free homotopy between the two cycles}
\label{contours:subsec:polefree}
Let $\Phi_u=Z_0+u\Upsilon$ for $0\le u\le1$, with $u$ an \emph{ordinary} real
homotopy parameter as in \eqref{contours:eq:homotopyoperator}; it does not index
a limit in $K$.

\begin{lemma}[Admissible straight homotopy]
\label{contours:lem:straighthomotopy}
For each $j$ and $0\le u\le1$,
\[
  F_j(\Phi_u)=t^{\kappa_j\rho}p_j\,(1+Q_{j,u}),
\]
where $Q_{j,u}$ has positive support contained in $\mathcal M$ and coefficients
smooth on $\mathsf X\times[0,1]$. Every Taylor substitution, geometric
inversion, differentiation and pullback for the untransformed form is strongly
summable there, and the total pullback has support contained in
$S_H-\kappa_\Sigma\rho+\mathcal M$.
\end{lemma}
\begin{proof}
The difference $f_j(Z_0+u\Upsilon)-f_j(Z_0)$ carries at least one factor
$\Upsilon$, so its support lies in $L\rho+(\mathcal M\setminus\{0\})$, and
division by $t^{\kappa_j\rho}p_j$ leaves positive support because
$L>\kappa_j$. The perturbation part has support in $S_E+\mathcal M$, and division
leaves positive support because $S_E-\kappa_j\rho\subset\mathcal M\setminus\{0\}$
by \eqref{contours:eq:protectingrho}. Expand $(1+Q_{j,u})^{-1}$ geometrically;
\cref{contours:lem:support} controls the whole indexed family, exactly as in
\cref{contours:prop:admissiblepullback}. At each exact exponent only finitely
many Taylor and geometric terms contribute, so the dependence on $u$ is
polynomial there, and ordinary differentiation in $u$ or in the angular
parameters cannot enlarge support. This covers the total pullback including its
$\dd u$ terms.
\end{proof}

\begin{theorem}[Untransformed homotopy invariance]
\label{contours:thm:untransformedhomotopy}
For every $H\in\A_n$,
\begin{equation}\label{contours:eq:untransformedhomotopy}
  \HInt_{\Phi[\mathsf T]}
  \frac{H\dd z_1\wedge\cdots\wedge \dd z_n}{F_1\cdots F_n}
  =\HInt_{Z_0[\mathsf T]}
  \frac{H\dd z_1\wedge\cdots\wedge \dd z_n}{F_1\cdots F_n} .
\end{equation}
Consequently \eqref{contours:eq:untransformed} holds with $\Phi[\mathsf T]$ in
place of $Z_0[\mathsf T]$: the equation-adapted cycle of
\cref{contours:thm:adaptedlift} also represents $\Res_F$.
\end{theorem}
\begin{proof}
A meromorphic top-degree form in $n$ complex variables is closed wherever its
denominators are invertible, since differentiating its coefficient produces an
extra $\dd z_j$ which wedges to zero against
$\dd z_1\wedge\cdots\wedge \dd z_n$; this is a differential-algebra identity and
needs no analytic hypothesis. The support-justified chain rule of
\cref{contours:lem:radiusfreesub} transports it to the total pullback under
$\Phi_u$. By \cref{contours:lem:straighthomotopy} that pullback is a
strongly summable family of smooth closed $n$-forms on the compact ordinary
cylinder $\mathsf T\times[0,1]$, so \cref{contours:thm:hahnstokes} applies to
each coefficient. Since $\partial\mathsf T=\varnothing$, the only boundary terms
are the two end cycles with opposite orientations, and their difference
vanishes. Combining with \cref{contours:thm:untransformed} gives the last
assertion.
\end{proof}

The ordinary pullback Jacobian $d\Phi_1\wedge\cdots\wedge d\Phi_n$ is of course
present in \eqref{contours:eq:untransformedhomotopy}; it is the
change-of-parameters factor that every pullback of a differential form carries,
and it is \emph{not} an auxiliary determinant arising from replacing the equation
system $F$ by $AF$. Neither the integrand nor its list of denominators has been
altered, which is the precise difference from
\cref{contours:thm:generalrepresentation} and
\cref{contours:thm:septorus}, where $\det A$ and $\det\mathsf C$ are not
optional.

\subsection{Degree normalization, a finite-sheet trace, and the two acceptance
tests}
\label{contours:subsec:normalization}
Write $J_F=\det(\partial F_i/\partial z_j)$ as in
\eqref{contours:eq:traceresidue}.

\begin{theorem}[Degree normalization on the adapted cycle]
\label{contours:thm:degreenormalization}
Under the hypotheses of \cref{contours:thm:untransformed},
\begin{equation}\label{contours:eq:jacobiannorm}
  \Res_F(J_F)=\mu ,
\end{equation}
and more generally, for $G\in\A_n$ in the target variables $w$,
\begin{equation}\label{contours:eq:targetnorm}
  \Res_F\bigl(G(F)\,J_F\bigr)=\mu\,G(0),
\end{equation}
where $G(F)$ is expanded about the ordinary map $f$ in the positive perturbation
$E$ and is therefore a well-defined radius-free Hahn germ.
\end{theorem}
\begin{proof}
Work on the adapted cycle, using
\cref{contours:thm:adaptedlift,contours:thm:untransformedhomotopy}. By
\eqref{contours:eq:adapted} and \cref{contours:cor:protectedsheets},
\[
  \Phi^*\Bigl(\frac{J_F\dd z_1\wedge\cdots\wedge \dd z_n}{F_1\cdots F_n}\Bigr)
  =\bigwedge_{j=1}^n\frac{d(F_j\circ\Phi)}{F_j\circ\Phi}
  =p^*\Bigl(\frac{\dd w_1}{w_1}\wedge\cdots\wedge\frac{\dd w_n}{w_n}\Bigr),
\]
whose normalized integral over the degree-$\mu$ covering $\mathsf T$ is $\mu$.
This proves \eqref{contours:eq:jacobiannorm} \emph{without} importing a Hahn
flatness theorem or a Nullstellensatz. For the second identity, write
$G=\sum_{\lambda\in S_G}t^\lambda g_\lambda$. The required composition is
\[
  G(F)=\sum_{\lambda,\alpha}
  t^\lambda\frac{(\partial^\alpha g_\lambda)\circ f}{\alpha!}\,E^\alpha.
\]
Since $f(0)=0$, every ordinary coefficient composition is a convergent germ.
\Cref{contours:lem:support} gives containing support $S_G+S_E^*$ and finitely
many source exponents and Taylor terms at every output exponent. Thus $G(F)$
belongs to $\A_n$, even though $F$ itself has an exponent-zero part and is not
a positive-support input to \cref{contours:lem:radiusfreesub}. On the adapted cycle,
that lemma and the finite Taylor composition identity give
$G(F\circ\Phi)=G\bigl(t^{\kappa_1\rho}p_1,\ldots,t^{\kappa_n\rho}p_n\bigr)$,
whose nonconstant target Taylor monomials integrate to zero against
$\bigwedge_j\dd p_j/p_j$ on the torus, while its constant monomial is $G(0)$.
Finite covering degree and strong summability give
\eqref{contours:eq:targetnorm}.
\end{proof}

The ordinary degree-residue identity behind \eqref{contours:eq:jacobiannorm} is
classical. What is added is that it survives for entirely unrestricted
radius-free positive-support perturbations and follows from an actual admissible
parameterized cycle rather than from a finite-algebra package. \emph{Calling
\eqref{contours:eq:jacobiannorm} a conservation law for local algebra length
would additionally require the relevant finite-flatness identification; that
identification is not smuggled into this proof.}

\begin{corollary}[Protected trace representation]
\label{contours:cor:protectedtrace}
For $H\in\A_n$ the fibrewise sum
\[
  \Theta_H(\xi)=\sum_{q\in p^{-1}(\xi)}
  \frac{H(\Phi(q))}{J_F(\Phi(q))}
\]
is a single-valued Hahn-holomorphic function on the polyannulus $\mathsf A$, and
\begin{equation}\label{contours:eq:finitetrace}
  \Res_F(H)=\frac1{(2\pi i)^n}
  \HInt_{(S^1)^n}\Theta_H(\xi)\,
  \frac{\dd\xi_1}{\xi_1}\wedge\cdots\wedge\frac{\dd\xi_n}{\xi_n} .
\end{equation}
\end{corollary}
\begin{proof}
Invertibility of $D_zF(\Phi)$ and the differential of
\eqref{contours:eq:adapted} give
\[
  \Phi^*\Bigl(\frac{H\dd z_1\wedge\cdots\wedge \dd z_n}{\prod_jF_j}\Bigr)
  =\frac{H(\Phi)}{J_F(\Phi)}\bigwedge_j\frac{\dd p_j}{p_j},
\]
and integration over a finite covering is integration of the fibrewise sum,
valid for each ordinary coefficient and hence for the Hahn family. The matrix
inverse used here is a fixed finite monomial shift times a positive-support
expansion, so it adds no summability question.
\end{proof}

This is a trace over \emph{protected regular fibres}. It is not a claim that all
individual zeros of $F$ at target value $0$ have been separated or
parameterized, and it does not identify $\Theta_H$ with a multiplication trace on
a finite algebra.

\begin{proposition}[Annihilation of the equation ideal, and descent]
\label{contours:prop:idealannihilation}
For every $H\in\A_n$ and every $j$, $\Res_F(F_jH)=0$. Hence $\Res_F$ descends to
a nonzero $K$-linear functional on $\A_n/(F_1,\ldots,F_n)$.
\end{proposition}
\begin{proof}
Classical residue annihilation gives
$\res_{(f_1^{k_1+1},\ldots,f_n^{k_n+1})}(f_jh)=0$ when $k_j=0$, and equals
$\res_{(\ldots,f_j^{k_j},\ldots)}(h)$ --- the same functional with $k_j$ lowered
by one --- when $k_j>0$. Substituting these identities into
\eqref{contours:eq:generalres} and shifting the $j$th summation index gives
$\Res_F(f_jH)=-\Res_F(E_jH)$, every reindexing being licensed by positive
support and \cref{contours:lem:support}; adding $\Res_F(E_jH)$ gives
$\Res_F(F_jH)=0$. The functional is nonzero by
\eqref{contours:eq:jacobiannorm} in characteristic zero. \emph{This is descent
only, and not nondegeneracy of a duality pairing.}
\end{proof}

\paragraph{The two acceptance tests, checked.}
\Cref{contours:subsec:problems} fixed two compatibility tests for any proposed
cycle: its integral should annihilate $(F)$, and it should satisfy
$\Res_F(HJ_F)=\Tr(M_H)$. The first is now proved for the cycle itself, by
\cref{contours:prop:idealannihilation}, from classical residue annihilation and
nothing else. The second is a statement about the functional $\Res_F$, and it is
the reported companion result \eqref{contours:eq:traceresidue}; what the new
theorem contributes is that the untransformed cycle integral \emph{is} that
functional, so the test is met exactly to the extent that
\eqref{contours:eq:traceresidue} holds --- and its $H=1$ case,
$\Res_F(J_F)=\mu=\Tr(M_1)$, is now proved directly on the cycle by
\cref{contours:thm:degreenormalization}, without the finite-algebra package.
\emph{Neither test is offered as a substitute for the construction, and passing
them was never in doubt for a functional defined by
\eqref{contours:eq:generalres}; their role is to certify that the cycle computes
the right functional and not some near neighbour of it.}

\subsection{Formal coefficient germs}
\label{contours:subsec:formalcoeff}
The radius-free theorem needs no common radius but still assumes each coefficient
germ converges. That hypothesis can be removed from $E$ and $H$, by a separate
argument. The leading map $f$ remains an ordinary \emph{convergent} map with an
isolated zero throughout; its convergence is what produced the bounded finite
covering, and no formal substitute for that object is offered here.

\begin{lemma}[Finite-jet extension of an ordinary local residue]
\label{contours:lem:finitejet}
For fixed $f$ and $k$ the functional
$\res_{(f_1^{k_1+1},\ldots,f_n^{k_n+1})}$ depends only on a finite Taylor jet of
its numerator, and has a unique extension to $\C[[z]]$ continuous for the formal
maximal-ideal filtration.
\end{lemma}
\begin{proof}
The ideal $I_k=(f_1^{k_1+1},\ldots,f_n^{k_n+1})\subset\C\{z\}$ has the same
radical as $(f)$, hence is primary to the maximal ideal
$\mathfrak n=(z_1,\ldots,z_n)$ of $\C\{z\}$ --- which is not the symbol $\m$ of
\cref{contours:subsec:conventions} --- so $\mathfrak n^{a_k}\subset I_k$ for some
finite $a_k$. Ordinary residue annihilation makes the functional vanish on $I_k$
and therefore on $\mathfrak n^{a_k}$, so
evaluation on the Taylor polynomial of degree less than $a_k$ defines the
extension, and any higher truncation gives the same value. Uniqueness follows
because polynomials are dense for that filtration. This is the classical
finite-jet property of local residue functionals
\cite{GriffithsHarris,TajimaNabeshima}, not a new duality theorem.
\end{proof}

For $E$ and $H$ with formal coefficient germs, \eqref{contours:eq:generalres}
still defines $\Res_F$ using these extended functionals, and positive Hahn
support still makes the series strongly summable.

\begin{theorem}[Formal-coefficient contour theorem]
\label{contours:thm:formalcontour}
In \cref{contours:thm:untransformed,contours:thm:adaptedlift,contours:thm:untransformedhomotopy}
the hypotheses $E\in\A_n^n$ and $H\in\A_n$ may be replaced by
$E\in\C[[z]]((t^\Gamma))^n$ and $H\in\C[[z]]((t^\Gamma))$, with $f$ still
convergent with an isolated zero. The same cycle, the same support certificates
and the same homotopy apply, the pullback coefficients are still ordinary
holomorphic functions on the fixed parameter domain, and
\eqref{contours:eq:untransformed} holds with the finite-jet definition of the
residue.
\end{theorem}
\begin{proof}
The construction of $p,\mathsf X,\mathsf T,Z$ uses only $f$ and is unchanged.
Every substitution of a formal coefficient germ into $Z_0$, into $\Phi$ or into
the homotopy is licensed by \cref{contours:lem:radiusfreesub}, which was stated
for formal $h_\lambda$ precisely for this purpose; the lift of
\cref{contours:thm:adaptedlift} uses only such substitutions, derivatives of
formal germs, and finite coefficient operations, so its proof and its support
estimates survive verbatim. What must be redone is the residue equality, because
\cref{contours:lem:constantperiod} cannot be quoted for a divergent formal
numerator.

Fix an output exponent $\nu$. On the contour side the accounting
\eqref{contours:eq:exponentaccounting} admits only finitely many numerator
exponents, perturbation exponents and radial degrees at $\nu$, and each
contributing radial coefficient reads only a finite Taylor jet of the selected
formal germs. On the series side \cref{contours:lem:support} admits only finitely
many coefficient selections and multi-indices $k$ at $\nu$, and by
\cref{contours:lem:finitejet} their residues also read only finite Taylor jets.
Take the union of these two finite lists of jet requirements, and replace each
selected formal germ by a Taylor polynomial meeting all of them, setting the
other coefficients to zero. This introduces no new Hahn exponents and changes the
coefficient at $\nu$ on neither side. The already proved analytic theorem applies
to these polynomial coefficient data and identifies the two coefficients at
$\nu$; as $\nu$ was arbitrary, the protecting-cycle identity follows. Finally the
homotopy proof is coefficientwise ordinary Stokes applied to expressions that are
polynomial in the real parameter $u$ at each exponent, and never uses convergence
of an original coefficient germ, so it identifies the adapted period with the
protecting period as before.
\end{proof}

\emph{The finite-jet argument is a bridge between two of the rings of
\cref{contours:subsec:rings}, not an identification of them.} It does not assert
that a formal germ is an ordinary function on a complex neighborhood; it becomes
a Hahn family of holomorphic functions only after the positive-support
substitution. Nor does it construct the initial ramified covering when $f$ itself
is formal: that is a different extension problem, recorded in
\cref{contours:subsec:untransformedscope}.

\subsection{A mixed six-sheeted example, and the shortcut that fails}
\label{contours:subsec:sixsheeted}

\begin{example}[A six-sheeted singularity, its adapted cycle, and all its
monomial residues]
\label{contours:ex:sixsheeted}
Work over $\Gamma=\Q$ and take
\begin{equation}\label{contours:eq:examplef}
  f_1=x^2+y^3,\qquad f_2=x^2+3y^3 .
\end{equation}
Then $(f_1,f_2)=(x^2,y^3)$, so $\mu=6$, and the ordered-denominator
transformation matrix is
$\left(\begin{smallmatrix}1&1\\1&3\end{smallmatrix}\right)$ with determinant
$2$; already at the ordinary origin
$\res_{(f_1,f_2)}(xy^2)=1/2$. \emph{That factor $2$ is the tripwire of this
example}: any calculation that silently replaces the integrand by an
ordered-denominator transform of it, without the compensating determinant of
\cref{contours:lem:residuetransform}, will differ from the correct answer by
exactly this factor.

For a target point $w$ the inverse equations are
$x^2=(3w_1-w_2)/2$ and $y^3=(w_2-w_1)/2$, so the branch divisor lies in
$(3w_1-w_2)(w_2-w_1)=0$. Choose target weights $(1,2)$ and radial ramification
$N=6$, giving $(\kappa_1,\kappa_2)=(6,12)$, and on a polyannulus in coordinates
$(a,b)$ put $p(a,b)=(a^6,b)$, a covering of degree $6$. The ordinary family of
\cref{contours:thm:ramifiedleray} may be written explicitly as
\begin{equation}\label{contours:eq:exampleZ}
  Z_x=\sqrt{\tfrac32}\,s^3a^3
    \Bigl(1-\frac{s^6b}{3a^6}\Bigr)^{1/2},
  \qquad
  Z_y=-2^{-1/3}s^2a^2
    \Bigl(1-\frac{s^6b}{a^6}\Bigr)^{1/3},
\end{equation}
both binomial branches being $1$ at $s=0$, single-valued and holomorphic for
small ordinary $s$ on a fixed collar of $|a|=|b|=1$. Substitution gives
$f_1(Z)=s^6a^6$ and $f_2(Z)=s^{12}b$. As $a$ runs over the six preimages of a
fixed $a^6$, the powers $a^3$ and $a^2$ realize every pair of square-root and
cube-root choices exactly once, \emph{so this is the full six-sheeted cycle and
not one selected branch}.

Now deform by
\begin{equation}\label{contours:eq:exampleF}
  F_1=x^2+y^3-t,\qquad F_2=x^2+3y^3-5t,
\end{equation}
so $\delta=1$ is an admissible threshold and
$J_F=\det\left(\begin{smallmatrix}2x&3y^2\\2x&9y^2\end{smallmatrix}\right)
=12xy^2$. On \eqref{contours:eq:exampleZ} the coordinate $x$ has radial order
$3$ and $y$ has radial order $2$, the inverse Jacobian has pole order at most
$4$, and distinct sheets separate at radial order at most $3$; so one may take
$b_{\mathrm{pole}}=4$, $c_{\mathrm{sep}}=3$, $L=13$ and $\rho=1/40$, the
inequality \eqref{contours:eq:Lrho} reading $17/40<1$. With $s=t^{1/40}$ the
adapted roots of \cref{contours:thm:adaptedlift} are characterized by
\begin{equation}\label{contours:eq:examplePhi}
  \Phi_x^2=\frac{3s^6a^6-s^{12}b}{2}-s^{40},
  \qquad
  \Phi_y^3=\frac{s^{12}b-s^6a^6}{2}+2s^{40},
\end{equation}
with the branches selected by their leading terms, and then
$F_1(\Phi)=s^6a^6$, $F_2(\Phi)=s^{12}b$ exactly. Their leading displacements are
\[
  \Phi_x-Z_x=-\frac{\sqrt6}{6a^3}\,s^{37}+\cdots,
  \qquad
  \Phi_y-Z_y=\frac{2\cdot2^{2/3}}{3a^4}\,s^{36}+\cdots,
\]
the omitted terms having strictly larger integer radial support --- \emph{this
is a statement about supports, not a limit in any surreal topology}. The general
estimate $\vH(\Upsilon)\ge\delta-b_{\mathrm{pole}}\rho=36/40$ is attained by the
second coordinate, whereas the uniqueness ball $\vH(\Upsilon)>13/40$ of
\cref{contours:thm:adaptedlift} is much larger than the actual correction: it is
a uniqueness region, not a sharp bound.

The displaced zeros of \eqref{contours:eq:exampleF} satisfy $x^2=-t$ and
$y^3=2t$, six distinct points of $K^2$, and the monomial residues are
\begin{equation}\label{contours:eq:examplemonomials}
 \boxed{\quad
  \Res_F(x^\alpha y^\beta)=
  \begin{cases}
    \tfrac12(-t)^{p}(2t)^{q},
      &\alpha=2p+1,\ \beta=3q+2,\ p,q\in\N,\\[2pt]
    0,&\text{otherwise},
  \end{cases}\quad}
\end{equation}
for $\alpha,\beta\in\N$. An independent verification transforms the ordered
equations to $(x^2+t,\,y^3-2t)$, includes the factor $1/2$, and uses the
one-variable residues $(-t)^p$ of $x^\alpha/(x^2+t)$ and $(2t)^q$ of
$y^\beta/(y^3-2t)$; equivalently it reduces in the quotient basis
$1,y,y^2,x,xy,xy^2$ by $x^2=-t$ and $y^3=2t$ and halves the coefficient of
$xy^2$. In particular
\[
  \Res_F(xy^2)=\tfrac12,\quad
  \Res_F(x^3y^2)=-\tfrac t2,\quad
  \Res_F(xy^5)=t,\quad
  \Res_F(J_F)=6 ,
\]
the last checking the total degree of the cycle and detecting the loss of a
sheet or a reversed orientation. For the entire numerator $H=e^{x+y}$ the
positive substitution is legitimate and gives the exact Hahn series
\begin{equation}\label{contours:eq:exampleexp}
  \Res_F(e^{x+y})
  =\frac12\sum_{p,q\ge0}
   \frac{(-1)^p2^q}{(2p+1)!\,(3q+2)!}\;t^{p+q},
\end{equation}
whose constant coefficient is $1/4$, not $1/2$, because the $xy^2$ coefficient
of the numerator is $1/2$. \emph{This is a strong Hahn sum; ordinary numerical
convergence at a complex value of $t$ is irrelevant to the statement.} Neither
of the two independent verifications substitutes for the arbitrary-support proof
of \cref{contours:thm:untransformed}.
\end{example}

\begin{example}[The original coordinate torus is the wrong shortcut]
\label{contours:ex:wrongtorus}
In the setting of \eqref{contours:eq:examplef}, integrate the original form over
a small coordinate product circle in $(x,y)$ instead. On a torus with
$|y|^3\ll|x|^2$ both denominators are dominated by $x^2$; on one with the
opposite inequality both are dominated by $y^3$. Neither is the
equation-adapted Leray cycle. In the first region
\[
  \frac{xy^2}{(x^2+y^3)(x^2+3y^3)}
  =x^{-3}y^2\sum_{r,\sigma\ge0}(-1)^r(-3)^\sigma
   \frac{y^{3(r+\sigma)}}{x^{2(r+\sigma)}},
\]
which has no $x^{-1}y^{-1}$ term, so its coordinate-torus integral is $0$ while
$\Res_F(xy^2)=1/2$. \emph{This is why \cref{contours:def:hahncycle} must allow a
finite ramified cover:} the ramification fixes the geometry without multiplying
the integrand by an auxiliary determinant, and a Jacobian from the actual
parameterization remains, as it must for any pullback of a differential form.
The example also shows that the objection recorded in
\cref{contours:subsec:scope} was correct about the coordinate torus: it is the
\emph{cycle}, not the theorem, that had to change.
\end{example}

\begin{remark}[What the executable checks do and do not establish]
\label{contours:rem:leraychecks}
The source archive ships an exact symbolic suite, preserved in this repository
as \texttt{code/05-untransformed-leray-cycles-checks.py} with its recorded run
in \texttt{data/05-untransformed-leray-cycles-checks.json}: $632$ exact checks
pass under SymPy 1.14.0, in eight groups --- $180$ binomial-jet identities,
$225$ monomial residues, $200$ ideal-annihilation instances, $9$ coefficients of
\eqref{contours:eq:exampleexp}, $7$ quotient relations, $5$ target equations, $4$
leading displacements and $2$ normalizations. The arithmetic is exact over
$\Q(t)$ with no floating point, and the matrix route is deliberately independent
of the contour derivation: with $M_x,M_y$ the multiplication matrices on the
basis $1,y,y^2,x,xy,xy^2$, the trace
$\Tr\bigl(M_x^\alpha M_y^\beta(12M_xM_y^2)^{-1}\bigr)$ computes the sum of the
simple-zero residues of $x^\alpha y^\beta$ and is compared with
\eqref{contours:eq:examplemonomials}, which tests the determinant factor and all
six sheets at once. \emph{These are finite examples only.} They do not verify
well-ordering for every input support, parameterized removable singularities for
every germ, the rooted-tree construction of
\cref{contours:lem:positiveimplicit}, or novelty, and the source's own README and
recorded output say so.
\end{remark}

\subsection{Radius collapse and non-Archimedean supports}
\label{contours:subsec:pathologies}
Two pathologies exercise the hypotheses, and both are handled by the same cycle.

\begin{example}[Infinitely shrinking coefficient radii in the perturbation]
\label{contours:ex:radiuscollapse}
Keep \eqref{contours:eq:examplef} and take
\[
  E_1=-\sum_{j\ge1}\frac{t^j}{1-jx}\ \in\A_2,\qquad E_2=0,
\]
that is, the witness \eqref{contours:eq:shrinking} of
\cref{contours:subsec:rings} used now as a \emph{perturbation} rather than as a
numerator. Its support is $\N_{>0}$, so $\delta=1$, but its $j$th coefficient has
a pole at $x=1/j$, so there is no common ordinary neighborhood of the origin
carrying all of its coefficient germs, and any argument resting on a uniformly
convergent family on a fixed source disk fails outright. The protecting cycle
\eqref{contours:eq:exampleZ} with $s=t^{1/40}$ nevertheless works: on it
\[
  \frac1{1-jZ_x}=\sum_{r\ge0}j^rZ_x^{\,r}
\]
is a strongly summable Hahn identity, each exact output exponent sees only
finitely many pairs $(j,r)$ and finitely many radial coefficient choices, and
every coefficient function is a finite polynomial expression in the coefficient
functions of $Z_x$ and therefore holomorphic on the same annular parameter
domain. \emph{Their sizes may grow without bound as the exponent varies; no
coefficient norm was imposed or used.} The exact adapted lift also exists, by
\cref{contours:thm:adaptedlift}, although no closed formula such as
\eqref{contours:eq:examplePhi} is available: its coefficients come from the
finite rooted trees of \cref{contours:lem:positiveimplicit}, all on the same
parameter domain, with the certificate \eqref{contours:eq:Upsilonsupport}.
\end{example}

\begin{example}[A value group in which radial powers are not cofinal]
\label{contours:ex:nonarchsupport}
Let $\Gamma=\Q^2$ with the lexicographic order, and put $\rho=(0,1)$ and
$\eta=(1,0)$, so that $r\rho<\eta$ for every positive integer $r$. No sequence of
radial Taylor truncations indexed by $r\in\N$ is then a valuation approximation
at all scales --- the phenomenon of \cref{contours:ex:rank} --- yet
$\sum_{r\ge0}t^{r\rho}$ is a perfectly good Hahn sum with well-ordered support
and a single radial degree at each exact exponent. Combining this with radius
collapse and a support that lies in no finitely generated positive semigroup,
take
\[
  S=\bigl\{(0,1-1/j):j\ge2\bigr\}\cup\bigl\{(1,0)\bigr\}
\]
as the perturbation support, of order type $\omega+1$ with least element
$\delta=(0,1/2)$, and assign to its first part the coefficient germs
$(1-jx)^{-1}$. Choosing $\rho=\delta/40$ supplies the same protecting and lifting
inequalities as before, and the support proofs remain valid although infinitely
many radial degrees lie below $(1,0)$. \emph{Every proof above uses finiteness at
an exact exponent, never finiteness of an initial segment}, which is exactly the
distinction recorded after \cref{contours:ex:rank} and in
\cref{contours:rem:notpicard}.
\end{example}

\begin{remark}[Why positivity cannot simply be dropped]
\label{contours:rem:positivityneeded}
The positive-support condition is structural, not cosmetic. If arbitrary Hahn
perturbations were allowed one could take $E=-f$, leaving $F=0$ and no
denominator at all. Less crudely, a support unbounded below destroys the
finite-jet property of substitution: an expression such as
$\sum_{r\ge0}t^{-r^2}z^r$ does not lie in $\C[[z]]((t^\Gamma))$, its coefficient
support failing to be well ordered, and it must not enter the argument merely
because each individual coefficient exists. \emph{The results above do not assert
that positivity is necessary in every individual example}; they give one uniform
sufficient condition under which the cycle, the substitution, the geometric
inversion and the homotopy all carry verifiable support certificates.
\end{remark}

\subsection{A construction guide, and an explicit refusal of effectivity}
\label{contours:subsec:construction}
When the leading germ is polynomial or effectively algebraic the proof is an
actionable recipe, and it is worth stating separately from the workflow of
\cref{contours:sec:procedure} because its steps are specific to this cycle.
Compute a finite-map representative or a local elimination description; find a
nonzero discriminant equation; choose positive integer weights exposing one of
its monomials, as in \cref{contours:lem:targetmonomial}; compute the radial
monodromy, or simply use an exponent dividing $\operatorname{lcm}(1,\ldots,\mu)$
that kills it; retain the finite annular cover rather than forcing globally
single-valued branches over the target; and expand the source-coordinate map in
the ramified radial variable. For a protecting cycle only
\eqref{contours:eq:protectingrho} is required. For an adapted cycle, compute the
inverse-Jacobian pole bound $b_{\mathrm{pole}}$ and the sheet-separation bound
$c_{\mathrm{sep}}$ of \cref{contours:lem:radialbounds}, choose $L$ and $\rho$ by
\eqref{contours:eq:Lrho}, and solve \eqref{contours:eq:Bequation} by coefficient
recursion; the certificate \eqref{contours:eq:Pfull} is part of the output.

\emph{This is not an effectivity claim, and the difference matters for computer
algebra.} A requested individual Hahn coefficient uses only finitely many tree
terms, but the existence of a finite dependency set supplies no decision
procedure for membership in an arbitrary well-ordered subset of a general
ordered group --- the same caution already recorded for Sturm and sign
algorithms in \cref{contours:sec:jordan}. Practical implementations should use
rational lattices, explicitly presented support semigroups, algebraic
coefficient germs, and exact branch specifications. Arbitrary convergent germs
need not admit finite effective descriptions, no uniform complexity bound is
asserted for any class of inputs, and the scale estimates \eqref{contours:eq:Lrho}
are transparent but can be far from sharp, as
\cref{contours:ex:sixsheeted} shows.

\begin{center}\small
\begin{longtable}{@{}L{0.24\textwidth}L{0.34\textwidth}L{0.34\textwidth}@{}}
\toprule
Step & Established input & What must be checked here\\\midrule\endhead
Finite ordinary map & Isolated convergent zero; finite mapping theorem
  & Proper bounded representative, degree $\mu$, a nonzero
    discriminant\\[0.4em]
Target weighting, \cref{contours:lem:targetmonomial} & Convergent nonzero
    discriminant
  & Unique least-weight monomial and uniform annular avoidance\\[0.4em]
Radial ramification, \cref{contours:thm:ramifiedleray} & Finite covering
    monodromy
  & Kill the \emph{radial} permutation only, not all angular monodromy\\[0.4em]
Extension across $s=0$ & Bounded source representative
  & Holomorphic parameter extension and $Z(q,0)=0$\\[0.4em]
Radius-free substitution, \cref{contours:lem:radiusfreesub} & Positive
    well-ordered support
  & Finitely many contributions at each \emph{exact} exponent\\[0.4em]
Protecting cycle, \eqref{contours:eq:protectingrho} & $\kappa_*\rho<\delta$
  & Positive denominator errors and the certificate
    \eqref{contours:eq:Pzero}\\[0.4em]
Residue equality, \cref{contours:thm:untransformed} & Fixed-germ Laurent period
    identity
  & Strong regrouping, with no common radius for all germs\\[0.4em]
Adapted cycle, \cref{contours:thm:adaptedlift} & Inverse-Jacobian bound
    $b_{\mathrm{pole}}$; $L>b_{\mathrm{pole}}$
  & Positive implicit coefficients; the nullary-vertex support bound\\[0.4em]
Homotopy, \cref{contours:thm:untransformedhomotopy} & $L>\kappa_*$ and positive
    support
  & No denominator zero; coefficientwise closedness and Stokes\\[0.4em]
Formal extension, \cref{contours:thm:formalcontour} & Finite jets of ordinary
    residues
  & A separate finite-jet comparison, not an analytic assertion about a formal
    germ\\
\bottomrule
\end{longtable}
\end{center}

\subsection{What is classical, what is new, and what remains open}
\label{contours:subsec:untransformedscope}
Ordinary Leray cycles and the Grothendieck residue are classical
\cite{GriffithsHarris,TajimaNabeshima,Tsikh}, as are finite holomorphic maps,
analytic discriminants and Weierstrass preparation \cite{GrauertRemmert}. The
passage from a finite unramified cover to ramified single-valued analytic data
is an established ingredient of Abhyankar--Jung theory; Parusi\'nski and Rond
use a finite covering, a power substitution and removable singularities in their
analytic proof \cite[Prop.~2.1]{ParRond}. The positive support lemma is
classical \cite{Neumann,Higman}. \emph{None of these ingredients is claimed here
as an independent discovery}, and the particular one-radial-parameter version of
the ramification step is proved above rather than quoted.

What is proposed as new is their combination: an original-form period on a
single fixed finite-cover cycle, uniformly for every perturbation above one
specified positive valuation and for every radius-free or formal numerator, with
an explicit well-ordered support certificate, an exact adapted lift, and an
admissible homotopy inside the same chain category. \emph{The targeted literature
search of September 21, 2026 did not locate a published theorem with these
arbitrary-support, radius-free, uniform-protecting-cycle hypotheses; that is not
a priority certification}, and it is subject to the same consolidated status
statement as the rest of this article, \cref{contours:app:status}.

Three comparisons must not be collapsed. M\"uller and Strohmaier's
Hahn-holomorphic and Hahn-meromorphic functions are another genuine use of Hahn
series in complex analysis \cite{MuellerStrohmaier}, but their analytic
summation category is \emph{not} identified with the coefficientwise ring $\A_n$
of arbitrary radius-free germs used here; the two are different objects and no
comparison theorem between them is asserted. Second, the determinant-transformed
formula \eqref{contours:eq:generalrepresentation} remains valid and useful; the
new theorem answers a different question and does not supersede it. Third, the
lift of \cref{contours:thm:adaptedlift} refines, but does not reprove, the
support-preserving lift of the ordinary root cover with unchanged monodromy of
the companion report \emph{Finite Hahn Deformations: Division, Conservation of
Multiplicity, and Residue Duality} \cite{FiniteDeformations} at
\texttt{docs/surcomplex/finite-deformations}: there a coherent root cover is
lifted with its monodromy intact, whereas here the finite \emph{radial}
monodromy is deliberately killed, before any Hahn substitution, in order to
obtain one fixed cover. That report is cited for the lift rather than having it
rebuilt.

\paragraph{Exact boundaries.}
All supports, covers and parameter manifolds are set-sized, and every
coefficientwise integral is an ordinary integral over a compact real manifold.
The Hahn parameterization is \emph{not} asserted continuous for the fine surreal
topology, and \cref{contours:thm:untransformedhomotopy} is not a
fine-topological homotopy; no surreal-valued measure and no unrestricted
Riemann-sum construction is invoked, consistently with
\cref{contours:prop:fullfine,contours:prop:mesh}. The input must have an isolated
\emph{convergent} ordinary leading map and a positive-support perturbation: an
arbitrary Hahn tuple with no such leading map, and a non-isolated leading zero,
are both out of scope. The theorems produce no canonical choice-independent
contour; they prove that the \emph{period} is independent of the choices, by
\cref{contours:cor:scalecancel}. The parameter space may be disconnected, and no
claim is made that every admissible cycle is homologous to the chosen one.
\Cref{contours:cor:protectedsheets} gives $\mu$ distinct regular points over each
protected target point but does not assert that every zero of $F$ is simple or
that the deformed zero algebra is reduced, and
\cref{contours:prop:idealannihilation} establishes nonzero descent without a
perfect-pairing or flatness theorem.

\paragraph{What is now more precise, and still open.}
First, whether the convergent leading map can be replaced by an arbitrary
\emph{formal} isolated map while keeping a fixed ordinary compact parameter
manifold. \Cref{contours:thm:formalcontour} removes convergence from $E$ and
$H$ only; the construction of the cycle depends on a bounded analytic finite
covering obtained before any Hahn substitution, and formal algebra alone does
not supply that object. Second, effective scale selection: an intrinsic
invariant of the finite map controlling the best protecting and lifting scales
would improve both the mathematics and exact symbolic implementations, and the
attained bound $36/40$ in \cref{contours:ex:sixsheeted} shows that the
inverse-Jacobian pole order carries real information even where the uniqueness
radius is conservative. Third, a homology comparison among different protecting
weights, covers and microscopic scales: equality of the periods is proved by
\cref{contours:cor:scalecancel}, but a chain-level comparison in a precisely
specified larger Hahn category is strictly stronger and does not follow from
period equality. This last item is the surviving part of the
\emph{admissible scale atlas} question of \cref{contours:subsec:problems}.

\section{An algebraic contour pairing for rational differentials}
\label{contours:sec:algebraic}

The Hahn constructions above rely on an ordinary base or a useful infinitesimal
chart. Rational functions admit a different, fully algebraic closed-period theory
over any fixed real closed workspace, requiring neither an infinite sum nor a
limiting partition. In this section the coefficient ring is the rational function
field $K(z)$ over $K=R[i]$, $R$ real closed containing $\R$; it is none of the four
rings of \cref{contours:subsec:rings}, and no statement below is transferred to or
from them except through the explicitly proved comparison
\cref{contours:prop:comparison}.

\subsection{Definition from definable winding}
Let $\gamma:[0,1]_R\to K$ be a semialgebraic continuous loop, piecewise smooth when
differential notation is used, or more generally a finite oriented semialgebraic
cycle. For $a\notin\gamma$ write $\wind_{\mathrm{sa}}(\gamma,a)\in\Z$ for its
semialgebraic winding number as in \cref{contours:sec:jordan}, normalized to $1$ for
a positively oriented circle about an inside point and $0$ outside.

\begin{definition}[Algebraic rational contour functional]
\label{contours:def:algpairing}
For a rational differential $\omega=Q(z)\dd z$, $Q\in K(z)$, with no pole on
$\gamma$, define
\begin{equation}\label{contours:eq:algebraicpairing}
 \boxed{\qquad
 I_{\mathrm{alg}}(\gamma,\omega)
 :=2\pi i\sum_{a\in K}\wind_{\mathrm{sa}}(\gamma,a)\,\res_a\omega .
 \qquad}
\end{equation}
\end{definition}

Only the finitely many finite poles contribute, so the sum is finite. As fixed in
\cref{contours:subsec:residuedef}, the factor $2\pi i$ is a normalization selected
to match ordinary complex integration using the distinguished copy of $\C$ inside
$K$; \emph{it is not forced by ordered-field arithmetic}, and over an abstract real
closed field with no distinguished copy of $\R$ one would normalize formally or use
$I/(2\pi i)$. A residue at infinity could be included after compactifying and
changing charts, but is unnecessary for this affine formula and supplies no new
generator, its residue being the negative sum of the finite ones.

\emph{Equation \eqref{contours:eq:algebraicpairing} is the definition of a pairing.}
It is not a residue theorem deduced from a previously constructed Riemann integral,
and it is not a claim that a finite-partition Riemann net converges to this value.
Its ingredients are algebraic local residues and the independently defined definable
degree. It is useful because it requires neither ordinary leading separation nor a
global logarithm: poles and contour coefficients may lie at arbitrary Hahn scales.

\begin{theorem}[Characterization and algebraic Cauchy theory]
\label{contours:thm:algpairing}
For a fixed semialgebraic loop $\gamma$, \eqref{contours:eq:algebraicpairing} is
\emph{the unique} $K$-linear pairing on rational differentials without poles on
$\gamma$ which vanishes on exact rational differentials and satisfies
\[
  I_{\mathrm{alg}}\Bigl(\gamma,\frac{\dd z}{z-a}\Bigr)
  =2\pi i\,\wind_{\mathrm{sa}}(\gamma,a).
\]
It is additive under concatenation of loops and additive in cycles, changes sign
under orientation reversal, and is invariant under definable homotopies avoiding the
poles. If all indices at the poles vanish, the integral is zero. If $\gamma$ is a
positively oriented semialgebraic Jordan curve and $Q$ has no poles on its bounded
side or on its boundary, then $I_{\mathrm{alg}}(\gamma,Q\dd z)=0$; and if $a$ lies on
that bounded side, then
\[
  I_{\mathrm{alg}}\Bigl(\gamma,\frac{Q(z)}{z-a}\dd z\Bigr)=2\pi i\,Q(a).
\]
Finally, for a nonzero $Q\in K(z)$ with neither zeros nor poles on $\gamma$,
\begin{equation}\label{contours:eq:algargument}
  \frac1{2\pi i}I_{\mathrm{alg}}\Bigl(\gamma,\frac{Q'}{Q}\dd z\Bigr)
  =\sum_a\wind_{\mathrm{sa}}(\gamma,a)\,\ord_a(Q),
\end{equation}
negative orders counting poles.
\end{theorem}
\begin{proof}
Since $K$ is algebraically closed of characteristic zero, partial fractions give
$Q(z)=P(z)+\sum_a\sum_{j=1}^{m_a}c_{a,j}(z-a)^{-j}$. The polynomial part has a
polynomial primitive, and for $j\ge2$
\[
  \frac{\dd z}{(z-a)^j}=d\Bigl(\frac{(z-a)^{1-j}}{1-j}\Bigr).
\]
Thus a pairing with the stated properties must retain exactly the $j=1$ terms with
the displayed winding coefficients; this proves uniqueness and simultaneously
verifies that \eqref{contours:eq:algebraicpairing} has those properties, since the
residue of a rational derivative is zero (differentiating a Laurent term cannot
produce exponent $-1$). Additivity, orientation behaviour and homotopy invariance
follow from the corresponding properties of $\wind_{\mathrm{sa}}$. For a positively
oriented Jordan curve, exterior poles have winding zero; in the Cauchy-kernel
expression the only contributing interior pole is $a$, with residue $Q(a)$. If
$Q=(z-a)^mV$ with $V(a)\ne0$, then $Q'/Q=m/(z-a)+V'/V$ has residue $m$ at $a$,
giving \eqref{contours:eq:algargument}.
\end{proof}

Equation \eqref{contours:eq:algargument} is a rational argument principle over the
full real closed workspace; for a simple positively oriented Jordan cycle its right
side is the number of inside zeros minus inside poles, with multiplicities, an exact
arbitrary-scale counterpart of the classical statement. Eisermann's polynomial
Cauchy and root-counting theorems provide the effective winding component for
polynomially described contours \cite{Eisermann,PR}, subject to the effectivity
caveat of \cref{contours:sec:jordan}.

\subsection{Residues are the complete rational obstruction}
\begin{proposition}[Rational de Rham quotient]
\label{contours:prop:rationalcohomology}
Over the algebraically closed characteristic-zero field $K$,
\begin{equation}\label{contours:eq:rationalcohomology}
  \frac{K(z)\dd z}{dK(z)}\ \simeq\ \bigoplus_{a\in K}K\Bigl[\frac{\dd z}{z-a}\Bigr],
\end{equation}
and the coordinates of a class are its residues at the finite points.
\end{proposition}
\begin{proof}
Partial fractions express every rational differential as a polynomial differential
plus finitely many pole terms. Polynomial terms integrate by division by positive
integers; for $j\ge2$, $(z-a)^{-j}\dd z$ is the derivative of
$-(z-a)^{-j+1}/(j-1)$. Only simple poles remain. The residue of a derivative is
zero, so a nonzero linear combination of distinct simple-pole terms cannot be exact.
This proves both spanning and independence.
\end{proof}

Consequently \eqref{contours:eq:algebraicpairing} is the winding pairing with the
\emph{entire} rational de Rham quotient, not an arbitrary selection of part of the
rational information. The corresponding local statement is purely formal: for
$\omega=\sum_{j\ge j_0}c_ju^j\dd u\in K((u))\dd u$ there is a formal Laurent
primitive exactly when $c_{-1}=0$, since every other term integrates by division by
$j+1$ while differentiating a Laurent series can never produce a $u^{-1}$
coefficient. Hence
\begin{equation}\label{contours:eq:formalresclass}
  K((u))\dd u\,\big/\,dK((u))\ \simeq\ K,\qquad[\omega]\longmapsto c_{-1},
\end{equation}
and the residue is invariant under an invertible formal coordinate change
$u=v\,h(v)$, $h(0)\ne0$: exact differentials remain exact, and
$\dd u/u=\dd v/v+h'(v)\dd v/h(v)$, whose second term has zero residue. Residues
therefore give a coordinate-invariant obstruction to exactness independently of any
connected topology. This also explains the difference from the fine-local category:
a global fine-analytic primitive of $1/z$ may exist after a phase convention, but no
rational primitive exists, and the cohomology computed in
\eqref{contours:eq:rationalcohomology} is that of a differential algebra, not
unrestricted fine-local cohomology.

Finite rational and semialgebraic data over $\SC$ can always be placed in a
set-sized real closed workspace: the relevant roots lie in its algebraic closure
$K$, and enlarging the workspace changes neither finite sign computations nor the
residue formula. This gives an especially robust full-surcomplex interpretation for
finite algebraic input.

\subsection{Agreement with Hahn integration on a verified overlap}
\label{contours:subsec:notconflated}
\begin{proposition}[Comparison on separated charts]
\label{contours:prop:comparison}
Work in the Hahn workspace $K=R_\Gamma[i]$ of \eqref{contours:eq:fields}.
Let $a\in K$, $r\in R_\Gamma$ with $r>0$, and let $C=a+r\psi_R$ be a
semialgebraic loop, where $\psi_R$ is piecewise polynomial over the workspace
on a finite partition with ordinary endpoints. Suppose its
restriction to ordinary parameters is an admissible Hahn-deformed closed contour
$\gamma_0+\eta$, and on the full field interval $\psi_R$ is uniformly
infinitesimally close to the semialgebraic extension of its ordinary
real-coefficient leading loop $\gamma_0$. Suppose each
pole $b$ of $\omega=Q\dd z$, $Q\in K(z)$, has either finite normalized position
$(b-a)/r$ with standard part off the ordinary trace of $\gamma_0$, or infinite
normalized position. Then, whenever the rational pulled-back differential is a
coherent family on a common ordinary collar,
\[
  I_{\mathrm{alg}}(C,\omega)=\HInt_{a+r(\gamma_0+\eta)}\omega .
\]
The same holds after a polynomial infinitesimal deformation when the homotopy is
both definably pole-avoiding and admissible for Hahn pullback.
\end{proposition}
\begin{proof}
Use finite partial fractions after the affine substitution. Polynomial terms have
coherent primitives and zero period; a term of order greater than one has a rational
primitive, coherent on the coefficient collar, and also zero period. For a finite
pole $b$ with normalized position $c+\varepsilon$, ordinary compactness gives a
positive ordinary separation between $\gamma_0$ and $c$; the real-coefficient
semialgebraic description has the same lower bound on the field interval, and
uniform infinitesimal closeness makes the straight definable homotopy from
$(\psi_R,c+\varepsilon)$ to $(\gamma_{0,R},c)$ avoid zero differences throughout,
where $\gamma_{0,R}$ is the semialgebraic extension of $\gamma_0$. Thus
$\wind_{\mathrm{sa}}(C,b)=\wind(\gamma_0,c)$. \Cref{contours:thm:winding} then
gives equality of the two pairings on $\dd z/(z-b)$. If the normalized pole is
infinite, its inverse is infinitesimal, the reciprocal has a coherent expansion in
nonnegative powers of the finite chart variable with zero period, and the
field-valued polynomial loop lies in the bounded part of that chart so the pole has
index zero. Exact rational differentials have zero period by the coefficientwise
fundamental theorem on the admitted pullbacks. Partial fractions now prove the
assertion term by term, the affine scale factors cancelling as in
\eqref{contours:eq:affineintegral}. For the last sentence, use the specified
definable homotopy for algebraic indices and
\cref{contours:thm:homotopy,contours:thm:endpoints} for the coherent integrals.
\end{proof}

The explicit overlap hypotheses prevent an unwarranted claim that every smooth
Hahn-deformed path extends to a semialgebraic map on $[0,1]_{R_\Gamma}$; most do
not. The proposition says the two theories agree when their objects and their
homotopies really can be compared. Agreement of the two constructions is
\emph{proved where both are defined}, not postulated across different chain
categories.

\Cref{contours:prop:overlap,contours:prop:comparison} are two separate
compatibility results and must not be conflated. The first compares two
\emph{chart descriptions of one chain} within the Hahn category; the second
compares \emph{two different integration theories}. Here is the explicit case where
conflating them fails. For the full semialgebraic unit circle,
\eqref{contours:eq:algebraicpairing} gives immediately
\begin{equation}\label{contours:eq:algcircleboundary}
  I_{\mathrm{alg}}\Bigl(S^1,\frac{\dd z}{z-(1-t)}\Bigr)=2\pi i,
  \qquad
  I_{\mathrm{alg}}\Bigl(S^1,\frac{\dd z}{z-(1+t)}\Bigr)=0 .
\end{equation}
These are legitimate values \emph{in exactly the case where the unmodified
macroscopic Hahn pullback fails}, namely the boundary-layer situation of
\cref{contours:subsec:boundarylayer} and \cref{contours:ex:boundarylayer}. They do
not contradict \cref{contours:prop:comparison}, whose leading-separation hypothesis
is absent here; and they are not obtainable from \cref{contours:prop:overlap}, which
presupposes admissible pullbacks that do not exist in this case.

\subsection{What this pairing does not determine}
For an \emph{open} path, the differential $\dd z/(z-a)$ requires a logarithmic
increment, not just a residue and a winding integer; a choice of path continuation
or a more elaborate period object is needed, and a global fine logarithm is not
automatically that choice by \cref{contours:sec:topologies}. For nonrational
analytic functions, partial fractions need not leave a rational exact remainder, so
\eqref{contours:eq:algebraicpairing} is \emph{not} a definition on all of $\HH(U)$
or $\A_n$; those larger classes are handled by the coherent and microscopic
constructions, not by pretending that the rational formula already covers them. The
functional also does not by itself extend to a general Hahn-holomorphic function
with an infinite collection of coefficient singularities.

\section{Root-of-unity sampling: coefficientwise and valuation limits}
\label{contours:sec:sampling}

Two of the merged manuscripts leave opposite impressions about root-of-unity
quadrature: one presents it as a negative result, the other devotes a section to its
positive side. These conclusions use different notions of convergence.
Common-support continuous coefficient functions suffice for coefficientwise
recovery. A separate decay condition guarantees a valuation limit; the examples
below show that ordinary coefficient convergence alone does not guarantee it.

\subsection{What the classical construction assumes}
For a complete algebraically closed field with a real-valued non-Archimedean
absolute value, a useful circle integral is a limit of the weighted samples
\begin{equation}\label{contours:eq:samples}
  S_N(f;a,r)=\frac rN\sum_{\xi^N=1}f(a+r\xi)\,\xi ,
\end{equation}
the $\xi$ being the roots of unity in that field. Cherry's presentation of the
Schnirelman integral uses this normalization, takes $N\to\infty$ with $|N|=1$, and
defines the integral only when the field-valued limit exists \cite[\S1.3]{Cherry};
its residue normalization is $1$, not $2\pi i$, and the chosen $r$ is part of the
datum for arbitrary functions, not just its absolute value. In our Hahn field with
coefficient field $\C$, these are exactly the ordinary complex roots of unity.
The natural valuation has residue characteristic zero and
$v(N)=0$ for every positive integer. Related averages appear in $p$-adic line
integral theories \cite{Diamond,Ettayb}.

The basic algebraic fact is exact and requires no limit. If
$f(a+w)=\sum_{j=j_0}^{j_1}c_jw^j$ is a Laurent polynomial, then
\begin{equation}\label{contours:eq:rootaverage}
  S_N(f;a,r)=\sum_{\substack{j_0\le j\le j_1\\ N\mid j+1}}c_jr^{j+1},
\end{equation}
from $N^{-1}\sum_{\xi^N=1}\xi^q=1$ when $N\mid q$ and $0$ otherwise. When $N$
exceeds every nonzero $|j+1|$ in the finite support, the average equals $c_{-1}$,
which is exactly the corresponding normalized contour integral. \emph{This is exact
algebraic quadrature for finite Laurent data; it is not a convergence theorem.} The
words ``complete Hahn field'' alone do not justify a limit: a common ordinary
analytic coefficient expansion and a convergent non-Archimedean Laurent expansion
are different hypotheses.

\subsection{A rank-independent sufficient condition for a valuation limit}
\begin{proposition}[A precise sampling criterion]
\label{contours:prop:samplingcriterion}
Suppose $A_k\in K$, $k\in\Z$, satisfy the \emph{decay condition}
\begin{equation}\label{contours:eq:decay}
  \text{for every }\gamma\in\Gamma,\qquad
  \{k\in\Z:v(A_k)\le\gamma\}\ \text{is finite}.
\end{equation}
Then $f(\xi)=\sum_{k\in\Z}A_k\xi^k$ is strongly defined at every root of unity, and
\begin{equation}\label{contours:eq:alias}
  S_N(f;0,1)=\sum_{\ell\in\Z}A_{\ell N-1}\ \longrightarrow\ A_{-1}
\end{equation}
in the valuation topology as the ordinary integer $N$ tends to infinity. This holds
for $\Gamma$ of arbitrary rank.
\end{proposition}
\begin{proof}
For a fixed exponent bound $\gamma$, only finitely many $A_k$ can contribute below
that bound; their supports form a finite union of well-ordered sets. This proves
that the full union is well ordered --- a descending sequence would be bounded above
by its first term and hence lie in that finite union --- and also proves local
finiteness of contributions. Multiplication by a root of unity does not change
support, so the series defining $f(\xi)$ is strong. The finite root-of-unity sum
gives $N^{-1}\sum_{\xi^N=1}\xi^{k+1}=1$ when $N\mid k+1$ and $0$ otherwise, and
strong summability permits applying this finite linear operation, proving the alias
formula. Given $\gamma$, choose $N_0$ so large that the finitely many $k\ne-1$ with
$v(A_k)\le\gamma$ all satisfy $|k+1|<N_0$; for $N\ge N_0$ none of them has
$N\mid k+1$, so every remaining nonprincipal term in \eqref{contours:eq:alias} has
valuation greater than $\gamma$, and so does the error. This is exactly valuation
convergence.
\end{proof}

\emph{The hypothesis is doing all the work, and its strength should be stated
plainly.} Condition \eqref{contours:eq:decay} demands eventual domination of
\emph{every} value-group threshold; it is strictly stronger than the existence of a
strong sum. Moreover, when no countable subset is cofinal in the positive value
group, a nontrivial infinite sequence of finite valuations \emph{cannot} satisfy
this tail condition. This is one reason sequential constructions have a limited role
at very large ranks.

\subsection{Coefficientwise recovery always has a different meaning}
Let $F=\sum_{\lambda\in S}f_\lambda t^\lambda$ have one well-ordered support and
ordinary continuous coefficient functions on the unit circle. In the fixed support
space $\C^S$ use the product topology of the ordinary complex coefficient
topologies; this is a topology on \emph{coefficient presentations} and is not the
natural valuation topology of $K$.

\begin{theorem}[Coefficientwise recovery of a Hahn contour]
\label{contours:thm:sampling}
The samples $S_N(F;0,1)$ have support contained in $S$ and
\begin{equation}\label{contours:eq:coefflimit}
  \operatorname*{lim}_{N\to\infty}^{\mathrm{coefficientwise}} S_N(F;0,1)
  =\frac1{2\pi i}\HCoint_{|z|=1}F(z)\dd z ,
\end{equation}
the right side being a Hahn series with support contained in $S$. If the
$f_\lambda$ are holomorphic on a common ordinary annulus --- that is, if
$F\in\HH$ of that annulus --- the coefficient at $\lambda$ is their ordinary Laurent
residue.
\end{theorem}
\begin{proof}
At exponent $\lambda$ the sample is
$N^{-1}\sum_{j=0}^{N-1}f_\lambda(e^{2\pi ij/N})e^{2\pi ij/N}$, an ordinary Riemann
sum for $(2\pi)^{-1}\int_0^{2\pi}f_\lambda(e^{i\theta})e^{i\theta}\dd\theta$, which
is the normalized ordinary contour integral; its convergence requires no uniformity
in $\lambda$. Taking these limits creates no exponents outside the fixed
well-ordered $S$, so the result is a legitimate element of $K$. The annulus
assertion is the ordinary Laurent coefficient formula.
\end{proof}

This realizes the already defined coefficientwise integral by finite samples in a
\emph{specified coefficient topology}. \emph{It does not make that integral a fine
limit, and in general it does not make it a valuation limit either.} Taking ordinary
coefficient limits is a legitimate alternative description of the Hahn contour
functional; it is not a limit in the internal valuation topology.

\subsection{Why coefficientwise recovery need not give a valuation limit}
\begin{example}[A failed valuation-limit contour definition]
\label{contours:ex:quadrature}
For $a=0$, $r=1$ and $f(z)=1/(z-b)$ with ordinary real $b=1/2$, logarithmic
differentiation of $X^N-1$ gives
\[
  S_N(f;0,1)=\frac1N\sum_{\xi^N=1}\frac{\xi}{\xi-b}=\frac1{1-b^N} .
\]
The ordinary complex limit is $1$, but $v(S_N(f)-1)=0$ for \emph{every} $N$, so
there is no convergence to $1$ in the Hahn valuation topology. In the full fine
topology it cannot converge at all, since it is not eventually constant. If instead
$b=t^\rho$ with $\rho>0$, then $v(S_N(f)-1)=N\rho$: a rank-one valuation can see
convergence, but a higher-rank group containing an element larger than every $N\rho$
still cannot.
\end{example}

\begin{example}[Ordinary convergence invisible to the Hahn valuation]
\label{contours:ex:ordinarysampling}
Take the ordinary $f(z)=1/(1-z/2)$ on a neighborhood of the unit disk, embedded as
the sole coefficient at exponent zero. Then
\begin{equation}\label{contours:eq:ordinarysample}
  S_N(f;0,1)=\frac{2}{2^N-1}.
\end{equation}
It tends to $0$ ordinarily and coefficientwise, agreeing with its zero contour
integral, but every displayed nonzero number has Hahn valuation $0$ and distinct
samples differ by nonzero real amounts of valuation $0$, so the sequence is not
valuation Cauchy. Its ordinary Taylor coefficients $2^{-k}$ do not tend to zero in
the natural non-Archimedean valuation. \emph{No classical Schnirelman theorem on
valuation-convergent analytic series is being contradicted.}
\end{example}

\begin{example}[An obstruction occurring only in higher rank]
\label{contours:ex:highranksampling}
In $\Gamma=\Q+\Q\omega$ let $f(z)=1/(1-tz)$. Its coherent unit-circle integral is
zero, but
\begin{equation}\label{contours:eq:highranksample}
  S_N(f;0,1)=\frac{t^{N-1}}{1-t^N},\qquad v(S_N)=N-1<\omega .
\end{equation}
The samples fail to converge to zero in the valuation topology, although each fixed
Hahn coefficient is eventually zero. In a rank-one workspace with value group $\Q$
the same formula does converge to zero. In the full fine topology it cannot
converge, since it is never eventually zero. \emph{One exact formula therefore
separates three notions of convergence.}
\end{example}

Both sample identities follow without any infinite-series argument: from
$\prod_{\xi^N=1}(1-q\xi)=1-q^N$, logarithmic differentiation gives
$N^{-1}\sum_{\xi^N=1}\xi/(1-q\xi)=q^{N-1}/(1-q^N)$; substitute $q=1/2$ and $q=t$.
This calculation is also an effective diagnostic: before interpreting root-of-unity
samples as an integral in a chosen topology, check whether their error valuations
dominate the thresholds that topology actually uses. Finite Laurent data give exact
algebraic quadrature; general coherent coefficients give ordinary coefficient
limits; valuation convergence needs the additional cofinality and regularity
hypothesis \eqref{contours:eq:decay}.

\section{Other established approaches and their exact scope}
\label{contours:sec:alternatives}

The following theories are genuine precedents or alternatives. They do not all
integrate the same forms, take values in the same field, or use the same notion of
path. The comparisons are intentionally narrower than a claim of equivalence, and
importing a theorem name without its hypotheses would be misleading.

\subsection{Semialgebraic chains, Stokes, and the target of integration}
Ohmoto and Shiota prove a $C^1$ triangulation theorem for semialgebraic sets and
formulate integration and Stokes using finite differentiable simplices
\cite[Thm.~4.1, Prop.~4.2]{OS}. This is an important existing counterpart, not
merely a proposal. Their non-Archimedean discussion explicitly warns that integral
values need not remain in the given field and that a larger saturated target may
enter; it should not be read as a construction of the particular $K$-valued contour
operator defined here.

There is a further codomain distinction in the cited literature.
Ma\v r\'{\i}kov\'a and Shiota construct a measure with values in an ordered
semiring, a Dedekind completion of a quotient, with a volume-preserving change of
variables \cite{MS}; Berarducci and Otero's earlier additive measure is real-valued
on definable subsets of the finite part \cite{BOmeasure}; Kaiser develops
semialgebraic measure and integration over non-Archimedean real closed fields
containing $\R$ with \emph{Archimedean} value group, using valuation-section data,
and his construction can enlarge the value ring --- for Puiseux series a polynomial
ring in an additional variable occurs \cite{Kaiser}. These are different targets
retaining different information, and \emph{a semiring-valued measure is not
automatically a signed $K$-valued complex line integral}. A reference to any of them
is not a proof that every signed or complex differential form has a canonical
integral in the original Hahn field. Kaiser's setting also warns that integrals of
rational functions may require logarithmic values outside the initial workspace ---
which the polynomial theorem \cref{contours:thm:polstokes} avoids by using finite
factorial identities, and which \cref{contours:thm:algpairing} avoids by using
residues rather than open-path logarithms.

For a surcomplex application one must therefore specify the scalar integration
operator on each simplex, its target, and an embedding or comparison map if values
in $\SC$ are wanted. \emph{Triangulation supplies a finite chain decomposition; it
does not by itself solve the scalar integration problem.} In particular, saturation
alone does not make an ordinary sequence $1/N$ converge to zero in the full ordered
topology of a field containing infinitesimals, so the finite-mesh obstruction
\cref{contours:prop:mesh} remains unless the limiting notion has been changed. A
promising combination is to use definable triangulation to organize domains and then
impose Hahn-coherent pullback conditions on each simplex; then
\cref{contours:thm:hahnstokes,contours:thm:deformedstokes} give Stokes whenever
those checks succeed. \emph{A theorem that all desired definable data admit such a
finite coherent atlas would be an additional result, not a consequence of
triangulation alone.}

\subsection{Berkovich geometry changes the space, not just the topology}
For a rank-one subgroup $\Gamma\subset\R$ put $|x|_v=\exp(-v(x))$. The Hahn field is
then complete and algebraically closed, hence a legitimate base for Berkovich
geometry. The Berkovich line adds multiplicative seminorms on $K[T]$, not merely
more field points; it is tree-like and uniquely arc-connected, and deleting a branch
point can create infinitely many components \cite[\S2.3.3]{Temkin}, \cite{BR}. Thus
the enriched line acquires meaningful paths but \emph{not} a planar Jordan circle:
two distinct arcs around a circle would contradict uniqueness, and an annulus has an
interval skeleton rather than a circular one. The appropriate boundary intuition is
a disk point with outgoing directions, or the ends of an annular skeleton, not a
fine-connected curve in $K$.

The rank restriction is substantive. An arbitrary higher-rank $\Gamma$ cannot be
embedded order-preservingly into $\R$; passing to a real-valued coarsening may lose
infinitesimal scales, while retaining them requires a different higher-rank
framework, and \cref{contours:ex:rankcircles} exhibits distinctions that no single
real-valued valuation can retain. Also $|c|_v=1$ for every nonzero ordinary complex
constant $c$, however small its usual complex modulus, so ordinary coefficient
convergence is not non-Archimedean convergence. \emph{No canonical reduction of all
surreal scales to one rank is asserted here.}

\subsection{Tropical differential forms and Stokes on Berkovich spaces}
Chambert-Loir and Ducros construct real $(p,q)$-forms and currents on Berkovich
spaces using tropicalizations and calibrated skeleta, with integration of
top-bidegree forms and a Stokes formula; the expanded 2025 revision develops this in
a broad tropical setting and includes a tropical Stokes formula
\cite[Prop.~5.11.1]{CLD}. This is a powerful established analogue of part of complex
differential geometry, and a reason not to regard $d^2=0$ as an integration theorem
by itself.

Its integrals are real-valued and its geometric objects are enriched analytic
spaces. The $K$-valued periods of $F(z)\dd z$ constructed here, with the chosen
normalization $2\pi i$, are different objects. A comparison could associate tropical
quantities to valuations or norms of analytic data, but one should not expect that
operation alone to recover complex phase or the cancellation of opposite residues
such as the $\pm1/(2t)$ of \cref{contours:ex:cluster}. A theorem connecting the
pairings would need additional hypotheses and an explicit map on both forms and
chains. \emph{A Berkovich Stokes theorem cannot be quoted as the proof of the
specific contour identities in the supplied manuscripts.}

\subsection{Perfectoid analytic cycles: an actual Stokes theorem, but a different
base}
Mihara constructs analytic singular homology from a cosimplicial perfectoid object
and integration over a restricted class of integrable analytic simplices, with
\cite[Thm.~5.15]{Mihara} proving a Stokes formula for those simplices; in the
integration part of the current text the base field is a finite extension of
$\Q_p$ and the target is the de Rham period ring $B_\dR$. This is directly relevant
as evidence that replacing ordinary simplices and carefully restricting
integrability can produce a non-Archimedean cycle pairing.

It is \emph{not} a theorem for the present natural Hahn valuation with residue field
$\C$: the residue characteristic and the integration target both differ, and the
positive-residue-characteristic hypothesis is decisive. Citing an older title or an
abstract about ``non-Archimedean integration'' would not justify silently applying
its period-ring construction to all surcomplex numbers. The common lesson is
methodological: a chain category and its face compatibility must be constructed
before a Stokes assertion is meaningful.

\subsection{Surreal antiderivatives, resurgent extensions, and nonstandard bridges}
Costin and Ehrlich construct extensions and integration for substantial classes of
resurgent real functions, with treatments of singularities and ordered exponential
subfields \cite{CE}. Their framework is a natural candidate when a contour is
parameterized by functions whose real and imaginary pullbacks belong to those
classes, and it can retain asymptotic information that a mere standard-part integral
would lose. It is not merely the coefficientwise integral on a fixed ordinary
domain.

To turn that into a contour theorem, however, one must prove closure under the
required compositions $F(\gamma(s))\gamma'(s)$, specify logarithmic branches and
endpoint conventions, and establish substitution and homotopy theorems; agreement
with Hahn integration should first be tested on common analytic germs, rational
forms, and the microscopic residues above. \emph{Extending a real one-variable
integral to surreal endpoints is not automatically a construction of complex contour
integration}: complexification, phase and branch data, pullback along a path, and
compatibility with a two-dimensional boundary operator remain to be established.
Outside such an overlap, equality of integration operators should be a theorem, not
an assumption. There is no conflict between the breadth of these function-extension
theories and the arbitrary smooth coefficient forms used on ordinary parameter
manifolds in \cref{contours:sec:derham}: the latter construction integrates ordinary
functions first and stores their values as Hahn coefficients, and does not assert a
canonical fine-smooth extension of every ordinary smooth function to every
surcomplex point.

Ducros, Hrushovski, and Loeser construct a comparison between complex and
non-Archimedean integrals using nonstandard models and compatible differential-form
structures \cite{DHL}. Their work shows that rigorous bridges exist, and also how
much additional structure a bridge requires. Real closedness alone is not a transfer
principle for arbitrary analytic functions, manifolds, or hyperfinite sums: a
nonstandard model with a specified star map can transfer ordinary smooth integration
and use internal hyperfinite sums, but the transferred domain, functions, and
partitions must be internal, and standard-part or summation operations then
determine what information the output retains. \emph{It is invalid to call a
proper-class surreal interval an internal hyperfinite interval merely because both
contain infinitesimals.} This identifies the extra data such a comparison would have
to carry; it does not rule one out.

\subsection{Comparison of the integration strategies}
\begin{center}\small
\begin{longtable}{@{}L{0.22\textwidth}L{0.34\textwidth}L{0.36\textwidth}@{}}
\toprule
Strategy & What it supplies & What must not be assumed\\\midrule\endhead
Polynomial simplex integration & Finite algebraic Stokes identities at arbitrary
  field scales & That it defines an integral for every semialgebraic function, or
  for forms with poles on the chain\\[0.45em]
Hahn coefficient integration & Exact $K$-valued periods and Stokes on admitted
  parameter chains & That fine continuity, or mere pole avoidance, is
  sufficient\\[0.45em]
Deformed chains and homotopy & Exact endpoint transport, integer winding, deformed
  Cauchy formulae & That all contour integrals are infinitesimally deformation
  invariant; only closed forms are\\[0.45em]
Microscopic rescaling & Contour realizations for radius-free analytic-germ residues
  & That the original coefficient germs share an ordinary radius\\[0.45em]
Determinant transformation & Coordinate-torus representation of a general isolated
  complete-intersection residue & That the untransformed integrand has the same
  torus integral, or that $\det A$ is optional\\[0.45em]
Algebraic winding and residues & Rational closed-contour pairing at arbitrary scales
  & A prior Riemann integral, or canonical open-path logarithms\\[0.45em]
Semialgebraic integration & Finite triangulated chains and Stokes formulations
  & That every version has the same scalar target\\[0.45em]
Berkovich--tropical integration & Real differential forms on enlarged analytic
  spaces & That the original point plane, or all higher valuation ranks, are
  preserved\\[0.45em]
Perfectoid cycle integration & An actual chain category with a Stokes formula
  & That residue characteristic and period-ring target can be changed for
  free\\[0.45em]
Surreal antiderivatives & Integration on distinguished extendable function classes
  & Automatic closure under arbitrary contour substitutions\\[0.45em]
Roots-of-unity averaging & Exact Laurent extraction; coefficientwise quadrature
  & That ordinary coefficient limits are valuation limits\\
\bottomrule
\end{longtable}
\end{center}

\section{A usable procedure, and research directions with exact missing steps}
\label{contours:sec:procedure}

\subsection{Choose the object before choosing the formula}
The most effective first decision is whether the input is a polynomial differential
on a nonstandard simplex, a rational function on a semialgebraic loop, a coherent
analytic family on an ordinary base, or a radius-free germ. Those choices lead
respectively to \eqref{contours:eq:simplexint},
\eqref{contours:eq:algebraicpairing}, \eqref{contours:eq:Hintegral}, and a
small-scale realization such as
\cref{contours:thm:microcircle,contours:thm:microtorus,contours:thm:generalrepresentation}. \emph{Merely declaring a function
fine-holomorphic selects none of these integrals.} A reliable workflow is the
following finite sequence of checks.
\begin{enumerate}
\item \textbf{Specify the workspace and the category of chain.} Record $\Gamma$, the
chosen monomials, and whether the parameters are ordinary real parameters,
field-valued semialgebraic parameters, or polynomial simplex data. State which
notion of interior or component is used, and which of the four rings of
\cref{contours:subsec:rings} the coefficients live in.
\item \textbf{Normalize the scale and the integrand.} In a chart $z=a+rw$, compute
the pulled-back differential \emph{including} its factor $r$, as in
\eqref{contours:eq:affineintegral}. For a meromorphic function, normalize its
denominator by a leading monomial and inspect the resulting ordinary leading
coefficient on the base contour.
\item \textbf{Prove the support certificate.} Identify the source support $S$ and the
positive perturbation set $E$, and check $S+E^*$. Positivity without well-ordering
is insufficient; a lower valuation bound without finite contribution at each
exponent is also insufficient.
\item \textbf{Choose the simplification justified by the hypotheses.} Use Stokes or
a coherent primitive for a bounding contour; \cref{contours:thm:cauchyformula} for an
evaluation; preparation and residues for finite moving poles; algebraic indices when
rational semialgebraic geometry is more informative than a coarse standard part.
\item \textbf{Resolve a boundary layer rather than assigning it a side.} When a
normalized pole reduces to the base contour, the coefficient inverse may not exist on
a common collar (\cref{contours:subsec:boundarylayer}). Move the contour by an
explicitly admissible detour, zoom into the offending scale, or use
\eqref{contours:eq:algebraicpairing}. State the resulting convention.
\item \textbf{Keep equality distinct from convergence.} A strong sum, a
coefficientwise limit, a rank-one valuation limit, and a full fine limit are
different assertions. Return the exact Hahn result with its support justification,
and add an asymptotic or numerical approximation only in an explicitly chosen
topology.
\end{enumerate}

At each step record the coefficient algebra, exponent group, parameter domain,
differential, and integral target. Those data remove most apparent paradoxes. A
proof that uses a derivative-zero fine function, an ordinary real contour, a
definable Jordan interior, and a valuation limit without comparison theorems is
generally mixing four different constructions.

\subsection{A compact decision table}
\begin{center}\small
\begin{longtable}{@{}L{0.28\textwidth}L{0.34\textwidth}L{0.30\textwidth}@{}}
\toprule
Input & Natural method & Check before applying it\\\midrule\endhead
Polynomial differential on a nonstandard simplex or polytope
  & Factorial simplex integral and \cref{contours:thm:polstokes}
  & Finite polynomial presentation; matching oriented faces.\\[0.45em]
Coherent holomorphic $F\in\HH(U)$ on an ordinary base
  & Hahn contour integral; \cref{contours:thm:cauchy,contours:thm:cauchyformula}
  & One common coefficient domain and one well-ordered support.\\[0.45em]
Uniform positive-support contour deformation
  & Taylor pullback and \cref{contours:thm:endpoints}
  & Common support throughout the deformation; nonsingular coefficient
    collar.\\[0.45em]
Rational $Q\in K(z)$ on a semialgebraic loop
  & Algebraic winding weighted by local residues
  & No pole on the loop; normalization and orientation fixed.\\[0.45em]
Radius-free one-variable germs in $\A_1$
  & Rescale to a protecting circle, \cref{contours:thm:microcircle}
  & The strict inequality $m\rho<\vH(E)$.\\[0.45em]
Radius-free coordinate-power system in $\A_n$
  & Scaled product torus, \cref{contours:thm:microtorus}
  & All $d_j\rho_j<\vH(E_j)$.\\[0.45em]
General isolated complete intersection in $\A_n$
  & Determinant transformation, \cref{contours:thm:generalrepresentation}
  & The integrand must be the \emph{transformed} one; see
    \cref{contours:subsec:scope}.\\[0.45em]
The same, with the \emph{untransformed} integrand
  & Ramified finite-cover cycle, \cref{contours:thm:untransformed}; adapted lift
    \cref{contours:thm:adaptedlift}
  & A convergent isolated $f$ and one threshold $\delta$ with
    $\kappa_*\rho<\delta$; the cycle is not a coordinate torus
    (\cref{contours:ex:wrongtorus}).\\[0.45em]
Root-of-unity sample sequence
  & Coefficientwise recovery, or a verified valuation-tail criterion
  & Never infer the topology of the limit from the sampling formula alone.\\
\bottomrule
\end{longtable}
\end{center}

\subsection{Research problems sharpened by these results}
\label{contours:subsec:problems}
\emph{A single-cycle realization for arbitrary finite Hahn zero schemes ---
posed here, and now solved.} The problem, as this article posed it, was: given
arbitrary isolated $f$ and radius-free positive-support $E$, construct one
admissible $n$-cycle on which the \emph{untransformed} integrand integrates to
\eqref{contours:eq:generalres}. A solution may require a non-diagonal analytic change
of variables, a finite resolution-type decomposition, or chains on an enlarged
analytic space, and it must prove that no omitted summation or radius assumption is
hidden in the parameterization. The finite-algebra and trace identities
\eqref{contours:eq:traceresidue} offer a robust target: any proposed cycle integral
should annihilate $(F)$ and satisfy $\Res_F(HJ_F)=\Tr(M_H)$. These are strong
compatibility tests; they do not by themselves construct the cycle. See
\cref{contours:subsec:scope} for exactly how this problem relates to
\cref{contours:thm:generalrepresentation}, which solves the transformed version.

\emph{The statement of the problem is retained above because it is the record of
what was asked.} It is answered by
\cref{contours:thm:untransformed,contours:thm:adaptedlift}, under the three
hypotheses set out in the closing paragraph of \cref{contours:subsec:scope}: a
convergent isolated leading map, one specified positive valuation lower bound
$\delta$ on the perturbation support, and a radius-free --- or, by
\cref{contours:thm:formalcontour}, formal --- numerator. The solution takes the
first of the three routes the problem allowed, in a form the problem did not
name: not a change of variables in the source but a monomial direction in the
\emph{target} combined with a ramified finite covering of the coordinate torus,
\cref{contours:lem:targetmonomial,contours:thm:ramifiedleray}. Both acceptance
tests are met, and are discussed where they are proved,
\cref{contours:subsec:normalization}: annihilation of $(F)$ is
\cref{contours:prop:idealannihilation}, and the trace test is met because the
cycle computes $\Res_F$ itself, with its $H=1$ case
$\Res_F(J_F)=\mu$ now proved directly on the cycle by
\cref{contours:thm:degreenormalization} without the finite-algebra package. What
was \emph{not} achieved, and remains open, is a canonical choice-independent
cycle and a chain-level homology comparison across protecting weights, covers and
scales; only equality of the periods is proved, by
\cref{contours:cor:scalecancel}. The formal-leading-map case, an intrinsic
invariant governing the best protecting and lifting scales, and the effectivity
question are the three further problems listed in
\cref{contours:subsec:untransformedscope}.

\emph{A boundary-layer atlas theorem.} Under explicit tame hypotheses, can rational
or selected analytic data with poles infinitesimally close to a leading contour be
covered by finitely many scale charts whose pullback integrals glue? The values for
rational closed contours are already fixed by
\eqref{contours:eq:algebraicpairing}; the issue is realizing them by a uniform chain
construction that also accommodates open paths and nonrational forms. More
generally, an \emph{admissible scale atlas} would need admissible transition
pullbacks, a common-refinement property for contour descriptions, and equality of
pulled-back integrands on each overlap; \cref{contours:lem:pullback,contours:thm:deformedstokes,contours:prop:overlap} supply local tools, not that
global theorem. If local primitives are used, their overlap differences must lie in a
controlled constant or period system, or the fine-local logarithm example
\cref{contours:ex:logarithm} destroys path invariance.

\emph{A comparison theorem with definable integration.} Identify a nontrivial common
class of functions and chains on which a chosen semialgebraic integration operator,
with its actual codomain specified, maps naturally to the Hahn functional. Polynomial
simplex forms are a necessary initial test; logarithmic forms and infinitesimal
boundary layers are more discriminating, because choices or loss of scale become
visible there. Agreement of rational periods is a beginning, not a general theorem.

\emph{An intrinsic coherent chain category.} The present construction uses explicit
ordinary parameterizations and support-controlled pullbacks. A global theory would
define admissible atlases and equivalence of chains, prove a gluing theorem, and
identify their homology independently of a chosen presentation, addressing possibly
unbounded collections of supports rather than assuming that local coherence gives
global common-support data. A related concrete question is to compare
\eqref{contours:eq:hahnforms} with the larger complex of forms whose supports are
uniform only near ordinary base points, especially on noncompact bases and for
locally finite chains.

\emph{Sectorial transcendental contour theory and higher-rank analytic chains.}
Essential singularities may require a selected summation or extension operator with
compatible monodromy: the formal Laurent expansion of $e^{1/z}$ at a nonzero
infinitesimal $z$ has descending valuations and is not strongly summable
\cite[\S13]{Analysis}, so fine analyticity alone does not identify an extension or
its periods. Separately, the microscopic torus formulas retain every exponent in
$\Gamma$, and any relationship with analytic skeleta or higher-rank valuations should
either preserve those exact periods or state precisely which quotient of the scale
information it keeps. Rank-one real-valued geometries and
positive-residue-characteristic perfectoid chains cannot be substituted unchanged;
the finite separating-torus construction provides test cases and exact invariants for
any proposed broader chain theory.

These are research directions formulated by the present analysis, \emph{not}
statements that a particular named published conjecture remains open or has been
solved.

\section{Conclusion}
\label{contours:sec:conclusion}

There is no single unqualified Jordan--Cauchy--Stokes theorem on the raw surcomplex
point class, and no justified reason to put the three classical statements under one
unspecified topology on $\No[i]$. There are, however, precise and useful
counterparts. Standard-part geometry restores an exact two-component theorem for a
thick boundary at any chosen affine scale. Definable topology restores an exact
two-component theorem with definable paths and definable connectedness, together
with integer winding. Polynomial simplex integration gives finite algebraic Stokes
identities at arbitrary field scales. Hahn differential forms give a
surcomplex-valued Stokes pairing with computable cohomology, and support-controlled
Taylor pullback extends that pairing to infinitesimally deformed chains, yielding
exact endpoint transport, Cauchy's theorem and formula, winding stability, and
deformed residue contours --- all without relying on fine-continuous paths. Positive
microscopic rescaling converts radius-free analytic germs into polynomial-coefficient
families and gives genuine small-circle and torus residue realizations; residue
transformation with its essential determinant represents arbitrary isolated
complete-intersection residues by controlled microscopic tori. Rational winding
pairings fill an important part of the gap where macroscopic coefficientwise
pullbacks fail. And the question that determinant transformation left open --- a
single cycle for the \emph{untransformed} integrand --- is answered by enlarging
the cycle rather than the theorem: a target monomial direction avoids the
ordinary discriminant, finite radial ramification supplies one fixed compact
parameter manifold, a sufficiently large infinitesimal radius protects the whole
perturbation, and the positive support lemma turns every radius-free
substitution into finite coefficient algebra. One such cycle, chosen from the
leading map and one valuation threshold alone, carries the original form to the
independently defined perturbation residue for an entire valuation ball of
perturbations whose coefficient germs need share no radius at all.

The unifying principle is constructive rather than prohibitive, and it is not to
discard infinitesimals or force every scale into one real-valued metric. It is to
specify which information each topology retains, to preserve the support algebra,
and to prove that boundary, pullback, and integration commute. Within those
conditions, nonstandard scales and infinite residues are manageable. Outside them,
changing topology or summation method is a new mathematical step that must be stated
and justified.

\appendix

\section{Support certificates and logical dependencies}
\label{contours:app:support}

The following table records the containing supports used in the proofs. $S$ is a
source support, $T$ a second source support, $E$ a positive well-ordered error
support, $E_a$ the union of the coordinate supports of a translation point, and
$P=\{\rho_1,\ldots,\rho_n\}$ a finite positive rescaling alphabet.

\begin{center}\small
\begin{longtable}{@{}L{0.46\textwidth}L{0.46\textwidth}@{}}
\toprule
Operation & Containing support or finite mechanism\\\midrule\endhead
Coefficientwise $d$, ordinary integration, differentiation, residue extraction, a
  chosen linear splitting, or the radial homotopy operator
  \eqref{contours:eq:radial} & $S$; no enlargement whatever.\\[0.45em]
Product of two Hahn families & $S+T$, with finite convolution at each
  exponent.\\[0.45em]
Infinitesimal Taylor evaluation \eqref{contours:eq:coherentevaluation} or deformed
  pullback \eqref{contours:eq:formpullback} & $S+E^*$.\\[0.45em]
Deformed homotopy operator \eqref{contours:eq:homotopyoperator}
  & The same $S+E^*$; ordinary integration in the homotopy parameter adds no
    exponents.\\[0.45em]
Endpoint correction \eqref{contours:eq:endpointT}
  & $S+(\supp\varepsilon)^*$, with at least one displacement factor in every
    term.\\[0.45em]
Positive monomial rescaling \eqref{contours:eq:rescaling}
  & $S+P^*$; a fixed output exponent has a \emph{polynomial} coefficient.\\[0.45em]
Rescaling after infinitesimal translation \eqref{contours:eq:translatedrescale}
  & $S+(E_a\cup P)^*$.\\[0.45em]
Normalized one-variable denominator, \cref{contours:lem:protect}
  & Positive error because $\vH(E)-m\rho>0$; inverse valid on an ordinary base
    annulus avoiding $w=0$.\\[0.45em]
Normalized coordinate-power denominator tuple, \eqref{contours:eq:protecting}
  & Component $j$ has error support $(\supp E_j-d_j\rho_j)+P^*$, positive by the
    protecting inequality; inverse valid on a common polyannulus.\\[0.45em]
Perturbation residue after integration,
  \eqref{contours:eq:toruscoeff}--\eqref{contours:eq:generalrepresentation}
  & $S+E^*$; before integration the numerator and volume factors may introduce a
    fixed shift, but a shift of a well-ordered support is still well
    ordered.\\[0.45em]
Substitution of a positive-support Hahn map into a radius-free or formal germ,
  \cref{contours:lem:radiusfreesub}
  & $S+P_W^*$, where $P_W$ is the union of the coordinate supports of the
    substituted map; each output coefficient is a finite sum of products of
    ordinary holomorphic functions on the one fixed $\mathsf X$.\\[0.45em]
Radial evaluation $s=t^\rho$ of an ordinary family vanishing at $s=0$,
  \cref{contours:cor:radialeval}
  & $\N_{>0}\rho$; a fixed output exponent sees finitely many radial
    degrees.\\[0.45em]
Untransformed pullback on the protecting cycle,
  \cref{contours:prop:admissiblepullback}
  & $S-\kappa_\Sigma\rho+\mathcal M_0$ with
    $P_0=\{\rho\}\cup\bigcup_j(S_E-\kappa_j\rho)$ positive by
    $\kappa_*\rho<\delta$; asserted for the \emph{fully expanded} family, before
    any regrouping, via the accounting
    \eqref{contours:eq:exponentaccounting}.\\[0.45em]
Rooted-tree solution of a positive implicit equation,
  \cref{contours:lem:positiveimplicit}
  & $W^*\setminus\{0\}$; at a fixed weight the vertex count is bounded, so
    finitely many trees contribute. \emph{Not} a limit of iterates.\\[0.45em]
Equation-adapted correction, \cref{contours:thm:adaptedlift}
  & $L\rho+(\mathcal M\setminus\{0\})$ and also
    $S_E-b_{\mathrm{pole}}\rho+\mathcal M$, with
    $P=\{\rho\}\cup(S_E-(L+b_{\mathrm{pole}})\rho)$; the second bound comes from
    the nullary vertex every finite tree has.\\[0.45em]
Straight homotopy of the two untransformed cycles,
  \cref{contours:lem:straighthomotopy}
  & $S-\kappa_\Sigma\rho+\mathcal M$, uniformly on the ordinary compact cylinder
    $\mathsf T\times[0,1]$; at each exact exponent the dependence on the ordinary
    parameter is polynomial.\\[0.45em]
Untransformed period after integration,
  \eqref{contours:eq:untransformed}
  & $S+E^*$; every auxiliary scale in
    \eqref{contours:eq:untransformedsupport} cancels after angular integration,
    by \cref{contours:cor:scalecancel}.\\[0.45em]
Algebraic rational contour \eqref{contours:eq:algebraicpairing}
  & Finite pole set and integer algebraic winding; no infinite sum or limit.\\[0.45em]
Polynomial simplex integral \eqref{contours:eq:simplexint}
  & Finite rational factorial identities in finitely many coefficients.\\[0.45em]
Coefficientwise sampling limit \eqref{contours:eq:coefflimit}
  & A fixed containing support $S$ and ordinary convergence in every coefficient; no
    valuation-convergence inference.\\
\bottomrule
\end{longtable}
\end{center}

For clarity, ``only finitely many dependencies at an exponent'' does not mean ``only
finitely many exponents below it''; the latter fails below $\omega$ in
$\Q+\Q\omega$. The proofs use the former, which follows from finite word
decompositions. This distinction is essential in
\cref{contours:lem:residuetransform} and in the passage from infinitely many germs to
one finite ordinary-parameter calculation.

Several hypotheses have deliberately \emph{not} been weakened. The curve in the
topological Jordan theorem may be arbitrary, but integration uses piecewise smooth or
polynomial parameterizations. The bounded Jordan domain in the Cauchy formula has
closure inside the common coefficient domain. Positive perturbation supports must be
well ordered, not just positive. Homotopies share one support on their whole ordinary
compact parameter space. \Cref{contours:thm:microtorus} assumes coordinate-power
leading equations. Rational closed-period integration is a definition, compared to
analytic integration only on an explicitly established overlap.
\Cref{contours:thm:untransformed} assumes a convergent isolated leading map and
\emph{one specified} positive lower bound $\delta$ on the perturbation support;
\cref{contours:thm:formalcontour} removes convergence from the perturbation and
the numerator, but not from the leading map, and
\cref{contours:rem:positivityneeded} records why positivity of the support cannot
simply be dropped.

The dependency structure is explicit. \Cref{contours:lem:support} defines the Hahn
operations. Ordinary Stokes proves \cref{contours:thm:hahnstokes}. Taylor pullback
and finite-jet identities prove \cref{contours:thm:deformedstokes}. Cauchy follows
from closedness and ordinary coefficient Cauchy. Independently, positive rescaling
and Laurent coefficient extraction realize the perturbation residue, and classical
finite-family residue transformation extends the realization to general isolated
complete intersections. Independently again, the untransformed realization runs
ordinary finite-map theory, a target monomial direction, radial ramification and
a removable-singularity extension entirely in an \emph{ordinary} complex radial
variable (\cref{contours:lem:targetmonomial,contours:thm:ramifiedleray,contours:lem:constantperiod,contours:lem:radialbounds}),
and only afterwards substitutes $s=t^\rho$; the Hahn half of that argument uses
\cref{contours:lem:support,contours:lem:radiusfreesub,contours:lem:positiveimplicit}
and ordinary Stokes on a compact cylinder, and nothing else. \emph{Reversing that
order of operations --- substituting the Hahn monomial before the covering and
the removable-singularity step --- would be exactly the interchange of analytic
convergence with Hahn summability that this article refuses throughout.} The
rational pairing is defined independently from
semialgebraic winding and compared with Hahn integration only afterward. \emph{No
Jordan statement in the fine topology, and no unproved general fine integration rule,
enters this chain of arguments.} The new proofs do not depend on the Noetherian or
Nullstellensatz claims of \cite{Geometry} except where that manuscript's finite-zero
and residue results are expressly invoked --- in
\eqref{contours:eq:traceresidue} and, conditionally and with that dependence stated,
in \cref{contours:thm:septorus}. In particular \cref{contours:thm:polstokes,contours:thm:hahnstokes,contours:thm:cohomology,contours:thm:homotopy,contours:thm:generalrepresentation,contours:thm:untransformed,contours:thm:adaptedlift,contours:thm:degreenormalization} and the sampling comparisons stand independently
of those deeper analytic-algebra claims; \cref{contours:thm:degreenormalization}
in particular proves $\Res_F(J_F)=\mu$ without any Hahn flatness or
Nullstellensatz input, and the only place where the new sections rely on
\cite{Geometry} at all is the second acceptance test discussed in
\cref{contours:subsec:normalization}, which is a statement about the functional
and not about the cycle.

\section{Why finite-parameter testing proves residue transformation}
\label{contours:app:transform}

This appendix supplies the technical point used in
\cref{contours:lem:residuetransform} and in \cref{contours:thm:septorus}, without
assuming an infinite-parameter contour for the original germ family. It is stated for
a transformation matrix $\mathsf C$ arising from a linear splitting, which is the
harder case; the ordinary analytic matrix $A$ of
\eqref{contours:eq:transformmatrix} is a special case in which the convergence
discussion of \cref{contours:appsub:formal} is unnecessary.

\subsection{Finite dependency}
In the supplied normal-form construction an ordinary linear splitting for
$\C\{z\}/(f)$ is extended coefficientwise, and the perturbed inverse is a Neumann
series in operators carrying positive perturbation support. Therefore each
coefficient of $M_i$, of $Q_i$, and of the chosen $C_{ij}$ is built from a finite sum
of positive-support words. The residue expansion
\eqref{contours:eq:generalres} has the same property: for a numerator with support
$S$, its coefficient at $\lambda$ uses only decompositions
\[
  \lambda=\sigma+e_1+\cdots+e_q,\qquad\sigma\in S,\ e_j\in E,
\]
of which there are finitely many, including finitely many possible lengths, by
\cref{contours:lem:support}. \emph{This does not say that only finitely many
exponents lie below $\lambda$.}

Fix a coefficient to be checked in the transformation identity. Retain all ordinary
source germs and perturbation germs occurring in its finite dependency sets on both
sides, treating numerator source coefficients separately by linearity. Replace the
finitely many retained positive perturbation monomials by independent parameters
$u_1,\ldots,u_b$. The normal-form formulae now define formal power series in $u$ with
ordinary analytic germ coefficients, and their formal identities include
$Q(u)=\mathsf C(u)F(u)$ with $Q_i(0,z)=z_i^\mu$. All omitted input terms are
irrelevant to the selected coefficient by construction.

\subsection{Why a formal matrix causes no convergence gap}
\label{contours:appsub:formal}
The auxiliary ordinary linear splitting need not be continuous in an analytic norm,
so $\mathsf C(u)$ cannot be asserted convergent as an analytic parameter series. Let
$N$ exceed every total parameter degree needed for the selected coefficient and let
$\mathsf C_{\le N}(u)$ be its finite truncation; its finitely many coefficient germs
\emph{do} have a common ordinary representative neighborhood. Define the actual
finite analytic family
\[
  \widetilde Q(u,z)=\mathsf C_{\le N}(u,z)\,F(u,z),
\]
which has ordinary leading tuple $(z_1^\mu,\ldots,z_n^\mu)$ and agrees with the
formal $Q(u,z)$ modulo $(u)^{N+1}$. The tuple $F(u,z)$ has leading isolated complete
intersection $f$, so after choosing ordinary small representatives both tuples have
finite clusters near zero.

The ordinary residue transformation rule, applied to these finite families and their
clusters, gives
\[
  \res^{\mathrm{cluster}}_{F(u)}(H)
  =\res^{\mathrm{cluster}}_{\widetilde Q(u)}
   \bigl(H\det\mathsf C_{\le N}(u)\bigr),
\]
the family form of classical residue transformation, which follows from the local
transformation law and additivity over the finitely many points \cite{CDS,Tsikh}.
Expanding the reciprocal denominators on ordinary protecting residue cycles
identifies their Taylor series with the geometric residue expansions, and a Taylor
coefficient through total parameter degree $N$ uses only denominator and numerator
coefficients through that degree. Consequently the last display proves the formal
transformation law through degree $N$ for the original $Q$ and $\mathsf C$.

Finally substitute $u_j=t^{e_j}$ and extract the chosen Hahn coefficient. Several
parameter monomials may contribute to the same exponent, but only finitely many do,
and their finite sum gives the desired coefficient equality. Repeating the argument
at each exponent establishes the transformation identity. The auxiliary ordinary
cycles may change from exponent to exponent; \emph{the final microscopic separating
torus is constructed independently}, by \cref{contours:thm:rescaling} and the radius
inequalities \eqref{contours:eq:protecting} or
\eqref{contours:eq:torusradius}.

\section{What is claimed, what is imported, and what is not proved}
\label{contours:app:status}

\subsection{Consolidated status statement}
Established results, statements reported in the two supplied manuscripts, and
deductions proved here are distinguished throughout. The two supplied manuscripts
\cite{Analysis,Geometry} are treated as research manuscripts, not as independently
refereed references; they are not modified or redistributed, and their source
filenames are recorded in the bibliography so that the direction of dependence can be
checked. The literature search was conducted on September 21, 2026 and was targeted
rather than exhaustive. \emph{This is not a claim of a unique accepted foundation for
surcomplex analysis, an exhaustive priority or novelty determination, or a formally
verified proof.} No named published conjecture is claimed to be solved. Symbolic
checks validate the displayed finite examples --- exact expressions, Taylor
coefficients, matrix identities, residue traces, finite roots-of-unity averages,
polynomial Stokes on monomial forms, and endpoint transport on polynomial
deformations --- and are \emph{not} machine verification of the general theorems: they
do not implement arbitrary Hahn fields, enumerate arbitrary ordered supports, decide
arbitrary surreal comparisons, or prove semialgebraic Jordan separation. All sums and
computational workspaces are set-sized; no proper-class summation is used. The
priority position is that the value of this package is its complete set of hypotheses;
\emph{no priority claim is based merely on not finding an identical formulation in the
literature}, and the constructions proved here are proposed extensions of the supplied
presentation, not a claim that their underlying coefficientwise or residue-theoretic
ideas are absent from the literature. This article is not an independent audit of
every theorem in either supplied manuscript.

\subsection{Per-theorem non-claims, preserved}
The following specific non-claims are part of the results they accompany and are
repeated here for reference.
\begin{itemize}
\item \emph{On the sources.} Neither supplied manuscript establishes a Jordan theorem
for arbitrary fine-topological curves, nor generalized Stokes for arbitrary
fine-smooth functions and arbitrary surcomplex domains. We do not silently attribute
either assertion to them.
\item \emph{On \cref{contours:prop:fullfine}.} This is not a theorem that a Hahn field
with its own valuation topology is discrete; the radius supplied by the cut may lie
outside the workspace.
\item \emph{On \cref{contours:rem:threetopologies}.} A claim about convergence or
connectedness in one of the three topologies is not automatically a claim in either of
the others.
\item \emph{On \cref{contours:thm:stdpartjordan}.} The separating object is the
thickened curve, not the ordinary point set; neither component is called
fine-connected; two infinitesimally separated points in one boundary monad cannot be
assigned different sides, and the theorem cannot distinguish two zeros lying in one
standard-part fibre; an unbounded exterior in the internal ordered-field sense has not
been obtained.
\item \emph{On \cref{contours:thm:sajordan}.} This is an application of semialgebraic
topology, not a new surreal theorem; the technical finite-topology input is imported,
not replaced by an informal appeal to transfer of arbitrary continuous maps; the
components are not connected in the unrestricted intrinsic topology; the o-minimal
generalization refers to that framework's differentiability and definability
hypotheses, not to arbitrary class functions or arbitrary fine-topological subsets;
definable results do not supply ordinary singular homology of a proper class.
\item \emph{On winding.} The definable degree is an ordinary integer, not a surreal
number of turns; it does not by itself define the integral of a general differential
form; the primitive of a semialgebraic function need not be semialgebraic; the
availability of a finite algebraic prescription does not make arbitrary surreal inputs
computable.
\item \emph{On \cref{contours:prop:mesh}.} A polynomial antiderivative still gives
$\int_0^1x\dd x=1/2$; the obstruction is to a particular unrestricted finite-mesh
definition, not to all integration. Hyperfinite partitions require an additional
nonstandard universe and transfer structure, supplied neither by the ordered-field
notation nor by the surreal normal form.
\item \emph{On \cref{contours:thm:polstokes}.} The proof uses classical Stokes to
certify a finite algebraic identity; it does not transfer a statement quantifying over
arbitrary integrable functions. It does not define an integral for every semialgebraic
function, nor for a rational form whose denominator vanishes on the chain.
\item \emph{On \cref{contours:thm:cohomology}.} For infinite-dimensional ordinary
cohomology, $V((t^\Gamma))$ must not be silently replaced by $K\otimes_\C V$; and the
global uniform-support assignment is not claimed to be a sheaf for arbitrary infinite
covers.
\item \emph{On \cref{contours:def:deformation,contours:lem:pullback}.} $\Phi$ describes
an admissible parameterized chain, not a fine-continuous map from the ordinary
parameter manifold; the jet prolongation is the \emph{specified} extension, with no
claim of ordinary Taylor convergence or of a canonical prolongation of every smooth
function on every fine domain.
\item \emph{On \cref{contours:thm:homotopy}.} It is incorrect to claim that all contour
integrals are infinitesimally deformation invariant; closedness is the controlling
condition.
\item \emph{On \cref{contours:prop:ml}.} Ordered estimates do not imply convergence of
set-indexed nets in the full fine topology.
\item \emph{On \cref{contours:thm:rescaling}.} The conclusion is a new coherent
representation in an infinitesimal chart, not ordinary convergence of the original
germs on a uniform disk; it neither establishes nor requires all-scale coherence.
\item \emph{On \cref{contours:thm:localresidue}.} An arbitrary Hahn family of
meromorphic germs is larger than a quotient of two elements of $\A_1$, and is not
claimed to admit the microscopic contour construction.
\item \emph{On \cref{contours:thm:generalrepresentation}.} The untransformed form is
not claimed to have the same torus integral; $\det A$ is not an optional Jacobian
factor, since $AF$ can introduce extra zeros; and $F$ and $G$ are not claimed to have
identical zero schemes. \Cref{contours:subsec:scope} records that one of the merged
sources judged the general realization unproved, and that the alternative route
\cref{contours:thm:septorus} is conditional on \cite{Geometry}.
\item \emph{On \cref{contours:rem:toruslimits}.} Formula
\eqref{contours:eq:septorus} does not assert regularity of the original quotient on
the chosen torus; omitting the determinant is generally wrong; a real $n$-torus is not
identified with the full boundary of a complex polydisk.
\item \emph{On \cref{contours:def:algpairing}.} Equation
\eqref{contours:eq:algebraicpairing} is a definition, not a claim that a
finite-partition Riemann net converges to this value; it does not determine open-path
logarithms and is not a definition on all of $\HH(U)$ or $\A_n$.
\item \emph{On \cref{contours:thm:sampling}.} Coefficientwise recovery does not make
the Hahn integral a fine limit, and in general does not make it a valuation limit
either.
\item \emph{On \cref{contours:thm:untransformed} --- hypotheses.} An isolated
\emph{convergent} ordinary leading map and a positive-support perturbation are
required. An arbitrary Hahn tuple with no such leading map, a non-isolated
leading zero, and a support not bounded below by a positive element are all out
of scope. Extending to a \emph{formal} isolated leading map is a different open
problem, because the proof needs a bounded analytic finite covering constructed
before any Hahn substitution; \cref{contours:thm:formalcontour} removes
convergence only from $E$ and $H$.
\item \emph{On \cref{contours:thm:untransformed} --- what independence means.}
Period independence is proved (\cref{contours:cor:scalecancel}); no canonical
choice-independent contour is produced. The parameter space may be disconnected,
and it is not claimed that every admissible cycle is homologous to the chosen
one. A chain-level homology comparison across weights, covers and scales is
listed as open.
\item \emph{On \cref{contours:thm:adaptedlift,contours:cor:protectedsheets}.} The
adapted parameterization gives $\mu$ distinct \emph{regular} points over each
protected target point. It does not assert that every zero of $F$ is simple, nor
that the deformed zero algebra is reduced, and \emph{its image is not identified
with the entire class of surcomplex solutions of $|F_j|=t^{\kappa_j\rho}$} --- no
such identification is needed for any integral proved here.
\item \emph{On \cref{contours:prop:idealannihilation}.} Nonzero descent of
$\Res_F$ to $\A_n/(F)$ is established; a perfect-pairing or flatness theorem is
not inferred from it. \Cref{contours:thm:degreenormalization} is not called a
conservation law for local algebra length, which would additionally require the
finite-flatness identification.
\item \emph{On \cref{contours:thm:untransformedhomotopy}.} The Hahn
parameterizations are not asserted continuous for the fine surreal topology, and
the homotopy is not a fine-topological homotopy; the parameter $u$ is an ordinary
real parameter. All supports, covers and parameter manifolds are set-sized.
\item \emph{On \cref{contours:subsec:construction}.} A finite coefficient
dependency set supplies no decision procedure for membership in an arbitrary
well-ordered subset of a general ordered group; no uniform complexity bound is
claimed; arbitrary convergent germs need not have finite effective descriptions;
and the scale estimates \eqref{contours:eq:Lrho} are transparent but can be far
from sharp.
\item \emph{On \cref{contours:rem:leraychecks}.} The $632$ exact symbolic checks
are finite examples only. They do not verify arbitrary-support Hahn summability,
parameterized removable singularities, the rooted-tree construction of
\cref{contours:lem:positiveimplicit}, or novelty; the source's README and its
recorded output both say so.
\item \emph{On \cref{contours:subsec:untransformedscope}.} M\"uller and
Strohmaier's Hahn-meromorphic category \cite{MuellerStrohmaier} is explicitly
\emph{not} identified with the coefficientwise ring $\A_n$ of radius-free germs,
and no comparison theorem between them is asserted. The targeted literature
search of September 21, 2026 found no published theorem with the
arbitrary-support, radius-free, uniform-protecting-cycle hypotheses proved here,
and \emph{that is not a priority certification}. The publisher endpoint for the
Neumann reference \cite{Neumann} could not be fetched during that search and its
details were corroborated indirectly.
\end{itemize}

\section{Notation}
\label{contours:app:notation}
\begin{center}\small
\begin{longtable}{@{}L{0.26\textwidth}L{0.66\textwidth}@{}}
\toprule
Symbol & Meaning\\\midrule\endhead
$\SC=\No[i]$ & The full surcomplex class field.\\[0.35em]
$\Gamma,\ R_\Gamma,\ K$ & A set-sized divisible ordered subgroup of $\No$, and the
  real and complex Hahn workspaces $\R((t^\Gamma))$ and $\C((t^\Gamma))$, with
  $t^\gamma=\omega^{-\gamma}$.\\[0.35em]
$v,\ \vH$ & Least exponent of a field element; least exponent of a Hahn
  coefficient-family datum.\\[0.35em]
$K^\circ,\ \m_K,\ \m,\ \st$ & Finite elements; infinitesimals of the workspace; the
  class of all infinitesimals of $\SC$; the standard-part (reduction) map
  $K^\circ\to\C$.\\[0.35em]
$U\sh$ & $\st^{-1}(U)=\{c+\varepsilon:c\in U,\ \varepsilon\in\m_K\}$.\\[0.35em]
Standard-part topology & The initial topology of $\st$; called the residue, reduction,
  or shadow topology elsewhere.\\[0.35em]
$E^*$ & The additive monoid of finite words in a well-ordered positive support $E$,
  including $0$.\\[0.35em]
$\HH(U)$ & $\OO(U)((t^\Gamma))$: one common ordinary coefficient domain.\\[0.35em]
$\A_n$ & $\C\{z_1,\ldots,z_n\}((t^\Gamma))$: radius-free coefficient germs.\\[0.35em]
$\C[w]((t^\Gamma))$ & Polynomial-coefficient ring; the target of
  \cref{contours:thm:rescaling}.\\[0.35em]
$\Omega_H^q(M)$ & $\Omega^q(M;\C)((t^\Gamma))$: globally supported Hahn series of
  ordinary smooth complex $q$-forms.\\[0.35em]
$\Phi=\phi_0+\eta$ & Ordinary base parameterization plus a uniformly supported
  infinitesimal deformation, \cref{contours:def:deformation}.\\[0.35em]
$\HInt,\ \HCoint$ & Hahn coefficientwise integration and its closed-contour
  form.\\[0.35em]
$\PInt,\ I_{\mathrm{alg}}$ & Polynomial-simplex integration
  \eqref{contours:eq:simplexint}; the algebraic rational contour pairing
  \eqref{contours:eq:algebraicpairing}.\\[0.35em]
$T_F(c,\varepsilon)$ & Microscopic endpoint Taylor correction
  \eqref{contours:eq:endpointT}.\\[0.35em]
$\wind,\ \wind_{\mathrm{sa}}$ & Ordinary complex winding; definable/algebraic integer
  winding.\\[0.35em]
$\res_a,\ \Res_F$ & Ordinary local residue (Laurent or Grothendieck); the strong Hahn
  perturbation residue \eqref{contours:eq:generalres}.\\[0.35em]
$m,\ \mu,\ e_a$ & Order of a one-variable germ; global colength; local multiplicity at
  an actual displaced zero, $\sum_ae_a=\mu$.\\[0.35em]
$\rho_j,\ r_j=t^{\rho_j}$ & Rescaling exponents and the corresponding microscopic
  radii.\\[0.35em]
$d_j$ & Coordinate powers in a coordinate-power leading system
  \eqref{contours:eq:powerleading}.\\[0.35em]
$J_F,\ M_H,\ B_F$ & Jacobian determinant; multiplication operator; the finite quotient
  algebra $\A_n/(F)$.\\[0.35em]
$\HH(\mathsf X)$ & $\OO(\mathsf X)((t^\Gamma))$ for one fixed ordinary complex
  manifold $\mathsf X$, possibly disconnected; $\HH(\mathsf X)_+$ for strictly
  positive support, \eqref{contours:eq:HXring}.\\[0.35em]
$\mathsf A,\ p,\ \mathsf X,\ \mathsf T$ & A polyannulus around the ordinary
  torus; a finite unramified covering $p:\mathsf X\to\mathsf A$ of degree $\mu$;
  its total space; and $\mathsf T=p^{-1}\bigl((S^1)^n\bigr)$, the compact
  oriented ordinary parameter manifold of
  \cref{contours:thm:ramifiedleray}.\\[0.35em]
$\kappa_j,\ \kappa_*,\ \kappa_\Sigma$ & Target radial weights with
  $f_j(Z(q,s))=s^{\kappa_j}p_j(q)$, their maximum and their sum. \emph{Written
  $m_j$ and $M$ in the source manuscript}; renamed because $m$ is the order of a
  one-variable germ.\\[0.35em]
$s,\ \rho,\ \delta$ & The \emph{ordinary} complex radial variable of
  \cref{contours:subsec:ramified}; the protecting exponent with
  $\kappa_*\rho<\delta$; the specified positive lower bound on
  $\supp E$.\\[0.35em]
$Z,\ Z_0,\ \Phi,\ \Upsilon$ & The ordinary ramified family; the protecting Hahn
  cycle $\sum_{\ell\ge1}t^{\ell\rho}Z_\ell$; the equation-adapted lift
  $Z_0+\Upsilon$; and its correction.\\[0.35em]
$b_{\mathrm{pole}},\ c_{\mathrm{sep}},\ L$ & Inverse-Jacobian pole order and
  sheet-separation degree of \cref{contours:lem:radialbounds}; the lift exponent
  $L=\max\{b_{\mathrm{pole}},c_{\mathrm{sep}},\kappa_*\}+1$.\\[0.35em]
$S_E,\ P_0,\ \mathcal M_0,\ P,\ \mathcal M$ & $\supp E=\bigcup_j\supp E_j$; the
  protecting and lifting generating sets \eqref{contours:eq:Pzero} and
  \eqref{contours:eq:Pfull}, and the monoids they generate.\\
\bottomrule
\end{longtable}
\end{center}

\begin{thebibliography}{99}
\addcontentsline{toc}{section}{References}

\bibitem{Analysis}
\emph{Surcomplex Analysis: Infinitesimal Calculus, Hahn-Coherent Holomorphy, and
Contour Theory}. Supplied research manuscript, September 21, 2026. Source file:
\texttt{surcomplex\_analysis(3).tex}. Used for its fine-topology obstructions,
coherent contour functionals, preparation and moving-pole residues, endpoint
corrections, affine scales, global fine logarithms, and all-scale rigidity results;
not treated as a published priority certification.

\bibitem{Leray}
\emph{Untransformed Leray Cycles for Radius-Free Hahn Analytic Residues}.
Supplied research manuscript, September 21, 2026, written against this report at
repository commit \texttt{39f2be66}. Its exact symbolic check script and that
script's recorded run are kept in this directory under \texttt{code/} and
\texttt{data/}, both carrying the file-name stem
\texttt{05-untransformed-\allowbreak leray-\allowbreak cycles}.
It is the source of
\cref{contours:sec:untransformed,contours:sec:adapted}, and is not treated as a
published priority certification.

\bibitem{FiniteDeformations}
\emph{Finite Hahn Deformations: Division, Conservation of Multiplicity, and
Residue Duality}. Companion merged report in this repository at
\texttt{docs/surcomplex/finite-deformations}. Cited for the Hahn local-cluster
residue defined by a \emph{series} rather than by a contour, for residue duality
and conservation of multiplicity, and for the support-preserving lift of the
ordinary root cover with unchanged monodromy. Its results are cited, not
reproved here.

\bibitem{Geometry}
\emph{Infinitesimal Analytic Geometry over the Surcomplex Numbers: Radius-Free Hahn
Germs, a Nullstellensatz, and Conservation of Singular Zeros}. Supplied research
manuscript, September 21, 2026. Source file:
\texttt{surcomplex\_analytic\_geometry.tex}. Used for the radius-free germ ring
$\A_n$, support-controlled evaluation, translation and division, the scaled Taylor
polynomial lemma, the finite complete-intersection deformation theorem, and the
perturbation residue.

\bibitem{Ahlfors}
L. V. Ahlfors, \emph{Complex Analysis}, 3rd ed., McGraw--Hill, 1979. Chapter 4
contains the classical contour, Cauchy, and residue theory.

\bibitem{BCR}
J. Bochnak, M. Coste, and M.-F. Roy, \emph{Real Algebraic Geometry}, Ergebnisse der
Mathematik und ihrer Grenzgebiete (3), vol.~36, Springer, 1998.
\href{https://doi.org/10.1007/978-3-662-03718-8}{doi:10.1007/978-3-662-03718-8}.
Semialgebraic triangulation and topology over real closed fields.

\bibitem{BO}
A. Berarducci and M. Otero, ``Intersection theory for o-minimal manifolds,''
\emph{Annals of Pure and Applied Logic} \textbf{107} (2001), 87--119.
\href{https://doi.org/10.1016/S0168-0072(00)00027-0}{doi:10.1016/S0168-0072(00)00027-0}.

\bibitem{BOhomology}
A. Berarducci and M. Otero, ``O-minimal fundamental group, homology and manifolds,''
\emph{Journal of the London Mathematical Society} (2) \textbf{65} (2002), 257--270.
\href{https://doi.org/10.1112/S0024610701003015}{doi:10.1112/S0024610701003015}.

\bibitem{BOmeasure}
A. Berarducci and M. Otero, ``An additive measure in o-minimal expansions of fields,''
\emph{Quarterly Journal of Mathematics} \textbf{55} (2004), 411--419.
\href{https://doi.org/10.1093/qmath/hah010}{doi:10.1093/qmath/hah010}.

\bibitem{BR}
Robert Rumely and Matthew Baker,
``Analysis and dynamics on the Berkovich projective line,'' 2004.
\href{https://arxiv.org/abs/math/0407433}{arXiv:math/0407433}.
Expanded lecture notes; used for the tree and unique-path geometry, not as a
contour-integral comparison theorem.

\bibitem{CDS}
E. Cattani, A. Dickenstein, and B. Sturmfels, ``Computing multidimensional residues,''
in \emph{Algorithms in Algebraic Geometry and Applications}, Progress in Mathematics,
vol.~143, Birkh\"auser, 1996, 135--164.
\href{https://arxiv.org/abs/alg-geom/9404011}{arXiv:alg-geom/9404011};
\href{https://doi.org/10.1007/978-3-0348-9104-2_8}{doi:10.1007/978-3-0348-9104-2\_8}.
\S0 records local residues, the transformation law with its determinant, coefficient
extraction, and univariate separating denominators.

\bibitem{CE}
O. Costin and P. Ehrlich, ``Integration on the surreals,'' \emph{Advances in
Mathematics} \textbf{452} (2024), article 109823.
\href{https://doi.org/10.1016/j.aim.2024.109823}{doi:10.1016/j.aim.2024.109823};
\href{https://arxiv.org/abs/2208.14331}{arXiv:2208.14331}.

\bibitem{Cherry}
William Cherry, ``Lectures on Non-Archimedean Function Theory,'' 2009.
\href{https://arxiv.org/abs/0909.4509}{arXiv:0909.4509}.
Section 1.3 gives the weighted Schnirelman normalization, its limit condition,
and the caution about dependence on the chosen radius element.

\bibitem{CLD}
A. Chambert-Loir and A. Ducros, \emph{Formes diff\'erentielles r\'eelles et courants
sur les espaces de Berkovich},
\href{https://arxiv.org/abs/1204.6277}{arXiv:1204.6277}, expanded version~2, July 25,
2025. Real forms, tropical spaces, integration, and Stokes formulas; Proposition
5.11.1 is the calibrated tropical Stokes formula cited here.

\bibitem{DHL}
A. Ducros, E. Hrushovski, and F. Loeser, ``Non-Archimedean integrals as limits of
complex integrals,'' \emph{Duke Mathematical Journal} \textbf{172} (2023), 313--386.
\href{https://doi.org/10.1215/00127094-2022-0052}{doi:10.1215/00127094-2022-0052};
\href{https://arxiv.org/abs/1912.09162}{arXiv:1912.09162}.

\bibitem{Diamond}
J. Diamond, ``$p$-adic Line Integrals and Cauchy's Theorems,'' preprint, 2016.
\href{https://arxiv.org/abs/1605.06484}{arXiv:1605.06484}.

\bibitem{Eisermann}
M. Eisermann, ``The Fundamental Theorem of Algebra Made Effective: An Elementary
Real-Algebraic Proof via Sturm Chains,'' \emph{American Mathematical Monthly}
\textbf{119} (2012), 715--752.
\href{https://doi.org/10.4169/amer.math.monthly.119.09.715}{doi:10.4169/amer.math.monthly.119.09.715};
\href{https://arxiv.org/abs/0808.0097}{arXiv:0808.0097}. Especially Theorem~1.2 on
algebraic winding and Theorem~1.5 on root counting.

\bibitem{Ettayb}
J. Ettayb, ``Some integrals for groups of bounded linear operators on
finite-dimensional non-Archimedean Banach spaces,'' \emph{Buletinul Academiei de
\c Stiin\c te a Republicii Moldova. Matematica} (2022), no.~3, 3--14.
\href{https://doi.org/10.56415/basm.y2022.i3.p3}{doi:10.56415/basm.y2022.i3.p3}. A
reference for contemporary use of Volkenborn and Schnirelman integration; no
convergence theorem is imported from it.

\bibitem{GH}
V. Guillemin and P. J. Haine, \emph{Differential Forms}, World Scientific, 2019.
\href{https://doi.org/10.1142/11058}{doi:10.1142/11058}. Theorem 4.6.1 is the ordinary
Stokes theorem used here.

\bibitem{GrauertRemmert}
H. Grauert and R. Remmert, \emph{Coherent Analytic Sheaves}, Grundlehren der
mathematischen Wissenschaften, vol.~265, Springer, 1984.
\href{https://doi.org/10.1007/978-3-642-69582-7}{doi:10.1007/978-3-642-69582-7}.
Chapters~II--III: local Weierstrass theory and finite holomorphic maps, the
ordinary finite-map input to \cref{contours:thm:ramifiedleray}.

\bibitem{GriffithsHarris}
P. Griffiths and J. Harris, \emph{Principles of Algebraic Geometry}, Wiley, 1978;
Wiley Classics Library edition, 1994.
\href{https://doi.org/10.1002/9781118032527}{doi:10.1002/9781118032527}.
Chapter~5, ``Residues,'' pp.~647--732: ordinary Leray cycles, the Grothendieck
residue, and its finite-jet and annihilation properties. \emph{Not to be confused
with \cite{GH}}, which is Guillemin--Haine.

\bibitem{Higman}
G. Higman, ``Ordering by divisibility in abstract algebras,'' \emph{Proceedings
of the London Mathematical Society}, series~3, \textbf{2} (1952), 326--336.
\href{https://doi.org/10.1112/plms/s3-2.1.326}{doi:10.1112/plms/s3-2.1.326}.
The word lemma giving the second proof of \cref{contours:lem:support} recorded in
\cref{contours:subsec:hahnsub}.

\bibitem{Kaiser}
Tobias Kaiser,
``Lebesgue measure theory and integration theory on non-archimedean real closed
fields with archimedean value group.''
\href{https://arxiv.org/abs/1409.2241}{arXiv:1409.2241}.
Used only with its stated value-group, valuation-section, and output-ring
qualifications.

\bibitem{Lee}
J. M. Lee, \emph{Introduction to Smooth Manifolds}, 2nd ed., Graduate Texts in
Mathematics, vol.~218, Springer, 2013.
\href{https://doi.org/10.1007/978-1-4419-9982-5}{doi:10.1007/978-1-4419-9982-5}.
Ordinary exterior calculus, Stokes, and de Rham theory.

\bibitem{Mihara}
T. Mihara, ``Non-Archimedean Analytic Singular Homology Based on Cosimplicial
Perfectoid Spaces and Integration along Cycles,'' \emph{The Ramanujan Journal}
\textbf{68} (2025), article 98.
\href{https://doi.org/10.1007/s11139-025-01255-8}{doi:10.1007/s11139-025-01255-8};
\href{https://arxiv.org/abs/1211.1422}{arXiv:1211.1422}. Theorem 5.15 gives the cited
Stokes identity; the residue-characteristic hypothesis is essential to the comparison.

\bibitem{MS}
J. Ma\v r\'{\i}kov\'a and M. Shiota, ``Measuring definable sets in o-minimal fields,''
\emph{Israel Journal of Mathematics} \textbf{209} (2015), 687--714.
\href{https://doi.org/10.1007/s11856-015-1234-0}{doi:10.1007/s11856-015-1234-0};
\href{https://arxiv.org/abs/1402.4787}{arXiv:1402.4787}. The measure target is an
ordered semiring, not the original scalar field.

\bibitem{MuellerStrohmaier}
J. M\"uller and A. Strohmaier, ``The theory of Hahn-meromorphic functions, a
holomorphic Fredholm theorem, and its applications,'' \emph{Analysis \& PDE}
\textbf{7} (2014), no.~3, 745--770.
\href{https://doi.org/10.2140/apde.2014.7.745}{doi:10.2140/apde.2014.7.745};
\href{https://arxiv.org/abs/1205.0236}{arXiv:1205.0236}. A genuine and different
use of Hahn series in complex analysis; its analytic summation category is
\emph{not} identified with $\A_n$ here. \emph{Not to be confused with
\cite{MS}}, which is Ma\v r\'{\i}kov\'a--Shiota.

\bibitem{Neumann}
B. H. Neumann, ``On ordered division rings,'' \emph{Transactions of the American
Mathematical Society} \textbf{66} (1949), 202--252. The ordered-series support lemma is
the algebraic summability input.

\bibitem{OS}
T. Ohmoto and M. Shiota, ``$C^1$-triangulations of semialgebraic sets,''
\emph{Journal of Topology} \textbf{10} (2017), 765--775.
\href{https://doi.org/10.1112/topo.12024}{doi:10.1112/topo.12024};
\href{https://arxiv.org/abs/1505.03970}{arXiv:1505.03970}. See \S4 for triangulated
integration and Stokes, and the introduction for the non-Archimedean codomain
qualification.

\bibitem{ParRond}
A. Parusi\'nski and G. Rond, ``The Abhyankar--Jung theorem,'' \emph{Journal of
Algebra} \textbf{365} (2012), 29--41.
\href{https://doi.org/10.1016/j.jalgebra.2012.05.003}{doi:10.1016/j.jalgebra.2012.05.003};
\href{https://arxiv.org/abs/1103.2559}{arXiv:1103.2559}. Proposition~2.1 and its
finite-covering, power-substitution and removable-singularity steps are the
classical mechanism behind \cref{contours:thm:ramifiedleray}. \emph{Not to be
confused with \cite{PR}}, which is Perrucci--Roy.

\bibitem{Poonen}
B. Poonen, ``Maximally complete fields,'' \emph{L'Enseignement Math\'ematique} (2)
\textbf{39} (1993), 87--106.
\href{https://math.mit.edu/~poonen/papers/amsval.pdf}{Author's full text}. See the
Hahn-field construction and Corollary~4.

\bibitem{PR}
D. Perrucci and M.-F. Roy, ``Algebraic Winding Numbers,'' preprint,
\href{https://arxiv.org/abs/2305.08638}{arXiv:2305.08638}, version~3, July 18, 2024.
Theorems 12--13 and the corrected product-formula discussion.

\bibitem{Spivak}
M. Spivak, \emph{Calculus on Manifolds: A Modern Approach to Classical Theorems of
Advanced Calculus}, W. A. Benjamin, 1965. Ordinary differential forms, integration,
and Stokes at the coefficient level.

\bibitem{TajimaNabeshima}
S. Tajima and K. Nabeshima, ``An effective method for computing Grothendieck
point residue mappings,'' preprint, 2020.
\href{https://arxiv.org/abs/2011.09092}{arXiv:2011.09092}. Used for the local
residue, its transformation law, and the finite local-algebra discussion behind
\cref{contours:lem:finitejet}.

\bibitem{Temkin}
M. Temkin, ``Introduction to Berkovich analytic spaces,''
\href{https://arxiv.org/abs/1010.2235}{arXiv:1010.2235}, version~2, December 16, 2011.
Especially \S2.3.3 on the analytic affine line and its tree structure.

\bibitem{Tsikh}
A. K. Tsikh, \emph{Multidimensional Residues and Their Applications}, Translations of
Mathematical Monographs, vol.~103, American Mathematical Society, 1992.
\href{https://doi.org/10.1090/mmono/103}{doi:10.1090/mmono/103}.

\bibitem{vdDE}
L. van den Dries and P. Ehrlich, ``Fields of surreal numbers and exponentiation,''
\emph{Fundamenta Mathematicae} \textbf{167} (2001), 173--188; erratum, \textbf{168}
(2001), 295--297.
\href{https://doi.org/10.4064/fm167-2-3}{doi:10.4064/fm167-2-3}.

\bibitem{Woerheide}
A. Woerheide, \emph{O-minimal Homology}, Ph.D. thesis, University of Illinois at
Urbana-Champaign, 1996. The o-minimal singular homology route to definable
Jordan--Brouwer separation.

\end{thebibliography}
\end{document}
